\documentclass[11pt]{article}
\usepackage[a4paper,margin=29mm]{geometry}
\usepackage[T1]{fontenc}
\usepackage{lmodern}
\usepackage{microtype}
\usepackage{amsmath,amssymb,amsthm,mathtools}
\usepackage{mathrsfs}
\usepackage{booktabs}
\usepackage{longtable}
\usepackage{array}
\usepackage{enumitem}
\usepackage[dvipsnames]{xcolor}
\usepackage[colorlinks=true,linkcolor=MidnightBlue,citecolor=BrickRed,urlcolor=RoyalBlue]{hyperref}
\usepackage[nameinlink,noabbrev]{cleveref}
\emergencystretch=3em
\hypersetup{
  pdftitle={Exponential-Integral Periods across Arithmetic, Pade, Hankel, Reciprocal Geometry, and Frobenius Framing},
  pdfauthor={KEIO},
  pdfsubject={Algebraic independence, arithmetic Gevrey series, Hankel determinants, and reciprocal entire functions}
}

\newtheorem{theorem}{Theorem}[section]
\newtheorem{proposition}[theorem]{Proposition}
\newtheorem{corollary}[theorem]{Corollary}
\newtheorem{lemma}[theorem]{Lemma}
\newtheorem{conjecture}[theorem]{Conjecture}
\theoremstyle{definition}
\newtheorem{definition}[theorem]{Definition}
\theoremstyle{remark}
\newtheorem{remark}[theorem]{Remark}

\newcommand{\Q}{\mathbb Q}
\newcommand{\Qb}{\overline{\mathbb Q}}
\newcommand{\Qbar}{\overline{\mathbb Q}}
\newcommand{\C}{\mathbb C}
\newcommand{\Z}{\mathbb Z}
\newcommand{\N}{\mathbb N}
\newcommand{\trdeg}{\operatorname{trdeg}}
\newcommand{\Ein}{\operatorname{Ein}}
\newcommand{\EinDepth}[1]{\operatorname{Ein}_{#1}}
\newcommand{\Cin}{\operatorname{Cin}}
\newcommand{\Ei}{\operatorname{Ei}}
\newcommand{\Ci}{\operatorname{Ci}}
\newcommand{\Chi}{\operatorname{Chi}}
\newcommand{\Eul}{\mathcal E}
\newcommand{\Eexp}{\mathcal E_{\mathrm{exp}}}
\newcommand{\ord}{\operatorname{ord}}
\newcommand{\NP}{\operatorname{NP}}
\newcommand{\OO}{\mathcal O}
\newcommand{\capc}{\operatorname{cap}}
\newcommand{\Span}{\operatorname{span}}
\newcommand{\Si}{\operatorname{Si}}
\newcommand{\Shi}{\operatorname{Shi}}
\newcommand{\Cinh}{\operatorname{Cinh}}
\newcommand{\Log}{\operatorname{Log}}

\title{\bfseries Exponential-Integral Periods across Arithmetic,\\
Pad\'e, Hankel, Reciprocal Geometry, and Frobenius Framing}
\author{KEIO}
\date{Version 6.3 --- 31 August 2026}

\begin{document}
\maketitle

\begin{abstract}
Write
\(\Ein(z)=\int_0^z(1-e^{-t})\,dt/t\) and
\(\delta=eE_1(1)\).  We first repair and globalize the multipoint relation
theorem
\[
 \trdeg_{\Qb}\Qb\bigl(e^\lambda,\Ein(\lambda):\lambda\in\Lambda\bigr)
 =\dim_\Q\langle\Lambda\rangle_\Q+|\Lambda|.
\]
It follows that all nonzero algebraic-point \(\Ein\)-values are jointly
algebraically independent over the field generated by every algebraic
exponential.  Linear projection then determines the exact graphic-matroid
relation ideal of definite \(\Ein\)-periods and yields algebraic
independence of the countable family
\(\{\Si(n),\Cin(n),\Shi(n),\Cinh(n):n\ge1\}\), together with the complete
first-jet ideal of these four functions.  A generic-point transfer principle
also gives algebraically independent Hankel determinants and formal Favard
parameters, exact Pl\"ucker relations, simple algebraically independent
symmetric spectra, and realizations of every \(\Qbar\)-representable matroid
by explicit \(\Ein\)-periods.  Finite differences of \(\Ein(b+am)\) yield
explicit positive Hausdorff measures whose infinite Hankel matrices are
strictly totally positive.  Their moments, principal Hankel determinants,
Favard coordinates, and every finite Jacobi spectrum are algebraically
independent; the associated self-adjoint Jacobi operators have simple purely
absolutely continuous spectrum, exact Nevai limits, capacity asymptotics,
and arcsine zero distribution.

This finite theorem extends through the classical one-level
polyexponential tower
\[
\begin{gathered}
 \EinDepth{r}(z)=\sum_{n\geq1}\frac{(-1)^{n-1}z^n}{n^r n!},\\
 \trdeg_{\Qbar}\Qbar\bigl(e^{\lambda_i},
 \EinDepth{r}(\lambda_i):
 1\leq i\leq s,\ 1\leq r\leq R\bigr)=d+Rs.
\end{gathered}
\]
Thus all nonzero algebraic-point values in every order are countably
algebraically independent over the pure exponential field.  A positive
Hausdorff--Jacobi realization exists at every order, with its arithmetic
coordinates jointly algebraically independent across orders; for a fixed
order the spectra of two consecutive truncations form a sharp
transcendence basis of size \(2N-1\) for the full finite spectral tower.
Splitting Euler's Gamma integral at one further gives an all-order
Gamma-jet/tail decomposition whose resulting combinations are countably
algebraically independent.

On the divergent side we prove, for every prime and every nonzero integral
parameter, an exact two-branch \(p\)-adic radius formula for the classical
Euler--Pad\'e remainder, an equivalent local denominator law, a nonzero
\(p\)-curvature obstruction, and a fixed logarithmic Borel singularity.
Geronimus--Markov reduction evaluates an entire Laguerre--Hankel tower as a
Kummer \(U\)-function for all sizes, while a confluent reciprocal-Gamma
determinant becomes a monomer--dimer polynomial whose complex zeros lie in
an explicit vertical strip.  Bell--Gould moment functionals admit monic
common annihilators modulo every prime power and, by coefficientwise CRT,
modulo every integer; the \(N!\)-annihilator has degree
\(3N+O(\log N)\).  A scalar one-pole Borel--Laguerre lift has exact least
denominator \((M-1)!\).  A corrected compact double integral has size
\(\sim \pi\,16^{-n}/(12n\log2)\), but its integer form has a strictly
positive logarithmic coefficient at every prime in \([n+1,2n+1]\).
Positive mixed Hankel determinants
\(\mathcal P_d(\gamma,\delta)\) have full degree, and their first two values
recover \(\Q(\gamma,\delta)\) through an extension of degree at most four;
in particular \(\mathcal P_1+i\mathcal P_2\) is transcendental.

A common-shift theorem shows that translating an independent family by
\(\gamma\) loses at most one transcendence degree: dependence in an
explicit Euler residual family forces \(\gamma,\delta\) to be algebraically
independent, while all anchored residual differences are unconditionally
algebraically independent.  We also prove pairwise algebraicity exclusion
among \(E_1,\Ei,\Ci,\Chi\) at every positive algebraic level and construct
a positive quadratically sampled \(E_1\)-family with exact
\(e^{-m^2}/m^2\) decay and countable algebraic independence.
Finite algebraic-point \(\Ein\)-representations are rigid: the one-variable
Euler level has at most one algebraic zero, and the rank-one witness family
\(n\Ein(z)-\Ein(z^n)-(n-1)\gamma\) has an exact real-zero classification
and at most one algebraic zero.  The first two nontrivial inverse traces of
the divisor of \(\Ein-\gamma\) recover \(\gamma\) by a rational formula.

Finally the reciprocal sum descends through the Joukowski map to an entire
function \(H\) with an exact product for \(H'\).  Its values at distinct
algebraic points are algebraically independent over their exponential
field; hence every preimage of an algebraic level is transcendental.  At
the Euler level, algebraicity of the unique real reciprocal orbit would
force \(\gamma\) and \(\delta\) to be algebraically independent and would
produce an explicit transcendental element of the intersection of the
\(E\)-value and arithmetic-Gevrey value rings.  A positive Bessel--Laplace
representation places every zero of \(H\) in a strict left half-plane and
every reciprocal zero in the open left half-plane.  A level-zero
de Branges--Clark construction produces algebraically independent scattering
values at algebraic sample points and recovers \(\gamma\) from their
asymptotic phase.  We further obtain an exact Hadamard trace at the Euler
level, recover
\(\gamma\) nontrivially from the asymptotic phase of the genuinely
\(\gamma\)-free reciprocal zero set, and prove that the full primitive
clearing polynomial for the symmetric Laurent truncations saturates its
defining-height budget: the factorial approximation gain is exactly
cancelled by denominator clearing, while the level first appears
in the full top norm.  On the dyadic subsequence the clearing polynomial is
irreducible by a \(2\)-Eisenstein argument, its reciprocal lift is irreducible
of degree \(2N\), and the same critical exponent persists for the actual
minimal Weil height.  Rational polynomial capacity pullbacks and unimodular
full-determinant normalizations are shown to remain exactly critical.
Finally, the rank-two scalar semilinear Frobenius model with \(|q|\ne1\) is
proved to split locally, so bare Frobenius data cannot by itself detect the
global extension class.  The
remaining motivic gap is isolated instead as a depth-one numerical-to-motivic
comparison/fullness principle, sufficient but presently unproved for the
transcendence of $\gamma$.  These are structural closures and conditional
detectors.  Nothing here proves the irrationality or transcendence of
\(\gamma\) or \(\delta\), and no named open problem is claimed solved.
\end{abstract}

\tableofcontents

\section{Introduction and scope}

Euler's constant and the Euler--Gompertz constant are
\[
 \gamma=\lim_{n\to\infty}\left(\sum_{k=1}^n\frac1k-\log n\right),
 \qquad
 \delta=\int_0^\infty\frac{e^{-t}}{1+t}\,dt=eE_1(1).
\]
Neither number is known to be irrational; see Lagarias' survey for context
\cite{Lagarias2013}.  The elementary connection formula
\begin{equation}\label{eq:hardy-connection}
 \Ein(1)=\gamma+E_1(1)=\gamma+\frac{\delta}{e}
\end{equation}
places this problem on the interface between Siegel \(E\)-functions and
arithmetic Gevrey series of order one.  Throughout, the pure exponential
field is
\[
 \Eexp:=\Qbar(e^\alpha:\alpha\in\Qbar).
\]
The purpose of this note is to make
several rigorous statements on that interface and, equally importantly, to
mark the point at which they stop.

The first result is a multipoint algebraic-independence theorem.  Its proof is
an application of Kolchin--Ostrowski followed by Beukers' refined
Siegel--Shidlovskii theorem, but a denominator step in the differential field
is essential: arbitrary rational functions cannot simply be expanded as
finite Laurent sums.  We supply the missing Darboux-polynomial argument.

\begin{theorem}[Multipoint \(\Ein\) theorem]\label{thm:master-intro}
For distinct \(\lambda_1,\ldots,\lambda_s\in\Qb^\times\), put
\(d=\dim_\Q\langle\lambda_1,\ldots,\lambda_s\rangle_\Q\).  Then
\[
 \trdeg_{\Qb}\Qb\bigl(e^{\lambda_i},\Ein(\lambda_i):1\le i\le s\bigr)=d+s.
\]
Moreover the \(\Ein(\lambda_i)\) are algebraically independent over
\(\Qb(e^{\lambda_1},\ldots,e^{\lambda_s})\), and every algebraic relation
among the displayed numbers is generated by the multiplicative, or toric,
relations among the exponential coordinates.
\end{theorem}

The most concrete consequence for the classical constants is the following.

\begin{theorem}[Six constants]\label{thm:six-intro}
Among
\[
 \gamma,\quad \delta,\quad E_1(1),\quad \Ei(1),\quad \Ci(1),\quad \Chi(1)
\]
at most one number is algebraic.  Consequently at least five are
transcendental.
\end{theorem}

This is a density statement, not an identification statement: it does not say
which five numbers are transcendental.  In particular it does not settle
\(\gamma\) or \(\delta\) individually.

The same multipoint theorem produces a genuinely positive spectral model,
not merely the formal Favard system obtained from the raw moments
\(\Ein(m+1)\).

\begin{theorem}[Positive Hausdorff--Jacobi realization]
\label{thm:positive-jacobi-intro}
For algebraic \(a,b>0\) and \(r\geq1\), put
\[
 \mu_m=(-1)^{r-1}\mathsf D^r\Ein(b+am).
\]
Then
\[
 \mu_m=\int_{e^{-a}}^1x^m
 \frac{x^{b/a-1}(1-x)^r}{-\log x}\,dx.
\]
The infinite Hankel matrix \((\mu_{i+j})_{i,j\geq0}\) is strictly totally
positive.  The moments, the principal Hankel determinants, and the complete
Favard coordinate family are each algebraically independent over
\(\Eexp\).  The associated Jacobi operator has simple purely absolutely
continuous spectrum \([e^{-a},1]\); every finite principal spectrum consists
of jointly algebraically independent real numbers, and the eigenvalue
counting measures converge to the arcsine law on that interval.
\end{theorem}

\begin{theorem}[Resonant polyexponential closure]
\label{thm:polyexponential-intro}
For \(q\geq1\), put
\[
 \EinDepth{q}(z)=
 \sum_{n\geq1}\frac{(-1)^{n-1}z^n}{n^q n!}.
\]
If \(\lambda_1,\ldots,\lambda_s\in\Qbar^\times\) are distinct and
\(d=\dim_\Q\langle\lambda_1,\ldots,\lambda_s\rangle_\Q\), then, for every
\(R\geq1\),
\[
 \trdeg_{\Qbar}\Qbar\bigl(
 e^{\lambda_i},\EinDepth{q}(\lambda_i):
 1\leq i\leq s,\ 1\leq q\leq R\bigr)=d+Rs.
\]
All the \(\EinDepth{q}(\lambda)\), jointly in \(q\geq1\) and
\(\lambda\in\Qbar^\times\), are algebraically independent over \(\Eexp\).
Their positive finite-difference moment systems have the same
Hausdorff--Jacobi conclusions as Theorem~\ref{thm:positive-jacobi-intro}.
For a fixed order, if \(\Sigma_j\) is the spectrum of the \(j\)-by-\(j\)
Jacobi truncation, then
\[
 \trdeg_{\Eexp}\Eexp(\Sigma_1,\ldots,\Sigma_N)=2N-1;
\]
the two consecutive spectra \(\Sigma_{N-1}\cup\Sigma_N\) are jointly
algebraically independent and form a sharp transcendence basis.
\end{theorem}

These are classical one-level polyexponentials, not newly defined special
functions.  The contribution here is the explicit resonant integral-shift
proof, its multipoint toric relation ideal, and the arithmetic spectral
applications; exact historical priority is treated conservatively below.
Splitting the Euler Gamma integral at one also yields
\[
 g_k-T_k=(-1)^{k+1}\EinDepth{k+1}(1),
\]
where \(\Gamma(s)=s^{-1}+\sum_{k\geq0}g_ks^k\) and
\(T_k=k!^{-1}\int_1^\infty e^{-u}(\log u)^k\,du/u\).  Hence the entire
family \(\{g_k-T_k:k\geq0\}\) is algebraically independent over
\(\Eexp\).

A second linear-algebra consequence is a precise defect-one principle.  A
common translation of \(r\) independent coordinates can lose at most one
transcendence degree.  Applied to explicit residual combinations of
\(E_1(2^N)\) and \(\Ei(n)\), it gives an unconditional countable family of
algebraically independent anchored differences; any relation among the
unanchored residuals would force \(\gamma\) and \(\delta\) to be
algebraically independent over \(\Eexp\).

Our second theme is zero removal.  An \(E\)-function is stable under division
by \(z-\alpha\) when it vanishes at a nonzero algebraic point \(\alpha\).
The analogous assertion for arithmetic Gevrey series is conjectural in the
form relevant to Euler's factorial series.  We prove an exact obstruction.

\begin{theorem}[Denominator obstruction]\label{thm:den-intro}
Let \(0\ne\xi\in\Z\), \(\alpha\in\Qb\), and
\[
 \Eul_\xi(z)=\sum_{n\ge0}(-1)^n n!\xi^n z^n.
\]
Write
\[
 \frac{\Eul_\xi(z)-\alpha}{z-1}=\sum_{n\ge0}n!b_nz^n.
\]
If \(K\) is a number field containing \(\alpha\), and \(\Delta_N\) is the
least positive integer such that
\(\Delta_Nb_0,\ldots,\Delta_Nb_N\in\OO_K\), then
\[
 \lim_{N\to\infty}\frac{\log\Delta_N}{N}=+\infty.
\]
Thus the quotient is not an arithmetic Gevrey series of order one and, in
particular, does not belong to the arithmetic-Gevrey \(\mathbf D\)-class.
\end{theorem}

For \(\xi=1\), the same Pad\'e construction has a sharp local signature at
every prime.

\begin{theorem}[Pad\'e radius transition]\label{thm:radius-intro}
Let
\[
 q_n=\sum_{k=0}^n k!\binom nk^2,
\]
let \(s_n=P_n(1)\) be the numerator in the classical Pad\'e approximation
\(Q_n(z)\Eul_1(z)-P_n(z)=R_n(z)\), and put
\[
 F_\alpha(z)=\sum_{n\ge0}(q_n\alpha-s_n)\frac{z^n}{n!}.
\]
For every prime \(p\) and \(\alpha\in\Q_p\), with
\(\delta_p=\sum_{m\ge0}(-1)^m m!\in\Z_p\),
\[
 R_p(F_\alpha)=
 \begin{cases}
 p^{1/(p-1)},&\alpha=\delta_p,\\
 p^{-1/(p-1)},&\alpha\ne\delta_p.
 \end{cases}
\]
Moreover \(F_\alpha\in\Z_p[[z]]\) if and only if \(\alpha=\delta_p\).
\end{theorem}

The local arithmetic can also be organized directly on the Bell and Gould
moment functionals.  If \(M=\prod p^{r_p}\), there is a monic polynomial
\(C_M\in\Z[X]\), of degree
\[
 \max_{p\mid M}\bigl(m_{p,r_p}+pr_p\bigr),
 \qquad m_{p,r}=\min\{m:v_p(m!)\geq r\},
\]
which annihilates both functionals modulo \(M\) after multiplication by
every polynomial in \(\Z[X]\).  For \(M=N!\) this degree is
\(3N+O(\log N)\).  This is an exact common-divisibility theorem; it does not
assert that the resulting \(\delta\)-linear forms are small.  A paired
logarithmic integral does, however, give an unconditional nonzero
\(\Eexp\)-transcendental mixed value with coefficient tending to infinity.

The third theme imports the capacity viewpoint.  The relevant determinant is
not merely small; it sits exactly on the critical constant.

\begin{theorem}[Shifted harmonic Hankel determinant]\label{thm:hankel-intro}
Let \(H_0=0\), \(H_n=\sum_{k=1}^n1/k\), and fix \(x\in\C\).  Then
\begin{equation}\label{eq:hankel-intro}
 \det(H_{i+j}-x)_{0\le i,j\le n}
 =(-1)^n\frac{2H_n-x}
 {n!\displaystyle\prod_{j=1}^n\binom{2j}{j}\binom{2j-1}{j}}.
\end{equation}
Consequently
\[
 \lim_{n\to\infty}
 \left|\det(H_{i+j}-x)_{0\le i,j\le n}\right|^{1/n^2}=\frac14
 =\capc([0,1]).
\]
The shift by \(x\), including \(x=\gamma\), changes only the
subquadratic part of the determinant asymptotic.
\end{theorem}

The mixed positive-moment construction gives a complementary positive
result.  It produces nonzero polynomials
\(\mathcal P_d(\gamma,\delta)>0\) of exact total degree \(2d\), with
\[
 [\gamma^d\delta^d]\mathcal P_d
 =(-1)^{d(d-1)/2}((d-1)!)^d.
\]
More unexpectedly, \(\Q(\gamma,\delta)\) is algebraic of degree at most
four over \(\Q(\mathcal P_1,\mathcal P_2)\).  Thus the first two positive
determinants preserve the full transcendence degree of the original pair;
at least one is transcendental over \(\Eexp\), and
\(\mathcal P_1+i\mathcal P_2\) is unconditionally transcendental.

The exact mixed family suggested by \eqref{eq:hardy-connection} is analysed in
\cref{sec:mixed}.  The conclusion is deliberately negative but sharp: the
proposed one-family zero-removal implication is equivalent to the
transcendence of \(\gamma\), rather than a smaller lemma from which that
transcendence follows cheaply.

The reciprocal quotient supplies the final theme.  The entire function \(H\)
is characterized on \(\C^\times\) by
\[
 H(z+z^{-1})=\Ein(z)+\Ein(z^{-1})
\]
has an explicit critical lattice and no algebraic value at an algebraic
point.  A positive Bessel--Laplace representation proves that the unique real
point in every real fibre is its unique rightmost point.  At level zero this
puts all zeros of \(H\), and all reciprocal zeros of \(H(z+z^{-1})\), in the
left half-plane and proves their simplicity and transcendence.  Rotation and
centering then give a Cartwright Hermite--Biehler function and a
zero-mean-type meromorphic inner quotient.  Its values at nonzero algebraic
real points are algebraically independent over the pure exponential field,
while their asymptotic phase recovers \(\gamma\) with an explicit first
correction.

The divisor of \(H-2\gamma\) has an exact second Hadamard trace, whereas the
upper exterior zeros of the genuinely \(\gamma\)-free level \(H=0\) recover
\(\gamma\) through a second-order phase shift.  The algebraic zeros of the
Laurent truncations approximate a fixed reciprocal zero with exponent one
relative to the full primitive clearing polynomial, but integerization makes
those auxiliary polynomial values grow.  On \(N=2^r\), a \(2\)-Eisenstein
argument and a norm-sign obstruction prove that the reciprocal clearing
polynomial is itself irreducible of degree \(2N\); hence exponent one persists
for the actual minimal Weil height.  Vieta's formulas separately show that
the level is invisible to every proper exterior power and first appears in
the factorially weighted top norm.

\paragraph{Status and priority.}
The one-point exponential-integral phenomenon is already present in
Mahler's work \cite{Mahler1968}.  The finite multipoint toric formulation
here packages a direct use of classical Kolchin--Ostrowski and
Siegel--Shidlovskii machinery \cite{Kolchin,Beukers2006}, with an explicit
Laurent-denominator repair.  The higher-order functions used here are
Hardy's classical one-level polyexponentials
\cite{Hardy1905PartII,Boyadzhiev2007}; Shidlovskii's classical theorem
covers the nearby nonintegral rational-shift family
\cite[Chap.~7, Sec.~3]{Shidlovskii1989}, and recent analytic context is
given by Aminov--Arnaudo \cite{AminovArnaudo2026}.  We record the regularized
resonant integral-shift subsystem considered here, its multipoint toric closure,
and its Jacobi consequences, without claiming priority for the underlying
all-order phenomenon.  The \(p\)-adic Pad\'e formulas are due to
Matala-aho and Zudilin \cite{MZ2018}.  The case \(x=0\) of
\eqref{eq:hankel-intro} appears in Cigler \cite{Cigler2017}; the rank-one shift and the capacity
interpretation are recorded here.  The logarithmic \(E\)-value corollary in
\cref{sec:log} is an application of a recent theorem of Fischler and Rivoal
\cite{FRlog}.  Related recent independence results for generalized constants
appear in Powers \cite{Powers2025}.
No claim of historical priority is made for the particular syntheses in this
draft.  Most importantly, nothing in this paper proves that \(\gamma\) or
\(\delta\) is irrational or transcendental.

\section{Preliminaries}

We use
\[
 \Ein(z)=\sum_{n\ge1}\frac{(-1)^{n-1}z^n}{n\,n!},
 \qquad
 \Ein'(z)=\frac{1-e^{-z}}{z}.
\]
This is a strict \(E\)-function.  We need two standard transfer principles.
In contrast, \(E_1(z)=\Ein(z)-\gamma-\Log z\) has a logarithmic
singularity at the origin and is not an \(E\)-function.  None of the
effective \(E\)-value arguments below is applied directly to \(E_1\).
First, the primitive form of the Kolchin--Ostrowski theorem \cite{Kolchin}
says that if a
differential field \(K\) of characteristic zero has algebraically closed
constant field \(C\), the ambient differential extension has the same
constant field, and \(y_i'\in K\), then the primitives \(y_i\) are
algebraically dependent over \(K\) only if
\(\sum c_i y_i'=Dg\) for some \(c_i\in C\), not all zero, and \(g\in K\).
Second, Beukers' refinement of the Siegel--Shidlovskii theorem
\cite{Beukers2006} lifts every
algebraic relation among values of a system of \(E\)-functions at a nonzero
algebraic ordinary point to a functional algebraic relation.

For the divergent side, call a formal series
\[
 G(z)=\sum_{n\ge0}n!a_nz^n,\qquad a_n\in\Qb,
\]
an \emph{arithmetic Gevrey series of order one} if it is holonomic and there
is \(C>0\) such that, for every \(N\), the conjugates of
\(a_0,\ldots,a_N\) have size at most \(C^{N+1}\) and a common integer
denominator at most \(C^{N+1}\).  This normalization is sufficient for every
nonmembership statement below.

\section{The Laurent denominator lemma}

The key differential-field point is isolated here.  Let
\(\nu_1,\ldots,\nu_d\in\Qb\) be \(\Q\)-linearly independent, put
\[
 R=\Qb(z)[X_1^{\pm1},\ldots,X_d^{\pm1}],\qquad K=\operatorname{Frac}(R),
\]
and extend \(D=d/dz\) by \(DX_j=-\nu_jX_j\).  Thus the specialization
\(X_j=e^{-\nu_jz}\) respects the derivation.

\begin{lemma}[No nonmonomial Darboux denominator]\label{lem:darboux}
If an irreducible \(P\in R\) satisfies \(P\mid DP\), then \(P\) is a unit
of the Laurent ring.
\end{lemma}

\begin{proof}
Write \(DP=AP\) with \(A\in R\).  If \(A=0\), Laurent-coefficient
comparison shows that every nonzero
coefficient \(a_u(z)\) would satisfy
\(a_u'/a_u=u\cdot\nu\); rationality forces \(u=0\), and then \(a_0\) is
constant.  Thus \(P\) is already a unit.  We may therefore assume
\(A\ne0\).  The Newton polytope of \(DP\) is
contained in that of \(P\), while
\(\NP(AP)=\NP(A)+\NP(P)\).  Hence \(\NP(A)=\{0\}\), so
\(A\in\Qb(z)\).  Write
\(P=\sum_{u\in S}a_u(z)X^u\).  For \(u,v\in S\), comparison of coefficients
gives
\[
 \frac{d}{dz}\log\frac{a_u}{a_v}=(u-v)\mathbin{\cdot}\nu.
\]
A nonzero rational function cannot have logarithmic derivative equal to a
nonzero constant: its logarithmic derivative tends to zero at infinity (and
has only simple polar principal parts).  Since the \(\nu_j\) are
\(\Q\)-linearly independent, \((u-v)\cdot\nu=0\) forces \(u=v\).
Thus \(S\) has one element and \(P\) is a Laurent monomial times an element
of \(\Qb(z)^\times\), hence a unit.
\end{proof}

\begin{lemma}[Denominator removal]\label{lem:den-removal}
If \(g\in K\) and \(Dg\in R\), then \(g\in R\).
\end{lemma}

\begin{proof}
Write \(g=A/B\) in lowest terms.  If an irreducible \(P\) occurs in \(B\),
let \(m=-\ord_P(g)>0\).  Unless \(P\mid DP\), differentiation lowers the
\(P\)-valuation from \(-m\) to \(-m-1\), contradicting \(Dg\in R\).
Thus \(P\mid DP\), impossible by \cref{lem:darboux} because a denominator
factor is not a unit.  Hence \(B\) is a unit.
\end{proof}

\begin{lemma}[Exact-derivative obstruction]\label{lem:exact}
Suppose \(c_u\in\Qb\) are finitely supported and
\[
 Dg=\frac1z\sum_{u\in\Z^d}c_uX^u,
 \qquad g\in K.
\]
Then every \(c_u\) is zero.
\end{lemma}

\begin{proof}
By \cref{lem:den-removal}, write \(g=\sum a_u(z)X^u\) with
\(a_u\in\Qb(z)\).  Coefficient comparison gives
\begin{equation}\label{eq:scalar-rat}
 a_u'(z)-(u\cdot\nu)a_u(z)=\frac{c_u}{z}.
\end{equation}
If \(u=0\), a derivative of a rational function has zero residue at every
finite pole, so \eqref{eq:scalar-rat} forces \(c_0=0\).  If \(u\ne0\), then
\(u\cdot\nu\ne0\).  A finite pole of \(a_u\) away from zero would produce an
uncancelled pole of one order higher in \(a_u'\).  A pole at zero produces the
same contradiction, whereas a function regular at zero makes the left-hand
side regular there.  Thus a nonzero right-hand side \(c_u/z\) is impossible.
\end{proof}

The differential constants of \(K\) are \(\Qb\).  One way to see this is to
apply the same denominator argument and then compare Laurent coefficients;
the equation \(a_u'-(u\cdot\nu)a_u=0\) has no nonzero rational solution for
\(u\ne0\), and \(a_0'=0\).

\section{Multipoint algebraic independence}

\begin{theorem}\label{thm:master}
Let \(\Lambda=\{\lambda_1,\ldots,\lambda_s\}\subset\Qb^\times\) consist of
distinct elements, and put
\(d=\dim_\Q\langle\Lambda\rangle_\Q\).  Then
\[
 \trdeg_{\Qb}\Qb\bigl(e^{\lambda_i},\Ein(\lambda_i):1\le i\le s\bigr)=d+s.
\]
If
\[
 L_\Lambda=\{m\in\Z^s:\textstyle\sum_i m_i\lambda_i=0\},
\]
the kernel of
\[
 \Qb[U_1^{\pm1},\ldots,U_s^{\pm1},V_1,\ldots,V_s]
 \longrightarrow\C,
 \quad U_i\mapsto e^{\lambda_i},\quad V_i\mapsto\Ein(\lambda_i),
\]
is generated by \(U^{m_+}-U^{m_-}\), \(m\in L_\Lambda\).
\end{theorem}

\begin{proof}
The additive group generated by \(\Lambda\) is free of rank \(d\).  Choose a
\(\Z\)-basis \(\nu_1,\ldots,\nu_d\), so that
\(\lambda_i=m_i\cdot\nu\) with \(m_i\in\Z^d\).  Equivalently, one may first
choose a \(\Q\)-basis and divide it by a common denominator; this lattice
clearing is necessary for the Laurent-ring argument.

Set \(X_j=e^{-\nu_jz}\) formally and let \(K\) be the differential field of
the previous section.  Distinct Laurent monomials \(X^u\) are linearly
independent over \(\Qb(z)\); after specialization this is the standard linear
independence of exponentials with distinct exponents.  Hence the exponential
torus has transcendence degree \(d\).

Put \(Y_i=\Ein(\lambda_i z)\).  Then
\[
 DY_i=\frac{1-X^{m_i}}{z}\in K.
\]
Under \(X_j=e^{-\nu_jz}\), the field generated by \(K\) and the \(Y_i\)
lies in \(\Qbar((z))\), whose differential constants are exactly
\(\Qbar\).  Thus the no-new-constants hypothesis in
Kolchin--Ostrowski is satisfied.
If the \(Y_i\) were algebraically dependent over \(K\), the primitive
Kolchin--Ostrowski theorem would give constants \(c_i\in\Qb\), not all zero,
and \(g\in K\) such that
\[
 Dg=\sum_{i=1}^s c_iDY_i
 =\frac{\sum_i c_i-\sum_i c_iX^{m_i}}{z}.
\]
Because the \(m_i\) are distinct, \cref{lem:exact} forces every \(c_i=0\), a
contradiction.  Thus the \(Y_i\) are algebraically independent over \(K\),
and the functional transcendence degree is \(d+s\).

The same degree holds over \(\C(z)\): after clearing denominators, a
relation there would give a finite homogeneous linear system with
\(\Qbar\)-valued Taylor coefficients, so scalar-extension invariance of
rank would produce a relation over \(\Qbar(z)\).

For specialization, include in the differential system the constant function
\(1\), the exponentials \(e^{\pm\nu_jz}\) and
\(e^{-\lambda_i z}\), and all \(Y_i\).  The additional characters
\(e^{-\lambda_i z}\) ensure that the displayed first-order system is closed
even when some \(\lambda_i\) is a nontrivial integral combination of the
\(\nu_j\).  They are
\(E\)-functions, the system has coefficients in \(\Qb(z)\), and its only
finite pole is at \(z=0\); hence \(z=1\) is ordinary.  Beukers' refined
Siegel--Shidlovskii theorem transfers the functional transcendence degree to
the values at \(1\).

Finally, since the \(V_i\) remain algebraically independent over the
exponential field, the full relation ideal is the relation ideal of the
exponential torus.  Its character lattice is exactly \(L_\Lambda\), which
gives the stated binomial generators.
\end{proof}

\begin{remark}
The exclusion of \(0\) from \(\Lambda\) is substantive only for notation:
\(e^0=1\) and \(\Ein(0)=0\) contribute fixed algebraic coordinates and should
simply be deleted before applying the theorem.
\end{remark}

% Standalone insertion for gamma_boundary_revision.tex.
% Required in the preamble if these commands are not already present:
% \providecommand{\Si}{\operatorname{Si}}
% \providecommand{\Cin}{\operatorname{Cin}}
% \providecommand{\Shi}{\operatorname{Shi}}
% \providecommand{\Cinh}{\operatorname{Cinh}}
% Suggested bibliography keys:
%   AdamczewskiFaverjon2025 = B. Adamczewski and C. Faverjon,
%     Algebraic independence measures for values of E-functions and
%     M-functions, arXiv:2502.09999 (2025).
%   Delaygue2025 = E. Delaygue, A Lindemann--Weierstrass theorem for
%     E-functions, arXiv:2210.12046v2 (2025).

\providecommand{\Si}{\operatorname{Si}}
\providecommand{\Cin}{\operatorname{Cin}}
\providecommand{\Shi}{\operatorname{Shi}}
\providecommand{\Cinh}{\operatorname{Cinh}}
\providecommand{\Eexp}{\mathcal E_{\mathrm{exp}}}

\section{Pure exponential base change and special-function periods}
\label{sec:efunction-consequences}

The toric relation theorem has a useful global form.  Define the pure
exponential field
\begin{equation}\label{eq:pure-exponential-field}
 \Eexp:=\Qbar\bigl(e^\alpha:\alpha\in\Qbar\bigr)\subset\C.
\end{equation}
Although this is an infinitely generated field, every coefficient of a
polynomial relation belongs to a subfield generated by only finitely many
algebraic exponentials.  This elementary observation removes the toric
coordinates altogether.

\begin{theorem}[Pure exponential base change]\label{thm:pure-exp-base}
The family
\[
 \bigl\{\Ein(\lambda):\lambda\in\Qbar^\times\bigr\}
\]
is algebraically independent over \(\Eexp\), in the finite-subset sense.
Equivalently, for any distinct
\(\lambda_1,\ldots,\lambda_s\in\Qbar^\times\),
\[
 \operatorname{trdeg}_{\Eexp}
 \Eexp\bigl(\Ein(\lambda_1),\ldots,\Ein(\lambda_s)\bigr)=s.
\]
\end{theorem}

\begin{proof}
Suppose that a polynomial relation exists.  Its finitely many coefficients
belong to
\(\Qbar(e^\omega:\omega\in\Omega)\) for some finite
\(\Omega\subset\Qbar\).  Clear coefficient denominators and apply
Theorem~\ref{thm:master} to
\[
 \Sigma=\{\lambda_1,\ldots,\lambda_s\}
        \cup(\Omega\setminus\{0\}).
\]
That theorem makes all \(\Ein(\sigma)\), \(\sigma\in\Sigma\),
algebraically independent over
\(\Qbar(e^\sigma:\sigma\in\Sigma)\).  In particular the selected
\(\Ein(\lambda_j)\) are independent over the coefficient field, a
contradiction.
\end{proof}

We shall also use the following elementary linear-image consequence.  It is
recorded because it identifies full relation ideals, rather than only their
transcendence degrees.

\begin{lemma}[Linear images of the \(\Ein\)-coordinates]
\label{lem:linear-image-ein}
Let \(Y=(\Ein(\lambda_1),\ldots,\Ein(\lambda_s))^t\), where the
\(\lambda_j\) are distinct nonzero algebraic numbers, and let
\(B\in M_{r,s}(\Qbar)\).  Put \(W=BY\).  The kernel of
\[
 \Eexp[T_1,\ldots,T_r]\longrightarrow\C,
 \qquad T_j\longmapsto W_j,
\]
is the linear ideal
\[
 \left(\sum_{j=1}^r c_jT_j:
       c\in\Qbar^r,\ c^tB=0\right).
\]
In particular, the \(W_j\) are algebraically independent over \(\Eexp\)
when \(B\) has full row rank.
\end{lemma}

\begin{proof}
By Theorem~\ref{thm:pure-exp-base}, the entries of \(Y\) are independent
indeterminates over \(\Eexp\).  After invertible row and column operations
over \(\Qbar\), the map \(Y\mapsto BY\) is a coordinate projection of rank
\(\operatorname{rank}B\).  Its polynomial kernel is therefore exactly the
displayed linear ideal.
\end{proof}

\subsection{A graphic matroid of definite periods}

Let \(G=(V,E)\) be a finite oriented graph, let \(c(G)\) be its number of
connected components, and attach pairwise distinct labels
\(\lambda_v\in\Qbar^\times\) to its vertices.  For an edge \(e\), write
\(s(e)\) and \(t(e)\) for its source and target and set
\begin{equation}\label{eq:graph-period}
 P_e:=\int_{\lambda_{s(e)}}^{\lambda_{t(e)}}
             \frac{1-e^{-z}}{z}\,dz
     =\Ein(\lambda_{t(e)})-\Ein(\lambda_{s(e)}).
\end{equation}
The integrand is entire after filling its removable singularity at zero, so
the integral is path independent.

Let \(\partial:\Qbar^E\to\Qbar^V\) be the oriented incidence map,
\(\partial e=t(e)-s(e)\).  Its kernel is the cycle space of \(G\).

\begin{theorem}[Graphic-period relation ideal]\label{thm:graphic-period}
The kernel of
\[
 \Eexp[X_e:e\in E]\longrightarrow\C,
 \qquad X_e\longmapsto P_e,
\]
is
\begin{equation}\label{eq:cycle-ideal}
 \left(\sum_{e\in E}c_eX_e:c\in\ker\partial\right).
\end{equation}
Thus every algebraic relation among the periods \(P_e\) is generated by the
signed cycle sums.  In particular,
\[
 \operatorname{trdeg}_{\Eexp}\Eexp(P_e:e\in E)
   =|V|-c(G),
\]
and the periods attached to any spanning forest are algebraically
independent over \(\Eexp\).
\end{theorem}

\begin{proof}
Put \(Y_v=\Ein(\lambda_v)\).  In vector notation,
\((P_e)_{e\in E}=\partial^t(Y_v)_{v\in V}\).  Theorem
\ref{thm:pure-exp-base} makes the \(Y_v\) algebraically independent over
\(\Eexp\), so Lemma~\ref{lem:linear-image-ein} identifies the kernel with
the linear relations among the rows of \(\partial^t\), namely
\(\ker\partial\).  The incidence matrix has rank \(|V|-c(G)\), and its
cycle space is spanned by signed circuit vectors.  This proves every claim.
\end{proof}

\subsection{The fourth-root orbit}

Define the four entire integral functions
\begin{equation}\label{eq:four-entire-integrals}
\begin{aligned}
 \Si(z)&=\int_0^z\frac{\sin t}{t}\,dt,&
 \Cin(z)&=\int_0^z\frac{1-\cos t}{t}\,dt,\\
 \Shi(z)&=\int_0^z\frac{\sinh t}{t}\,dt,&
 \Cinh(z)&=\int_0^z\frac{\cosh t-1}{t}\,dt.
\end{aligned}
\end{equation}
Here \(\Cinh(z)=\Chi(z)-\gamma-\Log z\) on a compatible logarithmic
branch; the integral in \eqref{eq:four-entire-integrals} is its entire
normalization.  Termwise comparison gives
\begin{equation}\label{eq:fourth-root-transform}
\begin{aligned}
 \Ein(z)&=\Shi(z)-\Cinh(z),&
 \Ein(-z)&=-\Shi(z)-\Cinh(z),\\
 \Ein(iz)&=\Cin(z)+i\Si(z),&
 \Ein(-iz)&=\Cin(z)-i\Si(z).
\end{aligned}
\end{equation}

\begin{theorem}[Fourth-root orbit]\label{thm:fourth-root-orbit}
For every \(a\in\Qbar^\times\), the six numbers
\begin{equation}\label{eq:six-orbit-values}
 e^a,\quad e^{ia},\quad
 \Si(a),\quad\Cin(a),\quad\Shi(a),\quad\Cinh(a)
\end{equation}
are algebraically independent over \(\Qbar\).  More strongly, the last four
numbers are algebraically independent over \(\Eexp\).
\end{theorem}

\begin{proof}
The four arguments \(a,-a,ia,-ia\) are distinct.  Their \(\Q\)-span has
dimension two, with basis \(a,ia\).  Theorem~\ref{thm:master} therefore
gives transcendence degree six for the corresponding exponential and
\(\Ein\)-coordinates.  The four-by-four transformation
\eqref{eq:fourth-root-transform} is invertible over \(\Qbar\), proving the
first assertion.  The stronger assertion follows directly from
Theorem~\ref{thm:pure-exp-base} and the same invertible transformation.
\end{proof}

\begin{corollary}[A countable independent family]
\label{cor:integer-integral-family}
The countable family
\begin{equation}\label{eq:integer-integral-family}
 \bigl\{\Si(n),\Cin(n),\Shi(n),\Cinh(n):n\ge1\bigr\}
\end{equation}
is algebraically independent over \(\Eexp\).
\end{corollary}

\begin{proof}
It suffices to treat a finite set of positive integers.  The fourth-root
orbits \(\{n,-n,in,-in\}\) are pairwise disjoint as \(n\) varies.  Apply
Theorem~\ref{thm:pure-exp-base} to their union and perform the invertible
transformation \eqref{eq:fourth-root-transform} separately in every block.
\end{proof}

\begin{remark}[Relation to recent sine-integral results]
Delaygue proved that, when
\(a_1^2,b_1^2,\ldots,a_m^2,b_m^2\) are pairwise distinct algebraic
numbers, the integrals
\(\int_{a_j}^{b_j}(\sin t/t)\,dt\) are linearly independent over
\(\Qbar\); see \cite[Proposition~4.4]{Delaygue2025}.  Theorem
\ref{thm:pure-exp-base} and Lemma~\ref{lem:linear-image-ein} upgrade this,
under the same nonzero square-distinctness hypothesis, to algebraic
independence over \(\Eexp\), because
\[
 \int_a^b\frac{\sin t}{t}\,dt
 =\frac{\Ein(ib)-\Ein(-ib)-\Ein(ia)+\Ein(-ia)}{2i}.
\]
This is a strict strengthening of the stated conclusion of Delaygue's
proposition, but a separate priority search in the older
Mahler--Shidlovskii literature is required before presenting the explicit
upgrade as historically new.
\end{remark}

\subsection{The exact first-jet ideal}

The preceding theorem also determines every algebraic relation involving
the first derivatives of the four functions.  This is stronger than a
transcendence-degree count.

\begin{theorem}[Exact first-jet ideal]\label{thm:first-jet-ideal}
Fix \(a\in\Qbar^\times\), and abbreviate
\[
 S=\Si(a),\quad C=\Cin(a),\quad H=\Shi(a),\quad K=\Cinh(a)
\]
and
\[
 S_1=\Si'(a),\quad C_1=\Cin'(a),\quad
 H_1=\Shi'(a),\quad K_1=\Cinh'(a).
\]
The kernel of the evaluation homomorphism
\[
 \Qbar[\mathsf S,\mathsf C,\mathsf H,\mathsf K,
       \mathsf S_1,\mathsf C_1,\mathsf H_1,\mathsf K_1]
 \longrightarrow\C
\]
at \((S,C,H,K,S_1,C_1,H_1,K_1)\) is generated by the two polynomials
\begin{align}
 (a\mathsf S_1)^2+(1-a\mathsf C_1)^2-1,
 \label{eq:trig-jet-relation}\\
 (1+a\mathsf K_1)^2-(a\mathsf H_1)^2-1.
 \label{eq:hyperbolic-jet-relation}
\end{align}
In particular, no nonzero polynomial relation mixes the four zeroth-order
values with the derivative coordinates.
\end{theorem}

\begin{proof}
Put \(x=e^{ia}\) and \(y=e^a\).  Differentiating
\eqref{eq:four-entire-integrals} gives
\begin{equation}\label{eq:jet-torus-coordinates}
\begin{aligned}
 aS_1&=\frac{x-x^{-1}}{2i},&
 1-aC_1&=\frac{x+x^{-1}}2,\\
 aH_1&=\frac{y-y^{-1}}2,&
 1+aK_1&=\frac{y+y^{-1}}2.
\end{aligned}
\end{equation}
Conversely,
\[
 x=1-aC_1+iaS_1,\qquad y=1+aK_1+aH_1,
\]
and the conjugate signs in these formulas give \(x^{-1}\) and \(y^{-1}\).
Thus the derivative coordinate ring is exactly
\(\Qbar[x^{\pm1},y^{\pm1}]\).  Its presentation in the four derivative
variables has precisely the two conic relations
\eqref{eq:trig-jet-relation}--\eqref{eq:hyperbolic-jet-relation}; the
inverse formulas show that there can be no further kernel.  Finally,
Theorem~\ref{thm:fourth-root-orbit} makes \(S,C,H,K\) algebraically
independent over the larger field \(\Eexp\), hence over
\(\Qbar(x,y)\).  Adjoining them is therefore a polynomial extension of the
derivative coordinate ring, and introduces no additional relations.
\end{proof}

\begin{remark}[All positive-order jets]
Repeated differentiation of the four integrands in
\eqref{eq:four-entire-integrals} shows that every derivative of positive
order at \(a\) belongs to \(\Qbar(e^a,e^{ia})\).  Conversely the first
derivatives recover \(e^a,e^{ia}\) by
\eqref{eq:jet-torus-coordinates}.  Hence Theorem
\ref{thm:first-jet-ideal} also says that the four values are algebraically
independent over the field generated by all of their positive-order jets.
\end{remark}

\subsection{A quantitative consequence}

The results above are qualitative.  A recent general theorem supplies
Liouville-type lower bounds, but the exponent is controlled by the
differential transcendence degree, not merely by the number of displayed
integral values.

\begin{proposition}[Quantitative orbit bound]\label{prop:AF-orbit-measure}
Let \(A=\{a_1,\ldots,a_N\}\) be distinct positive real algebraic numbers,
and put
\[
 r=\dim_{\Q}\Span_{\Q}\{a_1,\ldots,a_N\},\qquad
 \tau=4N+2r.
\]
Let \(\boldsymbol\xi_A\) be the \(4N\)-tuple consisting of
\[
 \Si(a_j),\ \Cin(a_j),\ \Shi(a_j),\ \Cinh(a_j)
 \qquad(1\le j\le N).
\]
There exist a constant \(C_2>0\) and positive constants \(C_1(D)>0\),
depending on \(A\) and on the degree bound \(D\), such that every nonzero
polynomial \(P\in\Z[\mathbf X]\) satisfies
\begin{equation}\label{eq:AF-orbit-bound}
 \bigl|P(\boldsymbol\xi_A)\bigr|
 \ge C_1(\deg P)\,
       H(P)^{-C_2(\deg P)^{\tau}}.
\end{equation}
In particular the alternative ``\(P(\boldsymbol\xi_A)=0\)'' in the general
measure theorem never occurs here.
\end{proposition}

\begin{proof}
Consider the \(4N\) E-functions obtained from
\eqref{eq:four-entire-integrals} by replacing their argument by \(a_jz\),
and evaluate them at \(z=1\).  The functional proof of
Theorem~\ref{thm:master}, followed by the blockwise transformation
\eqref{eq:fourth-root-transform}, shows that these functions are
algebraically independent over the exponential differential field generated
by \(e^{\pm a_jz}\) and \(e^{\pm ia_jz}\).  Their derivatives recover that
field, whose torus rank is \(2r\).  Consequently the differential field
generated by the functions and all their derivatives has transcendence
degree \(4N+2r=\tau\) over \(\Qbar(z)\).

Apply the E-function case of
\cite[Theorem~1.3]{AdamczewskiFaverjon2025} at the regular point \(z=1\).
It gives the alternative between vanishing and the lower bound
\eqref{eq:AF-orbit-bound}, with exponent \((\deg P)^\tau\).  Theorem
\ref{thm:pure-exp-base} rules out vanishing for every nonzero \(P\).
\end{proof}

\begin{remark}[Scope of the quantitative import]
Proposition~\ref{prop:AF-orbit-measure} is an application of a general
algebraic-independence measure, not a new measure theorem.  The constants in
\eqref{eq:AF-orbit-bound} are those furnished by
Adamczewski--Faverjon; their paper discusses effectivity separately.  The
same procedure applies to any fixed full-rank linear family supplied by
Lemma~\ref{lem:linear-image-ein}, but its exponent must be computed from the
associated differential field before a numerical value is stated.  No such
quantitative estimate bears on the arithmetic of \(\gamma\) or \(\delta\),
which are connection constants outside the finite-point E-function field.
\end{remark}


\section{Generic-point transfer and concrete algebraic models}
\label{sec:generic-ein-applications}

Theorem~\ref{thm:pure-exp-base} does more than certify isolated special
values.  It provides a concrete transcendence basis on which any finite
algebraic construction may be evaluated without acquiring accidental
relations.  We record this elementary transfer principle before applying it
to Hankel determinants, formal orthogonal polynomials, Grassmannians,
symmetric spectra, and representable matroids.

Throughout this section put $K=\Eexp$.

\begin{proposition}[Generic-point transfer]
\label{prop:generic-point-transfer}
Let $\lambda_1,\ldots,\lambda_s\in\Qbar^\times$ be distinct and put
$y_i=\Ein(\lambda_i)$.  For polynomials
$F_1,\ldots,F_r\in K[X_1,\ldots,X_s]$, the two homomorphisms
\[
 K[T_1,\ldots,T_r]\longrightarrow\C,
 \quad T_j\longmapsto F_j(y_1,\ldots,y_s),
\]
and
\[
 K[T_1,\ldots,T_r]\longrightarrow K[X_1,\ldots,X_s],
 \quad T_j\longmapsto F_j(X_1,\ldots,X_s),
\]
have the same kernel.  The analogous assertion holds for rational maps
after localizing at denominators which do not vanish identically.
\end{proposition}

\begin{proof}
Theorem~\ref{thm:pure-exp-base} says precisely that the evaluation map
$K[X_1,\ldots,X_s]\to\C$, $X_i\mapsto y_i$, is injective.  Composing the
second displayed homomorphism with this injection proves the assertion.
Localization gives the rational version.
\end{proof}

Thus the point $y=(y_1,\ldots,y_s)$ is a generic point of affine
$s$-space over $K$.  The proposition is classical generic-point algebra;
the arithmetic content is the explicit realization of the coordinates as
finite-point $E$-values.

\subsection{A Hankel--Favard transcendence basis}

Put
\begin{equation}\label{eq:ein-moments}
 \mu_m=\Ein(m+1),\qquad
 \Delta_n=\det(\mu_{i+j})_{0\le i,j\le n},
 \qquad \Delta_{-1}=1.
\end{equation}

\begin{theorem}[Hankel determinants and formal Jacobi parameters]
\label{thm:ein-hankel-jacobi}
The countable family
\[
 \Delta_0,\Delta_1,\Delta_2,\ldots
\]
is algebraically independent over $K$, in the finite-subset sense.

Moreover the functional $\mathcal L\in\C[x]^*$ defined by
$\mathcal L(x^m)=\mu_m$ is quasi-definite.  Let its monic formal orthogonal
polynomials be normalized by
\begin{equation}\label{eq:ein-favard-recurrence}
 p_{n+1}(x)=(x-a_n)p_n(x)-b_np_{n-1}(x),
 \qquad p_{-1}=0,\quad p_0=1.
\end{equation}
Then
\[
 \mu_0,a_0,a_1,\ldots,b_1,b_2,\ldots
\]
is algebraically independent over $K$.
\end{theorem}

\begin{proof}
Expansion of $\Delta_n$ at its lower-right entry gives
\begin{equation}\label{eq:hankel-triangular-expansion}
 \Delta_n=\Delta_{n-1}\mu_{2n}
 +R_n(\mu_0,\ldots,\mu_{2n-1})
\end{equation}
for a polynomial $R_n$.  The coefficient $\Delta_{n-1}$ is a nonzero
polynomial in the preceding independent moments.  Hence $\Delta_n$ is
transcendental over $K(\Delta_0,\ldots,\Delta_{n-1})$, and induction proves
the first assertion.  In particular every $\Delta_n$ is nonzero, which is
the quasi-definiteness criterion.

Finite Gram--Schmidt elimination expresses
\[
 (\mu_0,\ldots,\mu_{2N+1})
 \longmapsto
 (\mu_0,a_0,\ldots,a_N,b_1,\ldots,b_N)
\]
rationally.  Conversely, the matrix identity
$\mu_m=\mu_0(J^m)_{00}$ for the finite Jacobi truncation, equivalently the
weighted Motzkin-path expansion, expresses the moments rationally (indeed
polynomially after adjoining $\mu_0$) in the displayed Favard parameters.
The two lists both contain $2N+2$ entries and are birational.  Since the
moment list is algebraically independent, so is the parameter list.  Letting
$N$ vary proves the countable statement.
\end{proof}

The system in Theorem~\ref{thm:ein-hankel-jacobi} is formal rather than
positive-definite: for example $\Delta_1<0$.  It therefore does not define
a self-adjoint Jacobi operator.  Also the Hankel and Favard families must
not be joined into one independent family, since
\begin{equation}\label{eq:hankel-favard-dependence}
 b_n\Delta_{n-1}^2=\Delta_n\Delta_{n-2}.
\end{equation}

\subsection{Pl\"ucker coordinates, spectra, and matroids}

\begin{corollary}[Explicit generic Pl\"ucker data]
\label{cor:ein-pluecker}
Let $1\le r<n$, choose pairwise distinct nonzero algebraic numbers
$\lambda_{ij}$, and form the $r\times n$ matrix
$M=(\Ein(\lambda_{ij}))$.  If $p_I$ denotes its maximal minors, then the
complete relation ideal of the $p_I$ over $K$ is the classical Pl\"ucker
ideal.  Consequently
\[
 \trdeg_K K(p_I)=r(n-r)+1.
\]
The corresponding projective Pl\"ucker ratios have transcendence degree
$r(n-r)$.
\end{corollary}

\begin{proof}
Apply Proposition~\ref{prop:generic-point-transfer} to the maximal-minor
map on a generic $r\times n$ matrix.  Its kernel is the Pl\"ucker ideal and
its image is the affine cone over $\operatorname{Gr}(r,n)$; see, for
example, \cite{BrunsConcaVarbaro2013}.  The dimension statements are the
dimensions of that cone and its projectivization.
\end{proof}

\begin{theorem}[A simple real algebraically independent spectrum]
\label{thm:ein-symmetric-spectrum}
For $1\le i\le j\le n$, choose pairwise distinct positive algebraic
$\lambda_{ij}$ and define the real symmetric matrix
\[
 A_{ij}=A_{ji}=\Ein(\lambda_{ij}).
\]
Then $A$ has $n$ distinct real eigenvalues
$\rho_1,\ldots,\rho_n$, and these eigenvalues are algebraically independent
over $K$.
\end{theorem}

\begin{proof}
The characteristic coefficients of a generic symmetric matrix are
algebraically independent: restriction to diagonal matrices gives the
elementary symmetric polynomials in $n$ independent variables.  The
discriminant is likewise a nonzero polynomial, as diagonal matrices with
distinct entries show.  Proposition~\ref{prop:generic-point-transfer}
therefore makes the characteristic coefficients of $A$ independent and
its discriminant nonzero.  Symmetry makes all roots real, while the
nonzero discriminant makes them distinct.  The splitting field is finite
over the field generated by the $n$ characteristic coefficients, so
\[
 \trdeg_K K(\rho_1,\ldots,\rho_n)=n.
\]
As there are $n$ eigenvalues, they are jointly algebraically independent.
\end{proof}

\begin{corollary}[Representable matroids as $\Ein$-period matroids]
\label{cor:ein-representable-matroid}
Every finite matroid representable over $\Qbar$ occurs as the algebraic
matroid over $K$ of explicit algebraic linear combinations of
$\Ein$-values.  More precisely, if a rank-$m$ representation has columns
$b_e=(b_{1e},\ldots,b_{me})^t$ and
\[
 W_e=\sum_{i=1}^m b_{ie}\Ein(\lambda_i)
\]
for distinct $\lambda_i\in\Qbar^\times$, then for every subset $S$ of the
ground set
\[
 \trdeg_K K(W_e:e\in S)=\operatorname{rank}(b_e:e\in S).
\]
The complete relation ideal is generated by the linear forms
$\sum_e c_eT_e$ with $\sum_e c_eb_e=0$.
\end{corollary}

\begin{proof}
This is Lemma~\ref{lem:linear-image-ein} applied to the representing
matrix and to each column subfamily.
\end{proof}

These corollaries do not introduce new Grassmannian, spectral, or matroid
machinery: each ambient construction is classical.  They give exact
arithmetic realizations by special values and solve no named external open
problem.  The historical priority of these particular $\Ein$ realizations
should be checked separately before publication.


\section{Positive \texorpdfstring{$\Ein$}{Ein} moment systems and
algebraically independent Jacobi spectra}
\label{sec:positive-ein-jacobi}

The generic Hankel system built directly from the moments
$\Ein(m+1)$ is quasi-definite but not positive-definite.  A finite
difference removes this defect and produces a genuine Hausdorff moment
problem without sacrificing any algebraic independence.

Fix $a,b\in\overline{\mathbb Q}_{>0}$ and an integer $r\geq1$.  For a
sequence $f=(f_m)_{m\geq0}$ write
\[
 (\mathsf Df)_m=f_{m+1}-f_m,
\]
and put
\begin{equation}\label{eq:positive-ein-moments}
 f_m=\Ein(b+am),\qquad
 \mu_m^{(a,b,r)}=(-1)^{r-1}(\mathsf D^rf)_m\quad(m\geq0).
\end{equation}
When there is no danger of confusion, write simply $\mu_m$.

\begin{theorem}[Positive finite-difference $\Ein$ moment theorem]
\label{thm:positive-ein-jacobi}
Let $K=\Eexp$.  The moments in
\eqref{eq:positive-ein-moments} have the exact representation
\begin{align}
 \mu_m
 &=\int_0^1 e^{-(b+am)t}\frac{(1-e^{-at})^r}{t}\,dt
 \label{eq:positive-ein-laplace}\\
 &=\int_{e^{-a}}^1 x^m w_{a,b,r}(x)\,dx,
 \qquad
 w_{a,b,r}(x)=
 \frac{x^{b/a-1}(1-x)^r}{-\log x}.
 \label{eq:positive-ein-weight}
\end{align}
In particular $(\mu_m)_{m\geq0}$ is a strictly positive-definite
Hausdorff moment sequence.  Moreover:
\begin{enumerate}
\item The countable family $\mu_0,\mu_1,\ldots$ is algebraically
independent over $K$, in the finite-subset sense.  The infinite Hankel
matrix $(\mu_{i+j})_{i,j\geq0}$ is strictly totally positive: every one
of its finite minors is positive.
\item With
\begin{equation}\label{eq:positive-ein-Hankel}
 \mathcal H_n=\det(\mu_{i+j})_{0\leq i,j\leq n},
 \qquad \mathcal H_{-1}=1,
\end{equation}
one has
\begin{equation}\label{eq:positive-ein-Heine}
 \mathcal H_n=
 \frac1{(n+1)!}
 \int_{[e^{-a},1]^{n+1}}
 \prod_{0\leq i<j\leq n}(x_i-x_j)^2
 \prod_{j=0}^n w_{a,b,r}(x_j)\,dx_j>0.
\end{equation}
The family $\mathcal H_0,\mathcal H_1,\ldots$ is algebraically
independent over $K$, and its exact capacity scale is
\begin{equation}\label{eq:positive-ein-Hankel-capacity}
 \lim_{n\to\infty}\mathcal H_n^{1/n^2}
 =\operatorname{cap}([e^{-a},1])=\frac{1-e^{-a}}4.
\end{equation}
\item Let $p_{-1}=0$, $p_0=1$ be the monic orthogonal polynomials for
the functional $\mathcal L(x^m)=\mu_m$, and normalize the Favard
recurrence by
\begin{equation}\label{eq:positive-ein-Favard}
 p_{n+1}(x)=(x-\alpha_n)p_n(x)-\beta_np_{n-1}(x).
\end{equation}
Then $\beta_n>0$ for every $n\geq1$, and
\[
 \mu_0,\alpha_0,\alpha_1,\ldots,
 \beta_1,\beta_2,\ldots
\]
is algebraically independent over $K$.
\item The Jacobi matrix
\begin{equation}\label{eq:positive-ein-Jacobi}
 J=
 \begin{pmatrix}
 \alpha_0&\sqrt{\beta_1}&&\\
 \sqrt{\beta_1}&\alpha_1&\sqrt{\beta_2}&\\
 &\sqrt{\beta_2}&\alpha_2&\ddots\\
 &&\ddots&\ddots
 \end{pmatrix}
\end{equation}
defines a bounded self-adjoint operator on $\ell^2(\mathbb N_0)$.
It has simple purely absolutely continuous spectrum
\[
 \sigma(J)=[e^{-a},1]
\]
of multiplicity one.
Its monic recurrence coefficients lie in the Nevai class of this interval:
\begin{equation}\label{eq:positive-ein-Nevai}
 \alpha_n\longrightarrow\frac{1+e^{-a}}2,
 \qquad
 \beta_n\longrightarrow\frac{(1-e^{-a})^2}{16}.
\end{equation}
\item For every $N\geq1$, the $N$ eigenvalues of the leading
$N\times N$ principal truncation $J_N$ are distinct, belong to
$(e^{-a},1)$, and are jointly algebraically independent over $K$.
If they are denoted in increasing order by
$\lambda_{N,1},\ldots,\lambda_{N,N}$, then their normalized counting
measures satisfy
\begin{equation}\label{eq:positive-ein-arcsine}
 \frac1N\sum_{j=1}^N\delta_{\lambda_{N,j}}
 \xrightarrow[N\to\infty]{\mathrm{weak}}
 \frac{\mathbf 1_{(e^{-a},1)}(x)}
 {\pi\sqrt{(x-e^{-a})(1-x)}}\,dx.
\end{equation}
\end{enumerate}
\end{theorem}

\begin{proof}
The standard integral
\[
 \Ein(s)=\int_0^1\frac{1-e^{-st}}t\,dt\qquad(s>0)
\]
shows, after taking the $r$th forward difference in $m$, that
\[
 (-1)^{r-1}(\mathsf D^rf)_m
 =\int_0^1e^{-(b+am)t}\frac{(1-e^{-at})^r}{t}\,dt.
\]
The substitution $x=e^{-at}$ gives
\eqref{eq:positive-ein-weight}.  The density is positive on
$(e^{-a},1)$.  At $x=1$ it is asymptotic to
$x^{b/a-1}(1-x)^{r-1}$, while the other endpoint is strictly positive;
hence the density is integrable and has essential support
$[e^{-a},1]$.  The Gram determinant formula now gives
\eqref{eq:positive-ein-Heine}; its integrand is positive off the
diagonals, so every $\mathcal H_n$ is strictly positive.

More generally, choose strictly increasing nonnegative integer lists
$i_1<\cdots<i_k$ and $j_1<\cdots<j_k$.  Andreief's identity writes the
corresponding Hankel minor as
\[
 \frac1{k!}\int_{[e^{-a},1]^k}
 \det(x_\ell^{i_p})_{p,\ell=1}^k
 \det(x_\ell^{j_p})_{p,\ell=1}^k
 \prod_{\ell=1}^kw_{a,b,r}(x_\ell)\,dx_\ell.
\]
On the ordered simplex $e^{-a}<x_1<\cdots<x_k<1$, each generalized
Vandermonde determinant is a positive Vandermonde times a Schur
polynomial with positive values.  Their product is positive on a set of
positive measure, proving strict total positivity.

For algebraic independence, let
\[
 y_m=\Ein(b+am).
\]
The points $b+am$ are distinct nonzero algebraic numbers, so the
multipoint $\Ein$ theorem makes every finite subfamily of the $y_m$
algebraically independent over $K$.  For fixed $M$, the $M+1$ linear
forms $\mu_0,\ldots,\mu_M$ in
$y_0,\ldots,y_{M+r}$ are linearly independent: the coefficient of
$y_{M+r}$ occurs only in $\mu_M$, and descending induction finishes the
argument.  Extending these coefficient vectors to a basis gives an
invertible linear change of variables.  Thus
$\mu_0,\ldots,\mu_M$ are algebraically independent over $K$.

Expansion of the determinant at its lower-right entry gives
\begin{equation}\label{eq:positive-ein-Hankel-triangular}
 \mathcal H_n=\mathcal H_{n-1}\mu_{2n}
 +R_n(\mu_0,\ldots,\mu_{2n-1}).
\end{equation}
Induction in $n$, viewing a putative relation as a polynomial in
$\mu_{2n}$, proves the algebraic independence of the $\mathcal H_n$.

Strict positivity of all Gram determinants gives the monic Favard
recurrence, with
\begin{equation}\label{eq:positive-ein-beta-Hankel}
 \beta_n=\frac{\mathcal H_n\mathcal H_{n-2}}
 {\mathcal H_{n-1}^2}>0.
\end{equation}
For each $N\geq0$, finite Gram--Schmidt elimination gives a rational map
\[
 (\mu_0,\ldots,\mu_{2N+1})
 \longmapsto
 (\mu_0,\alpha_0,\ldots,\alpha_N,
             \beta_1,\ldots,\beta_N).
\]
Conversely, the weighted Motzkin-path expansion, equivalently
$\mu_m=\mu_0(J^m)_{00}$, recovers the moments through order $2N+1$
polynomially from the displayed Favard variables.  The two lists have
$2N+2$ entries and the transformations are inverse on the open set of
nonzero Hankel determinants.  They are therefore birational.  Algebraic
independence of the moments proves that of the Favard variables.

Normalize the measure in \eqref{eq:positive-ein-weight} by its mass
$\mu_0$.  Polynomials are dense in its $L^2$ space, and the orthonormal
polynomial transform identifies $J$ with multiplication by $x$ on
\[
 L^2\!\left([e^{-a},1],\mu_0^{-1}w_{a,b,r}(x)\,dx\right).
\]
This proves bounded self-adjointness.  Since the density is positive
almost everywhere on the whole interval, multiplication by $x$ has
purely absolutely continuous spectrum $[e^{-a},1]$.  The constant
function is cyclic, so the spectral multiplicity is one.  The
Denisov--Rakhmanov theorem~\cite{Denisov2004}, affinely rescaled from $[-2,2]$ to
$[e^{-a},1]$, gives
\[
 \alpha_n\to\frac{1+e^{-a}}2,
 \qquad
 \sqrt{\beta_n}\to\frac{1-e^{-a}}4,
\]
which is exactly \eqref{eq:positive-ein-Nevai} in the monic
normalization \eqref{eq:positive-ein-Favard}.
If $h_n=\mathcal L(p_n^2)$, then
$h_n/h_{n-1}=\beta_n$ and
$\mathcal H_n=\prod_{j=0}^nh_j$.  Consequently
\[
 \mathcal H_n
 =\mu_0^{n+1}\prod_{k=1}^n\beta_k^{n+1-k}.
\]
Taking logarithms and using \eqref{eq:positive-ein-Nevai} proves
\eqref{eq:positive-ein-Hankel-capacity}.

It remains to prove the arithmetic assertion about finite truncations.
The positivity of the off-diagonal entries makes every $J_N$ an
irreducible real symmetric tridiagonal matrix.  Hence its eigenvalues are
simple; the standard zero theorem for orthogonal polynomials places them
in $(e^{-a},1)$.  Write the characteristic polynomial as
\[
 \det(TI_N-J_N)=T^N+c_1T^{N-1}+\cdots+c_N.
\]
Each $c_j$ is a polynomial over $\mathbb Q$ in
$\alpha_0,\ldots,\alpha_{N-1},\beta_1,\ldots,\beta_{N-1}$.
These $N$ polynomials are algebraically independent: after formally
specializing all $\beta_i$ to zero, they become the signed elementary
symmetric polynomials in the independent diagonal variables
$\alpha_0,\ldots,\alpha_{N-1}$.  The algebraic independence of the
actual Favard variables makes their evaluation map injective, so the
actual $c_1,\ldots,c_N$ are algebraically independent over $K$.

If $\lambda_{N,1},\ldots,\lambda_{N,N}$ are the eigenvalues, then the
$c_j$ are their elementary symmetric functions and the splitting field
is finite over $K(c_1,\ldots,c_N)$.  Consequently
\[
 \operatorname{trdeg}_K
 K(\lambda_{N,1},\ldots,\lambda_{N,N})=N.
\]
There are exactly $N$ eigenvalues, so they are jointly algebraically
independent over $K$.

Finally, these eigenvalues are the zeros of $p_N$.  Formula
\eqref{eq:positive-ein-Nevai} gives their limiting distribution directly.
Indeed, put
\[
 c=\frac{1+e^{-a}}2,
 \qquad d=\frac{1-e^{-a}}4.
\]
For every fixed integer $k\geq0$, a closed-walk expansion of the trace,
with the $O(k)$ boundary rows discarded, gives
\[
 \lim_{N\to\infty}\frac1N\operatorname{tr}(J_N^k)
 =\frac1{2\pi}\int_0^{2\pi}(c+2d\cos\theta)^k\,d\theta.
\]
The matrices $J_N$ are uniformly bounded, so polynomial approximation
extends this convergence from moments to every continuous test function.
The push-forward of $d\theta/(2\pi)$ under
$\theta\mapsto c+2d\cos\theta$ is precisely the probability measure on
the right side of \eqref{eq:positive-ein-arcsine}.  This proves the weak
convergence.  Equivalently, it is the classical zero-counting theorem for
a regular measure on a single interval.
\end{proof}

\begin{remark}[What is new and what is imported]
The Hausdorff moment representation and the spectral consequences are
classical once \eqref{eq:positive-ein-weight} is known.  The arithmetic
content is that the explicit moments, all Hankel determinants, all
Favard coordinates, and every finite Jacobi spectrum acquire the stated
algebraic independence from the multipoint $\Ein$ theorem.  Historical
priority for this exact arithmetic realization should be audited before
publication.
\end{remark}


\providecommand{\EinDepth}[1]{\operatorname{Ein}_{#1}}
\providecommand{\Eexp}{\mathcal E_{\mathrm{exp}}}

\section{Multipoint one-level polyexponentials of arbitrary weight}
\label{sec:ein-depth-tower}

For \(r\geq1\), define the order-\(r\) complementary exponential
integral by
\begin{equation}\label{eq:ein-depth-definition}
 \EinDepth{r}(z)
 :=\sum_{n\geq1}\frac{(-1)^{n-1}z^n}{n^r n!}.
\end{equation}
We call \(r\) the \emph{order} (or weight), not the level: these are
classical one-level polyexponentials, not the multiple-index
polyexponentials of level greater than one
\cite{Hardy1905PartII,Boyadzhiev2007,AminovArnaudo2026}.
Thus \(\EinDepth{1}=\Ein\), and
\begin{equation}\label{eq:ein-depth-chain}
 z\EinDepth{1}'(z)=1-e^{-z},\qquad
 z\EinDepth{r}'(z)=\EinDepth{r-1}(z)\quad(r\geq2).
\end{equation}
We use the harmless notation
\(\EinDepth{0}(z)=1-e^{-z}\)
only when writing the chain compactly.  For each fixed \(r\),
\(\EinDepth{r}\) is an \(E\)-function.  Indeed, its coefficient sequence
is hypergeometric, hence holonomic.  Moreover, for every \(N\), the single
integer \(\operatorname{lcm}(1,\ldots,N)^r\) clears simultaneously all
coefficients \(n![z^n]\EinDepth{r}(z)\), \(0\leq n\leq N\); it grows at
most exponentially in \(N\).  The coefficients are rational of absolute
value at most one, so the conjugate-growth condition is automatic.

We first record the rational-denominator step needed in the iterated
primitive argument.  It is useful beyond the present family.

\begin{lemma}[Affine form of a primitive rational function]
\label{lem:depth-affine-primitive}
Let \(F\) be a differential field of characteristic zero whose constant
field \(C\) is algebraically closed.  Let
\(a_1,\ldots,a_s\in F\) have the property
\begin{equation}\label{eq:depth-no-base-exact-combination}
 \sum_{i=1}^s c_i a_i\in D(F),\quad c_i\in C
 \qquad\Longrightarrow\qquad c_1=\cdots=c_s=0.
\end{equation}
Adjoin algebraically independent primitives
\(u_1,\ldots,u_s\), with \(Du_i=a_i\), and put
\(L=F(u_1,\ldots,u_s)\).  If \(g\in L\) and \(Dg\in F\), then
\begin{equation}\label{eq:depth-affine-conclusion}
 g=f+\sum_{i=1}^s b_i u_i
 \qquad(f\in F,\ b_i\in C).
\end{equation}
In particular, the constant field of \(L\) is again \(C\).
\end{lemma}

\begin{proof}
Work first in \(F[u_1,\ldots,u_s]\).  There is no nonconstant
irreducible Darboux polynomial in this ring.  Indeed, suppose that an
irreducible nonconstant \(P\) divides \(DP\).  Degree in the \(u_i\)
shows that \(DP=hP\) for some \(h\in F\).  Fix a monomial order refining
total degree and divide \(P\) by the coefficient of its leading monomial.
For the resulting monic polynomial \(Q\), one has
\(DQ=(h-Dp/p)Q\), where \(p\) was that leading coefficient.  Since the
derivation lowers the \(u\)-degree of the monomial part and the leading
coefficient of \(Q\) is one, comparison of the leading monomial forces
\(h-Dp/p=0\); hence \(DQ=0\).
Write \(Q_n\) for its highest nonzero homogeneous part.  The degree-\(n\)
part of \(DQ=0\) shows that every coefficient of \(Q_n\) belongs to \(C\).
The degree-\((n-1)\) part is
\[
 DQ_{n-1}+\sum_{i=1}^s
 a_i\frac{\partial Q_n}{\partial u_i}=0.
\]
Since \(Q_n\) is nonconstant, at least one coefficient comparison in this
identity writes a nonzero constant linear combination of the \(a_i\) as
an element of \(D(F)\), contrary to
\eqref{eq:depth-no-base-exact-combination}.

Now write \(g=A/B\) in lowest terms.  If a nonconstant irreducible factor
\(P\) occurs in \(B\), put \(m=-v_P(g)>0\).  If \(P\nmid DP\), the usual
quotient-rule computation in the discrete valuation \(v_P\) gives
\(v_P(Dg)=-m-1\): the term obtained by differentiating \(P^{-m}\) has
that valuation and cannot cancel with terms of valuation at least \(-m\).
But every nonzero element of \(F\) has \(P\)-valuation zero.  Thus
\(Dg\in F\) forces \(P\mid DP\), which the preceding paragraph rules out.
Hence \(g\) is a polynomial in the \(u_i\).

Let \(g_n\) be its highest homogeneous part.  Since \(Dg\in F\), the
degree-\(n\) coefficient comparison makes the coefficients of \(g_n\)
constant.  If \(n\geq2\), the degree-\((n-1)\) comparison again makes a
nonzero constant linear combination of the \(a_i\) exact in \(F\), a
contradiction.  Thus \(n\leq1\), and
\eqref{eq:depth-affine-conclusion} follows.  If \(Dg=0\), then
\(Df+\sum_i b_i a_i=0\);
\eqref{eq:depth-no-base-exact-combination} gives \(b_i=0\), and then
\(f\in C\).  Thus \(L\) has no new constants.
\end{proof}

We now apply this lemma to the exponential torus used in
Lemma~\ref{lem:exact}.  Let
\(\lambda_1,\ldots,\lambda_s\in\Qbar^\times\) be distinct, let
\(\nu_1,\ldots,\nu_d\) be a \(\mathbb Z\)-basis of the additive group
generated by the \(\lambda_i\), and write
\(\lambda_i=m_i\mathbin{\cdot}\nu\), with
\(m_i\in\mathbb Z^d\).  Put
\[
 K=\operatorname{Frac}
 \bigl(\Qbar(z)[X_1^{\pm1},\ldots,X_d^{\pm1}]\bigr),
 \qquad DX_j=-\nu_jX_j.
\]
Because the \(\nu_j\) form a basis of the group generated by the
\(\lambda_i\), each \(\nu_j\) is also an integral combination of the
\(\lambda_i\).  Thus passing from \(\lambda\) to \(\nu\) changes neither
the generated exponential field nor its character lattice.  The vectors
\(m_i\) are pairwise distinct and nonzero.
The constant field of \(K\) is \(\Qbar\).  Lemma~\ref{lem:exact}
implies both
\begin{equation}\label{eq:depth-base-nonexact}
 \frac1z\notin D(K)
\end{equation}
and
\begin{equation}\label{eq:depth-first-layer-no-combination}
 \sum_{i=1}^s c_i\frac{1-X^{m_i}}z\in D(K),\quad c_i\in\Qbar
 \qquad\Longrightarrow\qquad c_1=\cdots=c_s=0.
\end{equation}
Indeed, after collecting the distinct Laurent characters, each assertion
is an immediate instance of that lemma.

All the functions below have Taylor coefficients in \(\Qbar\), and the
field they generate at the origin embeds in \(\Qbar((z))\).  Since the
constant field of \(\Qbar((z))\) under \(d/dz\) is \(\Qbar\), the ambient
extension used in every invocation of Kolchin--Ostrowski has no new
constants.  The affine lemma below also verifies this internally after each
primitive layer is adjoined.

For \(q\geq1\), set
\[
 Y_{i,q}(z)=\EinDepth{q}(\lambda_i z),
 \qquad
 K_q=K(Y_{i,j}:1\leq i\leq s,\ 1\leq j\leq q).
\]
Then
\begin{equation}\label{eq:depth-tower-derivatives}
 DY_{i,1}=\frac{1-X^{m_i}}z,
 \qquad
 DY_{i,q}=\frac{Y_{i,q-1}}z\quad(q\geq2).
\end{equation}

\begin{lemma}[Nonexactness throughout the tower]
\label{lem:depth-one-over-z-nonexact}
For every \(q\geq0\), with \(K_0=K\), the constant field of \(K_q\) is
\(\Qbar\), and
\begin{equation}\label{eq:depth-one-over-z-all-levels}
 \frac1z\notin D(K_q).
\end{equation}
Moreover, at each level the variables
\(Y_{1,q},\ldots,Y_{s,q}\) are algebraically independent over
\(K_{q-1}\).
\end{lemma}

\begin{proof}
The primitive form of the Kolchin--Ostrowski theorem~\cite{Kolchin} says that primitives
with derivatives in a differential field with algebraically closed
constants are algebraically dependent only if a nonzero constant linear
combination of their derivatives is exact in the base field.
Equation~\eqref{eq:depth-first-layer-no-combination} therefore makes the
first layer algebraically independent over \(K\).  Applying
Lemma~\ref{lem:depth-affine-primitive} also shows that \(K_1\) has no new
constants.

Suppose \(Dg=1/z\) with \(g\in K_1\).  Since \(Dg\in K\),
Lemma~\ref{lem:depth-affine-primitive} gives
\(g=k+\sum_i b_iY_{i,1}\), where \(k\in K\) and
\(b_i\in\Qbar\).  Hence
\[
 Dk=\frac{1-\sum_i b_i+\sum_i b_iX^{m_i}}z,
\]
contradicting Lemma~\ref{lem:exact}.  This proves
\eqref{eq:depth-one-over-z-all-levels} for \(q=1\).

Assume the assertions through level \(q\geq1\).  The extension
\(K_q/K_{q-1}\) admits all translations
\[
 \tau_h(Y_{i,q})=Y_{i,q}+h_i,\qquad
 h=(h_1,\ldots,h_s)\in\Qbar^s,
\]
as differential automorphisms.  If, for some \(g\in K_q\),
\[
 Dg=\sum_{i=1}^s c_i\frac{Y_{i,q}}z
\]
with \(c\ne0\), choose \(h\) with \(c\mathbin{\cdot}h\ne0\).  Then
\[
 D(\tau_hg-g)=\frac{c\mathbin{\cdot}h}{z},
\]
contradicting the inductive nonexactness of \(1/z\) in \(K_q\).
Thus no nonzero constant linear combination of the derivatives of the
next layer is exact.  Kolchin--Ostrowski gives the algebraic independence
of \(Y_{1,q+1},\ldots,Y_{s,q+1}\) over \(K_q\), and
Lemma~\ref{lem:depth-affine-primitive} shows that the constants remain
\(\Qbar\).

Finally, suppose \(Dg=1/z\) with \(g\in K_{q+1}\).  The affine lemma gives
\[
 g=k+\sum_i b_iY_{i,q+1},\qquad
 k\in K_q,\quad b_i\in\Qbar,
\]
and therefore
\[
 Dk=\frac{1-\sum_i b_iY_{i,q}}z.
\]
If \(b\ne0\), translate the \(q\)-th layer by an \(h\) with
\(b\mathbin{\cdot}h\ne0\); the difference makes a nonzero multiple of
\(1/z\) exact in \(K_q\), a contradiction.  If \(b=0\), the displayed
identity directly contradicts the inductive hypothesis.  This completes
the induction.
\end{proof}

For clarity, the first new case can also be checked without the induction.
For one nonzero algebraic \(\lambda\), put
\(Y=\Ein(\lambda z)\) and
\(Z=\EinDepth{2}(\lambda z)\).  Exactness in \(K\) first shows that
\(Y\) is transcendental over \(K\); the affine lemma then gives constants
\(\Qbar\) in \(K(Y)\), and its displayed calculation proves
\(1/z\notin D(K(Y))\).  If \(Z\) were algebraic over \(K(Y)\), the
one-primitive Kolchin--Ostrowski theorem would give
\(Dg=Y/z\) for some \(g\in K(Y)\), after rescaling a nonzero constant.
Because \(Y\) is transcendental, \(\tau_h(Y)=Y+h\) is a differential
automorphism of \(K(Y)/K\).  Applying it with \(h\ne0\) gives
\(D(\tau_hg-g)=h/z\), the required contradiction.  Thus
\(e^{\lambda z},Y,Z\) have functional transcendence degree three when
\(s=1\), and the order-two claim is independent of any higher-order
induction.

\begin{theorem}[Multipoint resonant polyexponential theorem]
\label{thm:ein-depth-multipoint}
Let \(\lambda_1,\ldots,\lambda_s\in\Qbar^\times\) be distinct, let
\[
 d=\dim_{\Q}\operatorname{span}_{\Q}
   \{\lambda_1,\ldots,\lambda_s\},
\]
and let \(R\geq1\).  Then
\begin{equation}\label{eq:ein-depth-value-trdeg}
 \operatorname{trdeg}_{\Qbar}
 \Qbar\bigl(
   e^{\lambda_i},\EinDepth{r}(\lambda_i):
   1\leq i\leq s,\ 1\leq r\leq R
 \bigr)
 =d+Rs.
\end{equation}
In particular, the \(Rs\) polyexponential values are algebraically independent over
\(\Qbar(e^{\lambda_1},\ldots,e^{\lambda_s})\).  All algebraic relations
among the displayed exponential and polyexponential coordinates are therefore the
toric relations among the exponential coordinates.
\end{theorem}

\begin{proof}
Lemma~\ref{lem:depth-one-over-z-nonexact} gives functional transcendence
degree \(d+Rs\) over \(\Qbar(z)\): the exponential torus contributes
\(d\), and each of the \(R\) primitive layers contributes \(s\).

The same functional transcendence degree holds over \(\C(z)\).  Indeed,
all Taylor coefficients are algebraic; after clearing rational-function
denominators, any finite relation over \(\C(z)\) would be a finite linear
dependence over \(\C\) among series with coefficients in \(\Qbar\), and
rank invariance under the scalar extension \(\Qbar\subset\C\) would produce
a relation over \(\Qbar(z)\).

For specialization, include in one vector the constant function \(1\),
the finitely many exponentials \(e^{\pm\lambda_i z}\), and all
\(\EinDepth{r}(\lambda_i z)\), \(1\leq r\leq R\).  The exponentials are
\(E\)-functions.  If \(q_i\in\mathbb Z_{>0}\) clears the algebraic
denominator of \(\lambda_i\), then for every \(N\) the integer
\[
 q_i^N\operatorname{lcm}(1,\ldots,N)^r
\]
clears simultaneously the first \(N\) normalized coefficients of
\(\EinDepth{r}(\lambda_i z)\); all of their conjugates have size at most
\(C_i^N\) for a fixed \(C_i\).  Thus these scaled order functions satisfy
the denominator and conjugate-growth axioms for \(E\)-functions.  By
\eqref{eq:ein-depth-chain}, the displayed vector satisfies a first-order
linear system over \(\Qbar(z)\), with diagonal constant entries for the
exponentials and only \(1/z\) entries in the order chains.  Consequently
its coefficient matrix has no finite pole except possibly \(z=0\), so the
nonzero algebraic point \(z=1\) is ordinary.  Beukers'
refined Siegel--Shidlovskii theorem~\cite{Beukers2006} transfers the functional
transcendence degree, and all relations, to the values at \(z=1\).  This
proves \eqref{eq:ein-depth-value-trdeg}; the final statement follows from
the algebraic independence of all polyexponential coordinates over the exponential
field.
\end{proof}

\begin{corollary}[Pure exponential base change in every order]
\label{cor:ein-depth-pure-exponential}
The countable family
\[
 \bigl\{\EinDepth{r}(\lambda):
      \lambda\in\Qbar^\times,\ r\geq1\bigr\}
\]
is algebraically independent over
\(\Eexp=\Qbar(e^\alpha:\alpha\in\Qbar)\), in the finite-subset sense.
\end{corollary}

\begin{proof}
Any proposed polynomial relation uses finitely many polyexponential values and
finitely many coefficients from \(\Eexp\).  Enlarge the finite set of
algebraic points to include the exponents occurring in those coefficients,
take the largest order appearing, clear coefficient denominators, and
apply Theorem~\ref{thm:ein-depth-multipoint}.  The selected order values
remain algebraically independent over the enlarged exponential field.
\end{proof}

\begin{corollary}[Hausdorff moments and consecutive finite spectra]
\label{cor:ein-depth-hausdorff-spectra}
Fix \(a,b\in\Qbar\cap\mathbb R_{>0}\) and an integer \(\ell\geq1\), and
write \(\Delta_a f(s)=f(s+a)-f(s)\).  For \(q\geq1\) and \(m\geq0\), put
\begin{equation}\label{eq:ein-depth-finite-difference-moments}
 \mu_m^{(q)}
 :=(-1)^{\ell-1}(\Delta_a^\ell\EinDepth{q})(b+am).
\end{equation}
Then \((\mu_m^{(q)})_{m\geq0}\) is a strict Hausdorff moment sequence on
\([e^{-a},1]\), with
\begin{equation}\label{eq:ein-depth-hausdorff-density}
 \mu_m^{(q)}=\int_{e^{-a}}^1x^m w_q(x)\,dx,
 \qquad
 w_q(x)=
 \frac{x^{b/a-1}(1-x)^\ell}{(q-1)!(-\log x)}
 \left(\log\frac{a}{-\log x}\right)^{q-1}.
\end{equation}
The entire countable family
\begin{equation}\label{eq:ein-depth-all-moments-independent}
 \{\mu_m^{(q)}:q\geq1,\ m\geq0\}
\end{equation}
is algebraically independent over \(\Eexp\).

For each \(q\), let
\[
 \mathcal H_n^{(q)}
  =\det(\mu_{i+j}^{(q)})_{0\leq i,j\leq n},
 \qquad \mathcal H_{-1}^{(q)}=1.
\]
All these Hankel determinants, jointly for \(q\geq1\) and \(n\geq0\),
are algebraically independent over \(\Eexp\), and they are strictly
positive.  Let the monic orthogonal polynomials for \(w_q(x)\,dx\) satisfy
\begin{equation}\label{eq:ein-depth-favard-recurrence}
 P_{n+1}^{(q)}(x)
 =(x-\alpha_n^{(q)})P_n^{(q)}(x)
  -\beta_n^{(q)}P_{n-1}^{(q)}(x),
 \qquad \beta_n^{(q)}>0\quad(n\geq1).
\end{equation}
Then the family
\begin{equation}\label{eq:ein-depth-favard-independent}
 \{\mu_0^{(q)},\alpha_n^{(q)},\beta_{n+1}^{(q)}:
        q\geq1,\ n\geq0\}
\end{equation}
is algebraically independent over \(\Eexp\).

For every fixed \(q\), the infinite Hankel matrix
\((\mu_{i+j}^{(q)})_{i,j\geq0}\) is strictly totally positive.  The
bounded Jacobi operator associated with
\eqref{eq:ein-depth-favard-recurrence} has simple, purely absolutely
continuous spectrum \([e^{-a},1]\) of multiplicity one, and
\begin{equation}\label{eq:ein-depth-nevai-limits}
 \alpha_n^{(q)}\longrightarrow\frac{1+e^{-a}}2,
 \qquad
 \beta_n^{(q)}\longrightarrow\frac{(1-e^{-a})^2}{16}.
\end{equation}
Consequently
\begin{equation}\label{eq:ein-depth-hankel-capacity}
 \lim_{n\to\infty}(\mathcal H_n^{(q)})^{1/n^2}
 =\frac{1-e^{-a}}4,
\end{equation}
and the normalized zero-counting measures of the monic orthogonal
polynomials \(P_N^{(q)}\) converge weakly to
\begin{equation}\label{eq:ein-depth-arcsine-law}
 \frac{\mathbf 1_{(e^{-a},1)}(x)}
 {\pi\sqrt{(x-e^{-a})(1-x)}}\,dx.
\end{equation}

Finally, let \(J_N^{(q)}\) be the \(N\)-by-\(N\) Jacobi truncation attached
to \eqref{eq:ein-depth-favard-recurrence}.  For every fixed \(q\) and
\(N\geq2\), the \(2N-1\) eigenvalues of
\(J_N^{(q)}\) and \(J_{N-1}^{(q)}\), taken together, are algebraically
independent over \(\Eexp\).  The same joint assertion holds for finitely
many distinct orders \(q\), with one chosen consecutive pair at each
order.  Each spectrum is simple and the two consecutive spectra strictly
interlace.  The count \(2N-1\) is sharp: if \(\Sigma_j^{(q)}\) denotes
the full spectrum of \(J_j^{(q)}\), then
\begin{equation}\label{eq:ein-depth-full-finite-spectral-rank}
 \operatorname{trdeg}_{\Eexp}
 \Eexp\bigl(\Sigma_1^{(q)},\ldots,\Sigma_N^{(q)}\bigr)=2N-1.
\end{equation}
Thus the consecutive pair
\(\Sigma_{N-1}^{(q)}\cup\Sigma_N^{(q)}\) is a transcendence basis for
the full finite spectral tower up to algebraic extension.
\end{corollary}

\begin{proof}
The Mellin identity
\[
 \frac1{n^q}=\frac1{(q-1)!}
 \int_0^1t^{n-1}(-\log t)^{q-1}\,dt
\]
and termwise integration give, first near the origin and then everywhere
by analyticity,
\begin{equation}\label{eq:ein-depth-mellin-integral}
 \EinDepth{q}(s)=\frac1{(q-1)!}
 \int_0^1(1-e^{-st})(-\log t)^{q-1}\frac{dt}{t}.
\end{equation}
Since
\[
 (-1)^{\ell-1}\Delta_a^\ell(1-e^{-st})
 =e^{-st}(1-e^{-at})^\ell,
\]
substitution of \(s=b+am\), followed by \(x=e^{-at}\), gives exactly
\eqref{eq:ein-depth-hausdorff-density}.  Its density is positive on
\((e^{-a},1)\).  At \(x=e^{-a}\) the logarithmic factor vanishes when
\(q>1\) (and is interpreted as one when \(q=1\)); as \(x\uparrow1\),
\[
 w_q(x)=O\!\left((1-x)^{\ell-1}
             \left|\log(1-x)\right|^{q-1}\right).
\]
Thus the density is integrable at both endpoints and has infinite support.
This proves the strict Hausdorff assertion.

For fixed \(q\), every finite set of the \(\mu_m^{(q)}\) is obtained from
the values
\(\EinDepth{q}(b+ak)\) by a matrix whose row indexed by \(m\) has last
nonzero entry, with nonzero coefficient, in column \(m+\ell\).  Distinct
rows therefore have distinct pivots.  For several orders the matrix is
block diagonal in \(q\).  Corollary~\ref{cor:ein-depth-pure-exponential}
then proves \eqref{eq:ein-depth-all-moments-independent}.

Strict positivity of the measure gives \(\mathcal H_n^{(q)}>0\).  Moreover,
\begin{equation}\label{eq:ein-depth-hankel-triangular-jacobian}
 \frac{\partial\mathcal H_n^{(q)}}
      {\partial\mu_{2n}^{(q)}}=\mathcal H_{n-1}^{(q)}.
\end{equation}
The Jacobian with respect to
\(\mu_0^{(q)},\mu_2^{(q)},\ldots,\mu_{2n}^{(q)}\) is triangular with
nonzero diagonal; together with the independence of the moments, this
proves the assertion about all Hankel determinants.

The finite Stieltjes--Favard algorithm gives, for every \(N\geq1\), the
birational equality
\begin{multline}\label{eq:ein-depth-moment-favard-fields}
 \Eexp(\mu_0^{(q)},\ldots,\mu_{2N-1}^{(q)})\\
 =\Eexp(\mu_0^{(q)},
   \alpha_0^{(q)},\ldots,\alpha_{N-1}^{(q)},
   \beta_1^{(q)},\ldots,\beta_{N-1}^{(q)}).
\end{multline}
Both sides have \(2N\) displayed generators.  Hence
\eqref{eq:ein-depth-favard-independent} follows, jointly in all orders.

Strict total positivity follows from Andreief's identity: every minor is
an integral of the product of two generalized Vandermonde determinants
against the positive density in
\eqref{eq:ein-depth-hausdorff-density}.  After ordering the integration
variables, both determinants have the same strict sign.  Normalizing the
measure by its mass identifies the infinite Jacobi operator with
multiplication by \(x\) on its \(L^2\)-space.  Since \(w_q\) is positive
almost everywhere on the whole interval, this gives simple purely
absolutely continuous spectrum \([e^{-a},1]\).  The
Denisov--Rakhmanov theorem~\cite{Denisov2004}, affinely rescaled to that
interval, gives \eqref{eq:ein-depth-nevai-limits}.  The identity
\[
 \mathcal H_n^{(q)}
 =(\mu_0^{(q)})^{n+1}
   \prod_{k=1}^n(\beta_k^{(q)})^{n+1-k}
\]
then proves \eqref{eq:ein-depth-hankel-capacity}; the standard fixed-walk
trace argument for Jacobi truncations with convergent coefficients gives
the arcsine law \eqref{eq:ein-depth-arcsine-law}.

The characteristic polynomial of \(J_N^{(q)}\) is \(P_N^{(q)}\).  Given
the two monic polynomials \(P_N^{(q)}\) and \(P_{N-1}^{(q)}\), successive
Euclidean divisions in \eqref{eq:ein-depth-favard-recurrence} recover
\(\alpha_0^{(q)},\ldots,\alpha_{N-1}^{(q)}\) and
\(\beta_1^{(q)},\ldots,\beta_{N-1}^{(q)}\) rationally; conversely the
recurrence constructs the two polynomials.  Their \(2N-1\) nonleading
coefficients are therefore algebraically independent.  Adjoining all
roots is an algebraic extension and introduces exactly \(2N-1\) displayed
roots, so those roots are algebraically independent as claimed.  Positivity
of the \(\beta_n^{(q)}\) gives simplicity and strict interlacing.
The same Euclidean divisions recover every earlier monic polynomial
\(P_j^{(q)}\), \(j<N-1\), from the coefficients of
\(P_N^{(q)}\) and \(P_{N-1}^{(q)}\).  All earlier eigenvalues are
therefore algebraic over the field generated by the consecutive pair,
which proves \eqref{eq:ein-depth-full-finite-spectral-rank}.
\end{proof}

\begin{corollary}[The order-two Basel split]
\label{cor:ein-depth-basel}
Put
\[
 \eta=\Ein(1)=\gamma+\frac{\delta}{e},\qquad
 T=\int_1^\infty\frac{e^{-u}\log u}{u}\,du.
\]
Then each of the following triples is algebraically independent over
\(\Qbar\):
\begin{equation}\label{eq:ein-depth-basel-independent-triples}
 \bigl(e,\eta,\EinDepth{2}(1)\bigr),
 \qquad
 \left(e,\gamma+\frac{\delta}{e},
       \gamma^2+\frac{\pi^2}{6}-2T\right).
\end{equation}
Consequently
\begin{equation}\label{eq:ein-depth-gamma-delta-lower-bound}
 \operatorname{trdeg}_{\Qbar}
 \Qbar\bigl(e,\gamma,\delta,\EinDepth{2}(1)\bigr)\geq3.
\end{equation}
\end{corollary}

\begin{proof}
The first triple is the case \(s=1\), \(\lambda_1=1\), and \(R=2\) of
Theorem~\ref{thm:ein-depth-multipoint}.  The exact identity
\[
 2\EinDepth{2}(1)+2T=\gamma^2+\frac{\pi^2}{6}
\]
turns the first triple into the second by replacing the third coordinate
with twice itself.  Finally
\(\eta=\gamma+\delta/e\), so the first triple is contained in the field
on the left of \eqref{eq:ein-depth-gamma-delta-lower-bound}.
No individual irrationality or transcendence conclusion for \(\gamma\)
or \(\delta\) is asserted.
\end{proof}

\begin{remark}[Classical lineage and scope]
The functions in \eqref{eq:ein-depth-definition} belong to Hardy's
classical polyexponential lineage.  In the notation
\[
 \mathfrak e_p(x,a)=\sum_{n\geq0}\frac{x^n}{n!(n+a)^p},
\]
one has
\(\EinDepth{r}(z)=z\mathfrak e_{r+1}(-z,1)\), whereas
\(\mathfrak e_1(x,1)=(e^x-1)/x\).  Thus \(a=1\) is a resonant integral
shift whose first coordinate is already elementary.  Shidlovskii proved an all-order
algebraic-independence theorem for a nearby shifted family with a
nonintegral rational shift; see \cite[Sec.~7.3.3]{Shidlovskii1989} and the
account \cite[Sec.~9]{Boyadzhiev2007}.  The argument above treats the
regularized integral-shift endpoint needed for \(\Ein\), simultaneously at
several algebraic dilation points, and identifies the complete toric
relation ideal.
It does not cover multiple-index polyexponentials, whose shuffle-type
relations lie outside the present one-level tower.  We make no broader
historical-priority claim.
\end{remark}


\begin{corollary}[All-order gamma-jet split]
\label{cor:all-order-gamma-jet-split}
Write the Laurent expansion at the origin as
\[
 \Gamma(s)=\frac1s+\sum_{k\geq0}g_ks^k
\]
and put
\[
 T_k=\frac1{k!}\int_1^\infty
       \frac{e^{-u}(\log u)^k}{u}\,du\qquad(k\geq0).
\]
Then, for every $k\geq0$,
\begin{equation}\label{eq:all-order-gamma-jet-split}
 g_k=(-1)^{k+1}\EinDepth{k+1}(1)+T_k.
\end{equation}
Consequently the countable family
\[
 \{g_k-T_k:k\geq0\}
\]
is algebraically independent over $\Eexp$.  In particular, for every
$R\geq1$ the $R+1$ numbers
\[
 e,\quad g_0-T_0,\quad g_1-T_1,\ \ldots,\quad g_{R-1}-T_{R-1}
\]
are algebraically independent over $\overline{\mathbb Q}$.

The coefficients $g_k$ are explicit polynomials in the classical gamma
jet.  Namely, if $Y_n$ denotes the complete exponential Bell polynomial,
then
\begin{equation}\label{eq:gamma-jet-bell-polynomial}
 g_k=\frac1{(k+1)!}
 Y_{k+1}\bigl(x_1,\ldots,x_{k+1}\bigr),
 \qquad
 x_1=-\gamma,\quad
 x_n=(-1)^n(n-1)!\zeta(n)\ (n\geq2).
\end{equation}
Thus \eqref{eq:all-order-gamma-jet-split} gives, at every order, an
explicit algebraically independent combination of
$\gamma,\zeta(2),\ldots,\zeta(k+1)$ and the tail $T_k$.
For $k=1$ it is precisely
\[
 \EinDepth{2}(1)+\int_1^\infty\frac{e^{-u}\log u}{u}\,du
 =\frac{\gamma^2}{2}+\frac{\pi^2}{12}.
\]
No individual arithmetic conclusion for a zeta value or a tail integral
is asserted.
\end{corollary}

\begin{proof}
For $\Re s>0$, split Euler's integral at $1$:
\[
 \Gamma(s)=\int_0^1e^{-u}u^{s-1}\,du
           +\int_1^\infty e^{-u}u^{s-1}\,du.
\]
After isolating the pole $1/s$, the resulting identities extend
meromorphically to a neighborhood of $s=0$, where termwise integration of
the first part gives
\begin{align*}
 \int_0^1e^{-u}u^{s-1}\,du
 &=\frac1s+\sum_{n\geq1}\frac{(-1)^n}{n!(n+s)}\\
 &=\frac1s+\sum_{k\geq0}(-1)^{k+1}
   \EinDepth{k+1}(1)s^k.
\end{align*}
Expansion under the exponentially convergent second integral gives
\[
 \int_1^\infty e^{-u}u^{s-1}\,du
 =\sum_{k\geq0}T_ks^k.
\]
Coefficient comparison proves \eqref{eq:all-order-gamma-jet-split}.
The algebraic-independence assertions are
Theorem~\ref{thm:ein-depth-multipoint} and
Corollary~\ref{cor:ein-depth-pure-exponential} at the point $1$.

Finally,
\[
 \Gamma(1+s)=
 \exp\!\left(-\gamma s+
   \sum_{n\geq2}\frac{(-1)^n\zeta(n)}n s^n\right).
\]
The defining generating identity for the complete exponential Bell
polynomials gives \eqref{eq:gamma-jet-bell-polynomial}.
\end{proof}


\section{Gamma-free curvature, logarithms, and a probabilistic transfer}
\label{sec:crossfield-transfers}

The pure-exponential base change is stable under every full-rank linear
projection.  This produces explicit gamma-free families among classical
logarithmic special functions and transfers the same algebraic matroid to a
probability law.

\subsection{A geometric-grid curvature family}

Fix an integer \(q\ge2\).  For \(n\ge0\), define
\begin{align*}
 A_n&=E_1(q^{n+2})-2E_1(q^{n+1})+E_1(q^n),\\
 B_n&=\Ei(q^{n+2})-2\Ei(q^{n+1})+\Ei(q^n),\\
 C_n&=\Ci(q^{n+2})-2\Ci(q^{n+1})+\Ci(q^n),\\
 S_n&=\Si(q^{n+2})-2\Si(q^{n+1})+\Si(q^n).
\end{align*}

\begin{theorem}[Gamma-free geometric curvature]
\label{thm:geometric-curvature}
The countable family
\begin{equation}\label{eq:geometric-curvature-family}
 \{A_n,B_n,C_n,S_n:n\ge0\}
\end{equation}
is algebraically independent over \(\Eexp\), in the finite-subset sense.
\end{theorem}

\begin{proof}
For positive \(x\), the connection formulas are
\[
 \Ein(x)=\gamma+\log x+E_1(x),\qquad
 \Ein(-x)=\gamma+\log x-\Ei(x),
\]
and
\[
 \frac{\Ein(ix)+\Ein(-ix)}2=\gamma+\log x-\Ci(x),
 \qquad
 \frac{\Ein(ix)-\Ein(-ix)}{2i}=\Si(x).
\]
The second-difference row \((1,-2,1)\) kills both the constant \(\gamma\)
and the affine sequence \(\log(q^n)=n\log q\).  Hence every displayed
quantity is a linear form in \(\Ein\)-values on one of the four disjoint
orbits \(\{q^j\}\), \(\{-q^j\}\), \(\{iq^j\}\), and \(\{-iq^j\}\).
Within each block, every finite family of second-difference rows is linearly
independent: inspect the largest terminal index in a linear combination.
The four blocks have disjoint support.  Lemma~\ref{lem:linear-image-ein}
therefore proves joint algebraic independence over \(\Eexp\).
\end{proof}

The geometric grid is essential to the uncorrected formula: on a general
grid one must add the logarithmic second difference
\(\log(x_{n+2}x_n/x_{n+1}^2)\).  Omitting it would reintroduce the connection
term and invalidate the linear-image argument.

\subsection{Logarithms of rational period expressions}

\begin{proposition}[Global logarithm principle]
\label{prop:global-log-principle}
Let \(\lambda_1,\ldots,\lambda_s\in\Qbar^\times\) be distinct, and let
\(R\in\Eexp(X_1,\ldots,X_s)\) be nonconstant.  Suppose
\[
 r=R(\Ein(\lambda_1),\ldots,\Ein(\lambda_s))
\]
is defined and nonzero.  Every solution of \(e^w=r\), hence every branch of
\(\log r\), is transcendental.
\end{proposition}

\begin{proof}
If \(w\) were algebraic, then \(e^w\in\Eexp\).  On the other hand,
Theorem~\ref{thm:pure-exp-base} makes the \(\Ein(\lambda_j)\) independent
indeterminates over \(\Eexp\), so a nonconstant rational expression in
them cannot belong to \(\Eexp\).  This contradiction proves the claim.
\end{proof}

The proposition applies after every invertible Fourier transform in
\eqref{eq:fourth-root-transform}.  For example, every logarithm of each of
\(\Si(a),\Cin(a),\Shi(a),\Cinh(a)\), with
\(a\in\Qbar^\times\), is transcendental.  It also proves the transcendence
of every branch of \(\log\Ein(1)\) without an exceptional-value analysis;
the non-ultra-Liouville refinement in \cref{sec:log} still uses the
quantitative theorem of Fischler--Rivoal.

More generally, let
\(G(W,e^{\beta_1W},\ldots,e^{\beta_mW})\) be a rational expression over
\(\Qbar\), where the \(\beta_j\) are algebraic.  Every solution \(w\) of
\[
 G(w,e^{\beta_1w},\ldots,e^{\beta_mw})=r
\]
is transcendental whenever both sides are defined: an algebraic \(w\)
would make the left side an element of \(\Eexp\).  Logarithms and Lambert
\(W\)-preimages of \(r\) are elementary special cases.

\subsection{The Dickman log-Laplace matroid}

Let \(\Phi\) denote the Laplace transform of the normalized Dickman law.
The classical de Bruijn formula is
\begin{equation}\label{eq:dickman-laplace}
 \Phi(s)=\exp\left(\int_0^1\frac{e^{-sx}-1}{x}\,dx\right)
        =e^{-\Ein(s)}\qquad(s>0);
\end{equation}
see, for example, \cite{Moree2014}.  Thus the L\'evy exponent, rather than
the transform value itself, is an \(\Ein\)-period.

\begin{corollary}[Dickman log-Laplace periods]
\label{cor:dickman-matroid}
For distinct positive algebraic \(s_1,\ldots,s_N\), the numbers
\[
 -\log\Phi(s_j)=\Ein(s_j)\qquad(1\le j\le N)
\]
are algebraically independent over \(\Eexp\).  If these parameters label
the vertices of an oriented graph, then the log-ratios
\[
 \log\frac{\Phi(s_{t(e)})}{\Phi(s_{s(e)})}
 =-\Ein(s_{t(e)})+\Ein(s_{s(e)})
\]
have precisely the signed cycle relations and no others.
\end{corollary}

\begin{proof}
Use \eqref{eq:dickman-laplace}, Theorem~\ref{thm:pure-exp-base}, and the
incidence-matrix argument of Theorem~\ref{thm:graphic-period}.
\end{proof}

This is an exact transfer to the logarithmic Laplace coordinates of a
classical smooth-number distribution.  It gives no arithmetic conclusion
about the values \(\Phi(s_j)=e^{-\Ein(s_j)}\) themselves, which lie outside
the finite-point Siegel--Shidlovskii mechanism.


\section{New classical families from balanced \texorpdfstring{\(\Ein\)}{Ein} coordinates}
\label{sec:new-classical-families}

The pure-exponential base-change theorem turns every full-rank linear image
of distinct \(\Ein\)-coordinates into an exact algebraic-independence
statement.  Two consequences are particularly concrete: a pairwise
algebraicity exclusion at every positive algebraic level, and a countable
family that is simultaneously algebraically independent and sharply small.

\subsection{Four classical values at every algebraic level}

\begin{theorem}[Four classical values at an algebraic level]
\label{thm:four-classical-algebraic-level}
For every positive algebraic number \(a\), among
\begin{equation}\label{eq:four-classical-level-values}
 E_1(a),\qquad \Ei(a),\qquad \Ci(a),\qquad \Chi(a)
\end{equation}
at most one is algebraic.

More precisely, each of the following six quantities is transcendental over
\(\Eexp\):
\begin{align}
 E_1(a)+\Ei(a),& & E_1(a)+\Ci(a),& & E_1(a)+\Chi(a),\notag\\
 \Ei(a)-\Ci(a),& & \Ei(a)-\Chi(a),& & \Chi(a)-\Ci(a).
 \label{eq:four-classical-witnesses}
\end{align}
\end{theorem}

\begin{proof}
Put
\[
 A=\Ein(a),\qquad B=\Ein(-a),\qquad
 C=\frac{\Ein(ia)+\Ein(-ia)}2.
\]
The connection formulas give, in the order displayed in
\eqref{eq:four-classical-witnesses},
\begin{equation}\label{eq:four-classical-witness-identities}
 A-B,\qquad A-C,\qquad \frac{A-B}{2},\qquad
 C-B,\qquad \frac{A-B}{2},\qquad
 C-\frac{A+B}{2}.
\end{equation}
The four values
\(\Ein(a),\Ein(-a),\Ein(ia),\Ein(-ia)\) are algebraically independent
over \(\Eexp\).  Every expression in
\eqref{eq:four-classical-witness-identities} is a nonzero linear form in
those independent coordinates and is therefore transcendental over
\(\Eexp\).  If any two members of
\eqref{eq:four-classical-level-values} were algebraic, the corresponding
witness in \eqref{eq:four-classical-witnesses} would be algebraic, a
contradiction.
\end{proof}

\subsection{A sharply decaying algebraically independent family}

\begin{theorem}[Quadratic balanced \(E_1\) family]
\label{thm:quadratic-balanced-E1}
For \(m\geq1\), put
\begin{equation}\label{eq:quadratic-balanced-E1-definition}
 Q_m=E_1(m^2)-2E_1\bigl(m(m+1)\bigr)+E_1\bigl((m+1)^2\bigr).
\end{equation}
Then the countable family \(\{Q_m:m\geq1\}\) is algebraically independent
over \(\Eexp\), in the finite-subset sense.  Every \(Q_m\) is positive, and
\begin{equation}\label{eq:quadratic-balanced-E1-asymptotic}
 Q_m=\frac{e^{-m^2}}{m^2}
 \left(1+O(m^{-2})+O(e^{-m})\right).
\end{equation}
Consequently
\begin{equation}\label{eq:quadratic-balanced-E1-saturation}
 \boxed{
 \lim_{m\to\infty}
 \frac{-\log Q_m}{(m+1)^2}=1.}
\end{equation}
\end{theorem}

\begin{proof}
The two cancellations
\[
 1-2+1=0,
 \qquad
 \log(m^2)-2\log\bigl(m(m+1)\bigr)+\log\bigl((m+1)^2\bigr)=0
\]
give the exact identity
\begin{equation}\label{eq:quadratic-balanced-Ein-form}
 Q_m=\Ein(m^2)-2\Ein\bigl(m(m+1)\bigr)
       +\Ein\bigl((m+1)^2\bigr).
\end{equation}
For \(m\geq1\), the integer \(m(m+1)\) is not a square: since
\(\gcd(m,m+1)=1\), a square product would force both factors to be
squares, and two consecutive positive squares do not exist.  Moreover
\(m(m+1)\) is strictly increasing with \(m\).  Hence the middle coordinate
\(\Ein(m(m+1))\) occurs in no row of
\eqref{eq:quadratic-balanced-Ein-form} other than the \(m\)-th row.  In any
finite linear combination of these rows, inspection of the middle
coordinates forces every row coefficient to vanish.  The rows therefore
have full rank.  The pure-exponential base-change theorem and
Lemma~\ref{lem:linear-image-ein} prove the asserted algebraic independence.

For \(x>0\), integration by parts gives the standard bounds
\begin{equation}\label{eq:E1-two-sided-bound}
 \frac{e^{-x}}{x+1}<E_1(x)<\frac{e^{-x}}x.
\end{equation}
Indeed, the upper bound is immediate, while
\[
 E_1(x)=\frac{e^{-x}}x-\int_x^\infty\frac{e^{-t}}{t^2}\,dt
 >\frac{e^{-x}}x-\frac{E_1(x)}x
\]
gives the lower bound.  Therefore
\[
 \frac{2E_1(m^2+m)}{E_1(m^2)}
 <2e^{-m}\frac{m^2+1}{m^2+m}
 \leq\frac2e<1.
\]
Thus \(E_1(m^2)>2E_1(m^2+m)\), and the last positive term in
\eqref{eq:quadratic-balanced-E1-definition} proves \(Q_m>0\).

Finally, the standard expansion
\[
 E_1(x)=\frac{e^{-x}}x\left(1+O(x^{-1})\right)
\]
shows that
\[
 \frac{E_1(m^2+m)}{E_1(m^2)}=O(e^{-m}),
 \qquad
 \frac{E_1((m+1)^2)}{E_1(m^2)}=O(e^{-2m}).
\]
This proves \eqref{eq:quadratic-balanced-E1-asymptotic}; taking logarithms
gives \eqref{eq:quadratic-balanced-E1-saturation}.
\end{proof}

\begin{remark}[Uniform moving-point boundary]
The scale \((m+1)^2\) in
\eqref{eq:quadratic-balanced-E1-saturation} is the largest evaluation
argument, not a Weil height or a relation-polynomial height.  The theorem
therefore gives a sharp counterexample to naive lower bounds that are
uniform only in the moving archimedean evaluation scale.  It does not
contradict fixed-point algebraic-independence measures for \(E\)-functions.
\end{remark}

\subsection{A conditional \texorpdfstring{\(E\)--\(D\)}{E--D} collision}

\begin{corollary}[An explicit conditional \(E\)--\(D\) collision]
\label{cor:gamma-algebraic-ED-collision}
Let \(\mathbf E\) and \(\mathbf D\) be the rings of finite \(E\)-values
and arithmetic-Gevrey values of order one, respectively, in the sense of
Fischler--Rivoal \cite{FR2024ED}.  Then
\begin{equation}\label{eq:gamma-algebraic-ED-collision}
 \boxed{
 \gamma\in\Qbar
 \quad\Longrightarrow\quad
 \delta\in(\mathbf E\cap\mathbf D)\setminus\Qbar.}
\end{equation}
Consequently the conjectural equality
\(\mathbf E\cap\mathbf D=\Qbar\) implies that \(\gamma\) is
transcendental.
\end{corollary}

\begin{proof}
Assume \(\gamma\in\Qbar\).  Then
\[
 f(z)=e^z\bigl(\Ein(z)-\gamma\bigr)
\]
is an \(E\)-function with algebraic Taylor coefficients, and
\(f(1)=\delta\).  Hence \(\delta\in\mathbf E\).  The standard Borel-sum
realization of Euler's factorial series gives \(\delta\in\mathbf D\).
If \(\delta\) were algebraic, then
\[
 \Ein(1)=\gamma+e^{-1}\delta
\]
would belong to \(\Eexp\), contradicting the pure-exponential base-change
theorem.  Thus \(\delta\notin\Qbar\), proving the assertion.
\end{proof}


\section{Common-shift defect one and residual families}
\label{sec:common-shift-residuals}

The following field-theoretic lemma isolates exactly what a common
translation can and cannot lose.  The independent anchor is essential: merely
excluding the possibility that the shift parameter belongs to the base field
would not prove the final two-variable independence statement.

\begin{theorem}[Common-shift defect one]
\label{thm:common-shift-defect}
Let \(K\subset\C\) be a field of characteristic zero, and let
\(z_0,z_1,\ldots,z_r\) be algebraically independent over \(K\), where
\(r\geq1\).  Fix \(c,e_0\in K^\times\), and put
\begin{equation}\label{eq:common-shift-coordinates}
 T_j=c(z_j-\gamma)\quad(1\leq j\leq r),
 \qquad
 \Delta=e_0(z_0-\gamma).
\end{equation}
Then
\begin{equation}\label{eq:common-shift-lower-bound}
 \trdeg_K K(T_1,\ldots,T_r)\geq r-1.
\end{equation}
More precisely, the \(r-1\) anchored differences
\[
 T_2-T_1,\ldots,T_r-T_1
\]
are algebraically independent over \(K\).  Consequently the \(T_j\) are
algebraically dependent over \(K\) if and only if
\[
 \trdeg_K K(T_1,\ldots,T_r)=r-1.
\]
In that case \(\gamma\) and \(\Delta\) are algebraically independent over
\(K\).
\end{theorem}

\begin{proof}
The linear forms \(z_j-z_1\), \(2\leq j\leq r\), have row rank \(r-1\)
in the independent coordinates \(z_1,\ldots,z_r\).  Hence
\[
 T_j-T_1=c(z_j-z_1),\qquad 2\leq j\leq r,
\]
are algebraically independent over \(K\), proving
\eqref{eq:common-shift-lower-bound}.

Suppose now that the \(T_j\) are algebraically dependent and write
\(L=K(T_1,\ldots,T_r)\).  The lower bound gives
\(\trdeg_K L=r-1\).  Since \(z_j=\gamma+c^{-1}T_j\), one has
\begin{equation}\label{eq:common-shift-field-identity}
 K(z_1,\ldots,z_r,\gamma)=L(\gamma).
\end{equation}
The left side contains the purely transcendental field
\(K(z_1,\ldots,z_r)\), of transcendence degree \(r\), whereas the right
side has transcendence degree at most \((r-1)+1=r\).  Thus both sides of
\eqref{eq:common-shift-field-identity} have transcendence degree exactly
\(r\).  It follows that \(\gamma\) is algebraic over
\(K(z_1,\ldots,z_r)\), and also that \(\gamma\) is transcendental over
\(L\), hence over \(K\).

Because \(z_0,z_1,\ldots,z_r\) are algebraically independent, \(z_0\) is
transcendental over the algebraic closure of
\(K(z_1,\ldots,z_r)\).  In particular it is transcendental over
\(K(\gamma)\).  Therefore \(z_0-\gamma\) is transcendental over
\(K(\gamma)\).  Since \(\gamma\) is transcendental over \(K\) and
\(e_0\in K^\times\), the two numbers
\(\gamma\) and \(\Delta=e_0(z_0-\gamma)\) are algebraically independent
over \(K\).
\end{proof}

We apply the theorem over \(K=\Eexp\).  For \(N,n\geq1\), define
\begin{align}
 Y_N&=(N+1)\Ein(2^N)-N\Ein(2^{N+1}),
 &\xi_n&=\Ein(-n),\label{eq:residual-Ein-coordinates}\\
 R_N&=e\bigl((N+1)E_1(2^N)-N E_1(2^{N+1})\bigr),
 &S_n&=e\bigl(\Ei(n)-\log n\bigr).
 \label{eq:residual-classical-coordinates}
\end{align}
The connection formulas give
\begin{equation}\label{eq:residual-common-shift}
 R_N=e(Y_N-\gamma),
 \qquad
 -S_n=e(\xi_n-\gamma),
 \qquad
 \delta=e(\Ein(1)-\gamma).
\end{equation}

\begin{lemma}[Independent residual coordinates]
\label{lem:residual-Ein-independent}
The countable family
\[
 \{\Ein(1)\}\cup\{Y_N:N\geq1\}\cup\{\xi_n:n\geq1\}
\]
is algebraically independent over \(\Eexp\), in the finite-subset sense.
\end{lemma}

\begin{proof}
The pure-exponential base-change theorem makes the distinct coordinates
\[
 \Ein(1),\qquad \Ein(2^k)\ (k\geq1),\qquad \Ein(-n)\ (n\geq1)
\]
algebraically independent over \(\Eexp\).  The negative coordinates are
disjoint from the positive ones.  It remains only to check the row rank of
the \(Y_N\).  In a nonzero finite linear combination of their defining
rows, choose the smallest occurring index \(N\).  The coordinate
\(\Ein(2^N)\) occurs only in that smallest row and has coefficient \(N+1\),
so its row coefficient must vanish.  Iteration proves that every row
coefficient vanishes.  Lemma~\ref{lem:linear-image-ein} now gives the
assertion.
\end{proof}

\begin{corollary}[Countable residual dichotomy]
\label{cor:countable-residual-dichotomy}
Exactly one of the following alternatives holds.
\begin{enumerate}[label=\textup{(\roman*)}]
 \item The family
 \[
  \{R_N:N\geq1\}\cup\{S_n:n\geq1\}
 \]
 is algebraically independent over \(\Eexp\).
 \item Some finite subfamily is algebraically dependent over \(\Eexp\),
 and \(\gamma,\delta\) are then algebraically independent over
 \(\Eexp\).  In particular \(e,\gamma,\delta\) are algebraically
 independent over \(\Qbar\).
\end{enumerate}
Unconditionally, the anchored family
\begin{equation}\label{eq:anchored-residual-family}
 \boxed{
 \{R_N-R_1:N\geq2\}\cup\{S_n+R_1:n\geq1\}}
\end{equation}
is algebraically independent over \(\Eexp\).
\end{corollary}

\begin{proof}
Apply Theorem~\ref{thm:common-shift-defect} to every finite selection from
\eqref{eq:residual-common-shift}, with
\(z_0=\Ein(1)\), common scale \(c=e\), and \(\Delta=\delta\).
Lemma~\ref{lem:residual-Ein-independent} verifies the hypothesis.  A
countable family is algebraically dependent exactly when one of its finite
subfamilies is dependent, which proves the dichotomy.

For the final assertion, anchor the common-shift variables at \(R_1\).  Their
differences are
\[
 R_N-R_1=e(Y_N-Y_1),
 \qquad
 (-S_n)-R_1=-(S_n+R_1),
\]
and the anchored-difference assertion of
Theorem~\ref{thm:common-shift-defect} proves their joint algebraic
independence.
\end{proof}

\begin{remark}[Logical scope]
The corollary is an exact alternative, not a construction of the relation in
its second branch.  If no finite residual relation is supplied, the result is
the unconditional algebraic independence of the anchored family
\eqref{eq:anchored-residual-family}; it does not decide \(\gamma\) or
\(\delta\) individually.
\end{remark}


\section{A six-constant consequence}

Put
\[
 A=\Ein(1),\qquad B=\Ein(-1),\qquad
 C=\Cin(1)=\frac{\Ein(i)+\Ein(-i)}2.
\]
Applying \cref{thm:master} to \(\{1,-1,i,-i\}\) shows that
\begin{equation}\label{eq:eabc-ind}
 e,A,B,C\quad\text{are algebraically independent over }\Qb.
\end{equation}
The identities needed below are
\begin{align}
 E_1(1)&=A-\gamma,&
 \Ei(1)&=\gamma-B,\label{eq:six-identities}\\
 \Ci(1)&=\gamma-C,&
 \Chi(1)&=\gamma-\frac{A+B}{2},&
 \delta&=e(A-\gamma).\nonumber
\end{align}

\begin{proposition}[The remaining dimension]\label{prop:one-dimension}
For
\[
 \mathcal S=\{\gamma,\delta,E_1(1),\Ei(1),\Ci(1),\Chi(1)\}
\]
one has
\[
 \Qb(\mathcal S)=\Qb(\gamma,e,A,B,C)
\]
and therefore
\[
 4\le \trdeg_{\Qb}\Qb(\mathcal S)\le5.
\]
In particular,
\[
 \trdeg_{\Qb}\Qb\bigl(\gamma,\delta,E_1(1),\Ei(1)\bigr)\ge3.
\]
\end{proposition}

\begin{proof}
The forward inclusion follows from
\(e=\delta/E_1(1)\), \(A=\gamma+E_1(1)\),
\(B=\gamma-\Ei(1)\), and \(C=\gamma-\Ci(1)\); here
\(E_1(1)>0\).  The reverse inclusion follows from
\eqref{eq:six-identities}.  Algebraic independence of \(e,A,B,C\) gives the
lower bound, and adjoining \(\gamma\) gives the upper bound.  For the final
claim, the four-number field contains
\(e=\delta/E_1(1)\), \(A=\gamma+E_1(1)\), and
\(B=\gamma-\Ei(1)\), which are algebraically independent.
\end{proof}

\begin{theorem}\label{thm:six}
No two members of
\[
 \mathcal S=\{\gamma,\delta,E_1(1),\Ei(1),\Ci(1),\Chi(1)\}
\]
are algebraic.  Hence at least five members of \(\mathcal S\) are
transcendental.
\end{theorem}

\begin{proof}
First omit \(\delta\).  If any two of the remaining five numbers were
algebraic, adding or subtracting their expressions in
\eqref{eq:six-identities} would put one of the following nonzero linear forms
in \(A,B,C\) in \(\Qb\):
\[
 A,\ B,\ C,\ A+B,\ A-B,\ A-C,\ C-B,\ \frac{A+B}{2}-C.
\]
This contradicts \eqref{eq:eabc-ind}.

Now pair \(\delta\) with each of the other constants.  If \(\delta\) and
\(\gamma\) were algebraic, the identity
\(\delta=e(A-\gamma)\) would be a nonzero algebraic relation between \(e\)
and \(A\).  If \(\delta\) and \(E_1(1)\) were algebraic, positivity of
\(E_1(1)\) would give \(e=\delta/E_1(1)\in\Qb\).  Finally, algebraicity of
\(\delta\) together with, respectively, \(\Ei(1),\Ci(1),\Chi(1)\), would give
one of the relations
\[
 \delta=e\bigl(A-B-\Ei(1)\bigr),\qquad
 \delta=e\bigl(A-C-\Ci(1)\bigr),
\]
or
\[
 \delta=e\left(\frac{A-B}{2}-\Chi(1)\right),
\]
contradicting algebraic independence of the indicated subset of
\(e,A,B,C\).  All fifteen pairs are therefore excluded.
\end{proof}

\begin{corollary}\label{cor:conditional-triples}
If \(\gamma\) is algebraic, then \(E_1(1),\Ei(1),\delta\) are algebraically
independent.  More generally, if one of \(E_1(1),\Ei(1),\delta\) is algebraic,
then \(\gamma\) and the other two are algebraically independent.
\end{corollary}

\begin{proof}
Use \(A=\gamma+E_1(1)\), \(B=\gamma-\Ei(1)\), and
\(e=\delta/E_1(1)\) when the relevant constant is assumed algebraic and
nonzero.  A polynomial relation among the claimed triple would then induce a
polynomial relation among an appropriate subset of \(e,A,B\).  The changes of
coordinates are birational, so nonzero relations remain nonzero.
\end{proof}

\begin{theorem}[Conditional relation ideal after one algebraic specialization]
\label{thm:six-one-algebraic}
Write
\[
 (\Gamma,D,X,Y,Z,W)
 =
 (\gamma,\delta,E_1(1),\Ei(1),\Ci(1),\Chi(1))
\]
and let \(I_{\mathcal S}\) be the kernel of the corresponding evaluation
homomorphism from \(\Qbar[\Gamma,D,X,Y,Z,W]\) to \(\C\).  If one displayed
constant \(\kappa\) is algebraic, say \(\kappa=c\in\Qbar\), then
\[
 \trdeg_{\Qbar}\Qbar(\mathcal S)=4
\]
and
\begin{equation}\label{eq:six-one-algebraic-ideal}
 \boxed{
 I_{\mathcal S}
 =\bigl(Y-X-2W,\ T_\kappa-c\bigr),}
\end{equation}
where \(T_\kappa\) denotes the coordinate occupied by \(\kappa\).  In
particular the Euler-level point \(w_*\) of
Theorem~\ref{thm:H-euler-level} is transcendental.
\end{theorem}

\begin{proof}
Recall that \(e,A,B,C\) are algebraically independent and that
\(\Qbar(\mathcal S)=\Qbar(\gamma,e,A,B,C)\).  In the six possible cases,
algebraicity of \(\kappa=c\) expresses \(\gamma\) inside
\(\Qbar(e,A,B,C)\) as follows:
\[
\begin{array}{c|c@{\qquad}c|c@{\qquad}c|c}
 \kappa&\gamma&\kappa&\gamma&\kappa&\gamma\\ \hline
 \gamma&c&\delta&A-c/e&E_1(1)&A-c\\
 \Ei(1)&B+c&\Ci(1)&C+c&\Chi(1)&(A+B)/2+c.
\end{array}
\]
Hence \(\Qbar(\mathcal S)=\Qbar(e,A,B,C)\), of transcendence degree four.
The evaluation kernel therefore has height two.  It contains the two
independent linear polynomials in \eqref{eq:six-one-algebraic-ideal}; these
generate a height-two prime ideal and hence the complete kernel.  If \(w_*\)
were algebraic, Corollary~\ref{cor:witness-six-constant-field} would instead
give transcendence degree five, a contradiction.
\end{proof}

\section{Conditional relation ideals and gamma detectors}
\label{sec:gamma-detectors}

The structural theorems above do not decide the arithmetic nature of
\(\gamma\).  They do, however, turn several concrete relations into exact
certificates against algebraicity.  The point of this section is to state
those certificates with their full quantifiers and relation ideals.

Put
\[
 D=\delta,\qquad X=E_1(1),\qquad Y=\Ei(1),\qquad
 Z=\Ci(1),\qquad W=\Chi(1),
\]
and retain the pure exponential field
\(\Eexp=\Qbar(e^\alpha:\alpha\in\Qbar)\).

\begin{proposition}[The conditional ideal over the pure exponential field]
\label{prop:conditional-pure-exp-ideal}
If \(\gamma\in\Qbar\), then \(X,Y,Z\) are algebraically independent over
\(\Eexp\), and the complete relation ideal of \(D,X,Y,Z,W\) over
\(\Eexp\) is
\begin{equation}\label{eq:conditional-pure-exp-ideal}
 (\,\mathsf D-e\mathsf X,\;2\mathsf W-\mathsf Y+\mathsf X\,).
\end{equation}
Consequently
\begin{equation}\label{eq:gamma-or-delta-global}
 \boxed{\ \gamma\text{ is transcendental, or }
 \delta\text{ is transcendental over }\Eexp.\ }
\end{equation}
\end{proposition}

\begin{proof}
Under the hypothesis, the connection identities give
\[
 X=\Ein(1)-\gamma,\qquad
 Y=\gamma-\Ein(-1),\qquad
 Z=\gamma-\frac{\Ein(i)+\Ein(-i)}2.
\]
They are three linearly independent combinations of the four
\(\Ein\)-coordinates at \(1,-1,i,-i\).  Theorem
\ref{thm:pure-exp-base} and Lemma~\ref{lem:linear-image-ein} therefore make
them algebraically independent over \(\Eexp\).  The remaining coordinates
satisfy, and are determined by,
\[
 D=eX,\qquad 2W=Y-X.
\]
This proves the ideal statement.  In particular \(D=eX\) is transcendental
over \(\Eexp\) whenever \(\gamma\) is algebraic, proving
\eqref{eq:gamma-or-delta-global}.
\end{proof}

For \(m\ge1\), set
\[
 T_m=e^{\gamma^m}.
\]

\begin{theorem}[Exact \(T_m\)-ideals under algebraic \(\gamma\)]
\label{thm:Tm-exact-ideals}
Assume \(\gamma\in\Qbar\), fix \(m\ge1\), and put
\(\eta=\gamma^m\).  If \(\eta\notin\Q\), then
\[
 T_m,D,X,Y,Z
\]
are algebraically independent over \(\Qbar\).  If
\(\eta=a/b\in\Q_{>0}\), where \((a,b)=1\), then the complete relation
ideal of these five coordinates is the principal prime ideal
\begin{equation}\label{eq:Tm-rational-ideal}
 (\,\mathsf D^a-\mathsf T^b\mathsf X^a\,).
\end{equation}
After adjoining \(W\), in the first case one adds only
\(2\mathsf W-\mathsf Y+\mathsf X\); in the second case the full ideal is
generated by that linear form and \eqref{eq:Tm-rational-ideal}.

In either case \(T_m\) and \(D\) are algebraically independent.  Hence,
unconditionally,
\begin{equation}\label{eq:Tm-delta-detector}
 \boxed{\ P(e^{\gamma^m},\delta)=0\text{ for some }
 0\ne P\in\Qbar[U,V]\quad\Longrightarrow\quad
 \gamma\text{ is transcendental}.\ }
\end{equation}
\end{theorem}

\begin{proof}
Since \(0<\gamma<1\), the number \(\eta\) is nonzero and differs from
\(1,-1,i,-i\).  Apply Theorem~\ref{thm:master} to
\(\{1,-1,i,-i,\eta\}\).  With
\[
 A=\Ein(1),\qquad B=\Ein(-1),\qquad
 C=\frac{\Ein(i)+\Ein(-i)}2,
\]
the relevant coordinates are birationally related by
\[
 e=\frac D X,\qquad A=X+\gamma,\qquad
 B=\gamma-Y,\qquad C=\gamma-Z.
\]
If \(1\) and \(\eta\) are \(\Q\)-independent, the exponential coordinates
\(e,T_m\) have no toric relation, and the five displayed coordinates are
independent.  If \(\eta=a/b\), their only toric relation is
\(T_m^b=e^a\), which becomes
\(D^a=T_m^bX^a\).  The exponent vector is primitive, so the binomial is
prime; the transcendence-degree formula gives height one and excludes any
additional relation.  The identity \(2W=Y-X\) handles \(W\).

In both cases contraction to \(\Qbar[\mathsf T,\mathsf D]\) is zero.  Thus
a relation as in \eqref{eq:Tm-delta-detector} contradicts the temporary
hypothesis \(\gamma\in\Qbar\).
\end{proof}

An exact zero is not the only possible certificate.  A quantitative
\(E\)-function theorem also rules out fixed-degree super-Liouville
approximation under the algebraicity hypothesis.

\begin{corollary}[Fixed-degree approximation detector]
\label{cor:fixed-degree-gamma-detector}
Fix \(m,d\ge1\).  Suppose that nonzero polynomials
\(P_n\in\Z[U,V]\), of degree at most \(d\), satisfy
\[
 H(P_n)\longrightarrow\infty,\qquad
 P_n(T_m,D)\ne0,
\]
and
\begin{equation}\label{eq:fixed-degree-super-small}
 \frac{-\log|P_n(T_m,D)|}{\log H(P_n)}\longrightarrow\infty.
\end{equation}
Then \(\gamma\) is transcendental.
\end{corollary}

\begin{proof}
Assume \(\gamma\in\Qbar\).  Then
\[
 f_1(z)=e^{\gamma^m z},\qquad
 f_2(z)=e^z\bigl(\Ein(z)-\gamma\bigr)
\]
are \(E\)-functions with algebraic coefficients and
\((f_1(1),f_2(1))=(T_m,D)\).  Theorem~\ref{thm:Tm-exact-ideals} makes the
two values algebraically independent.  The fixed-degree case of the
Adamczewski--Faverjon algebraic-independence measure
\cite{AdamczewskiFaverjon2025} supplies constants \(c_d,\kappa_d>0\) with
\[
 |P(T_m,D)|\ge c_dH(P)^{-\kappa_d}
\]
for every nonzero integral polynomial of degree at most \(d\).  This
contradicts \eqref{eq:fixed-degree-super-small}.
\end{proof}

\begin{remark}[Logical boundary]
Theorems~\ref{thm:Tm-exact-ideals} and
Corollary~\ref{cor:fixed-degree-gamma-detector} are one-way detectors.
They do not construct a polynomial relation or a super-small sequence.
Moreover \(e^{\gamma^m}\) is an \(E\)-value with algebraic coefficients
only inside the contradiction hypothesis.  Treating it as an unconditional
\(E\)-value would reverse the logic.
\end{remark}


\section{Finite \texorpdfstring{$\Ein$}{Ein} witnesses and
one-variable Euler traces}
\label{sec:finite-witness-zero-traces}

The pure-exponential base-change theorem also controls finite
representations of constants by algebraic-point \(\Ein\)-values.  We apply
that observation first to \(\Ein(z)-\gamma\), and then to the rank-one
witness family
\(n\Ein(z)-\Ein(z^n)-(n-1)\gamma\).  No appeal to Picard's theorem or to a
numerical zero count is needed.

\subsection{Rigidity of finite algebraic-point witnesses}

Put
\[
 \mathscr V=\operatorname{span}_{\Qbar}
 \{\Ein(\alpha):\alpha\in\Qbar^\times\}\subset\C,
\]
and let \(\Eexp^{\mathrm{alg}}\) denote the relative algebraic closure of
\(\Eexp\) in \(\C\).

\begin{proposition}[Finite-witness rigidity]
\label{prop:finite-ein-witness-rigidity}
The linear map
\[
 \bigoplus_{\alpha\in\Qbar^\times}\Qbar e_\alpha
 \longrightarrow\C,
 \qquad
 \sum_\alpha c_\alpha e_\alpha\longmapsto
 \sum_\alpha c_\alpha\Ein(\alpha),
\]
where the direct sum means finite support, is injective.  Moreover
\begin{equation}\label{eq:V-Eexp-intersection}
 \mathscr V\cap\Eexp^{\mathrm{alg}}=\{0\}.
\end{equation}
Consequently, if \(\gamma\in\mathscr V\), then its finite
\(\Ein\)-representation is unique and \(\gamma\) is transcendental over
\(\Eexp\).
\end{proposition}

\begin{proof}
Every finite family of distinct nonzero algebraic-point values
\(\Ein(\alpha)\) is algebraically independent over \(\Eexp\) by
Theorem~\ref{thm:pure-exp-base}; it is therefore linearly independent.
This proves injectivity.  A nonzero linear form in algebraically independent
indeterminates is itself transcendental over the base field: extend that
linear form to an invertible linear coordinate system.  Hence every nonzero
element of \(\mathscr V\) is transcendental over \(\Eexp\), proving
\eqref{eq:V-Eexp-intersection} and the final assertion.
\end{proof}

Let \(q_*\in(0,1)\) be the unique real solution of
\begin{equation}\label{eq:qstar-ein-level}
 \Ein(q_*)=\gamma.
\end{equation}
Indeed, \(\Ein'(x)=(1-e^{-x})/x>0\) on \(\mathbb R\), after filling the
removable value at \(0\), while
\(\Ein(0)-\gamma=-\gamma<0\) and
\(\Ein(1)-\gamma=\delta/e>0\).

\begin{theorem}[The one-variable Euler-level dichotomy]
\label{thm:qstar-algebraic-dichotomy}
Exactly one of the following alternatives holds.
\begin{enumerate}[label=\textup{(\alph*)}]
 \item Every zero of \(\Ein(z)-\gamma\) is transcendental.
 \item The number \(q_*\) is algebraic and is the unique algebraic zero of
       \(\Ein(z)-\gamma\).  In this case \(\gamma\) and \(\delta\) are
       algebraically independent over \(\Eexp\).
\end{enumerate}
Under alternative \textup{(b)}, the six-constant field
\[
 \Qbar\bigl(\gamma,\delta,E_1(1),\Ei(1),\Ci(1),\Chi(1)\bigr)
\]
has transcendence degree \(5\), and its complete relation ideal is
generated by
\begin{equation}\label{eq:qstar-six-relation}
 \Ei(1)-E_1(1)-2\Chi(1).
\end{equation}
\end{theorem}

\begin{proof}
Two distinct algebraic zeros \(a,b\) would give
\(\Ein(a)=\Ein(b)\), contradicting
Theorem~\ref{thm:pure-exp-base}.  A nonreal algebraic zero is also
impossible, since its distinct conjugate would be a second algebraic zero.
Thus the only possible algebraic zero is the unique real zero \(q_*\).

Assume that \(q_*\) is algebraic and put \(Y_a=\Ein(a)\).  The values
\(Y_{q_*}\) and \(Y_1\) are independent over \(\Eexp\), and
\[
 \gamma=Y_{q_*},\qquad
 \frac{\delta}{e}=Y_1-Y_{q_*}.
\]
This invertible linear change proves the asserted independence of
\(\gamma,\delta\) over \(\Eexp\).

For the last assertion, put
\[
 A=\Ein(1),\qquad B=\Ein(-1),\qquad
 C=\frac{\Ein(i)+\Ein(-i)}2.
\]
The five arguments \(q_*,1,-1,i,-i\) are distinct.  Pure-exponential
base change shows that \(\gamma=\Ein(q_*),A,B,C\) are independent over
\(\Eexp\); together with the transcendence of \(e\), the five numbers
\(e,\gamma,A,B,C\) are algebraically independent over \(\Qbar\).
The six-constant field is \(\Qbar(e,\gamma,A,B,C)\).  Identity
\eqref{eq:qstar-six-relation} is therefore a height-one prime relation and
generates the complete kernel.
\end{proof}

\subsection{The rank-one witness family}

For \(n\ge2\), define
\begin{equation}\label{eq:Gn-witness-family}
 G_n(z)=n\Ein(z)-\Ein(z^n)-(n-1)\gamma.
\end{equation}

\begin{proposition}[Real-zero classification for \(G_n\)]
\label{prop:Gn-real-zero-classification}
For every \(n\ge2\), the function \(G_n\) has exactly one positive real zero
\(\beta_n\in(0,1)\).  If \(n\) is even, it has no negative real zero.  If
\(n\) is odd, it has exactly one negative real zero
\(\lambda_n\in(-\infty,-1)\).  All these real zeros are simple.
\end{proposition}

\begin{proof}
For \(x\ne0\), direct differentiation gives
\begin{equation}\label{eq:Gn-real-derivative}
 G_n'(x)=\frac n{x}\bigl(e^{-x^n}-e^{-x}\bigr).
\end{equation}
Thus \(G_n\) is strictly increasing on \((0,1)\).  Since
\[
 G_n(0)=-(n-1)\gamma<0,
 \qquad
 G_n(1)=\frac{(n-1)\delta}{e}>0,
\]
there is exactly one zero \(\beta_n\) in that interval.  For \(x>1\),
\[
 G_n(x)=nE_1(x)-E_1(x^n)>(n-1)E_1(x)>0,
\]
so there are no further positive zeros.

If \(n\) is even and \(x<0\), then \(x^n>0>x\); hence the numerator in
\eqref{eq:Gn-real-derivative} is negative and \(G_n'(x)>0\).  Moreover
\(G_n(x)\to-\infty\) as \(x\to-\infty\), while
\(G_n(0)=-(n-1)\gamma<0\).  Strict monotonicity therefore excludes a
negative zero.  The asserted limit also follows directly, on writing
\(x=-t\), from
\[
 G_n(-t)=-n\Ei(t)-E_1(t^n)\qquad(t>0).
\]

If \(n\) is odd, \eqref{eq:Gn-real-derivative} shows that \(G_n\) is
strictly decreasing on \((-\infty,-1)\) and strictly increasing on
\((-1,0)\).  One has
 \begin{align*}
 G_n(x)&\longrightarrow+\infty\quad(x\to-\infty),\\
 G_n(-1)&=-(n-1)\Ei(1)<0,\qquad
 G_n(0)=-(n-1)\gamma<0.
 \end{align*}
Here the first limit follows from
\[
 G_n(-t)=\Ei(t^n)-n\Ei(t)\qquad(t>0)
\]
and the standard positive-axis asymptotic for \(\Ei\).
This gives exactly one zero in \((-\infty,-1)\) and none in \((-1,0)\).
Formula \eqref{eq:Gn-real-derivative} is nonzero at each listed root, so
they are simple.
\end{proof}

\begin{theorem}[Algebraic-zero rigidity for \(G_n\)]
\label{thm:Gn-algebraic-zero-rigidity}
For every \(n\ge2\), the function \(G_n\) has at most one algebraic zero,
and every algebraic zero is real.  Consequently:
\begin{enumerate}[label=\textup{(\roman*)}]
 \item if \(n\) is even, either every zero is transcendental or
       \(\beta_n\) is the unique algebraic zero;
 \item if \(n\) is odd, every nonreal zero is transcendental and at most
       one of \(\beta_n,\lambda_n\) is algebraic.
\end{enumerate}
If \(\alpha\) is an algebraic zero and
\(D_\alpha=\Ein(\alpha)-\Ein(\alpha^n)\), then
\begin{equation}\label{eq:Gn-witness-amplification}
 \gamma,\quad\delta,\quad D_\alpha
 \quad\text{are algebraically independent over }\Eexp.
\end{equation}
The six-constant field then has transcendence degree \(5\), with complete
relation ideal generated by \eqref{eq:qstar-six-relation}.
\end{theorem}

\begin{proof}
Suppose that distinct algebraic numbers \(\alpha,\beta\) were zeros.  They
are nonzero, and subtraction gives
\[
 n\Ein(\alpha)-\Ein(\alpha^n)
 =n\Ein(\beta)-\Ein(\beta^n).
\]
After collecting equal arguments, this is a linear relation among distinct
algebraic-point \(\Ein\)-values.  Its coefficient at
\(\Ein(\alpha)\) is
\[
 n-\mathbf1_{\alpha^n=\alpha}
   +\mathbf1_{\beta^n=\alpha}\ge n-1>0,
\]
because \(\beta\ne\alpha\).  This contradicts
Proposition~\ref{prop:finite-ein-witness-rigidity}.  Hence there is at most
one algebraic zero.  Since \(G_n\) has real coefficients, a nonreal
algebraic zero and its conjugate would give two, so every algebraic zero is
real.  Proposition~\ref{prop:Gn-real-zero-classification} now gives the two
alternatives.

For either possible algebraic real zero, the three arguments
\(\alpha,\alpha^n,1\) are distinct.  The corresponding three
\(\Ein\)-values are independent over \(\Eexp\), while
\[
 \gamma=\frac{n\Ein(\alpha)-\Ein(\alpha^n)}{n-1},\qquad
 \frac\delta e=\Ein(1)-\gamma,\qquad
 D_\alpha=\Ein(\alpha)-\Ein(\alpha^n).
\]
This is an invertible linear change of three coordinates, proving
\eqref{eq:Gn-witness-amplification}.  Moreover
\(\{\alpha,\alpha^n\}\) is disjoint from \(\{1,-1,i,-i\}\), so the final
six-constant assertion follows exactly as in
Theorem~\ref{thm:qstar-algebraic-dichotomy}.
\end{proof}

\subsection{Two exact traces of the one-variable Euler divisor}

Let \(Z_\gamma\) be the multiset of zeros of
\(F_\gamma(z)=\Ein(z)-\gamma\), and put
\begin{equation}\label{eq:one-variable-traces}
 S_m=\sum_{\rho\in Z_\gamma}\rho^{-m}\qquad(m\ge2).
\end{equation}
The sums converge absolutely because \(F_\gamma\) is entire of order one.

\begin{theorem}[Trace recovery for \(\Ein-\gamma\)]
\label{thm:ein-gamma-trace-recovery}
One has
\begin{align}
 S_2&=\frac1{\gamma^2}-\frac1{2\gamma},
 \label{eq:ein-gamma-S2}\\
 S_3&=\frac1{\gamma^3}-\frac3{4\gamma^2}+\frac1{6\gamma}.
 \label{eq:ein-gamma-S3}
\end{align}
Conversely,
\begin{equation}\label{eq:ein-gamma-trace-inversion}
 \boxed{\displaystyle
 \gamma=\frac{24S_2+1}{24S_3+6S_2}.}
\end{equation}
In particular
\[
 \Qbar(S_2,S_3)=\Qbar(\gamma),
\]
and the two traces satisfy
\begin{equation}\label{eq:S2-S3-curve}
 144S_2^3+21S_2^2+S_2-144S_3^2+3S_3=0.
\end{equation}
Moreover \(S_2\in\Qbar\) if and only if \(\gamma\in\Qbar\), and the same
equivalence holds with \(S_3\) in place of \(S_2\).
\end{theorem}

\begin{proof}
A genus-one Hadamard factorization of \(F_\gamma\) gives, for \(m\ge2\),
the inverse power sums from the Taylor coefficients of
\(\log(F_\gamma(z)/F_\gamma(0))\).  The origin jet is
\[
 F_\gamma(0)=-\gamma,\quad
 F_\gamma'(0)=1,\quad
 F_\gamma''(0)=-\frac12,\quad
 F_\gamma'''(0)=\frac13.
\]
Hence
\[
 S_2=-\left(\log F_\gamma\right)''(0),
 \qquad
 S_3=-\frac12\left(\log F_\gamma\right)'''(0),
\]
which gives \eqref{eq:ein-gamma-S2} and
\eqref{eq:ein-gamma-S3}.  With \(x=1/\gamma\), these read
\[
 S_2=x^2-\frac x2,
 \qquad
 S_3=x^3-\frac34x^2+\frac x6.
\]
Eliminating \(x^2\) gives
\[
 x\left(S_2+\frac1{24}\right)=S_3+\frac{S_2}{4},
\]
which proves \eqref{eq:ein-gamma-trace-inversion}; its denominator is
nonzero because \(x=1/\gamma>1\).  A resultant calculation, or direct
substitution, gives \eqref{eq:S2-S3-curve}.  Finally, each \(S_j\) is a
nonconstant polynomial in \(x\) over \(\Qbar\), so algebraicity of either
trace forces \(x\), hence \(\gamma\), to be algebraic.  The converse is
immediate.
\end{proof}

\begin{remark}[All higher traces]
The same calculation gives
\[
 S_m=-m[z^m]\log\left(
  1+\frac1\gamma\sum_{r\ge1}
       \frac{(-1)^r z^r}{r\,r!}\right)
 \qquad(m\ge2).
\]
Thus every higher inverse trace is a rational polynomial in \(1/\gamma\).
Unlike the phase recovery from the zero set of the \(\gamma\)-free
reciprocal function, this is a spectral reparametrization of a
\(\gamma\)-dependent level.
\end{remark}


\section{Euler zero lifting and denominator explosion}

Fix \(0\ne\xi\in\Z\) and set
\[
 \Eul_\xi(z)=\sum_{n\ge0}(-1)^n n!\xi^nz^n,
 \qquad
 S_n(\xi)=\sum_{k=0}^n(-1)^k k!\xi^k.
\]
Coefficient comparison gives the exact identity
\begin{equation}\label{eq:quotient-coeff}
 \frac{\Eul_\xi(z)-\alpha}{z-1}
 =\sum_{n\ge0}n!b_nz^n,\qquad
 b_n=\frac{\alpha-S_n(\xi)}{n!}.
\end{equation}

\begin{theorem}\label{thm:denominator}
Let \(K\) contain \(\alpha\), and let \(\Delta_N\) be the least positive
integer with \(\Delta_Nb_0,\ldots,\Delta_Nb_N\in\OO_K\).  Then
\[
 \lim_{N\to\infty}\frac{\log\Delta_N}{N}=+\infty.
\]
The series in \eqref{eq:quotient-coeff} is therefore not arithmetic Gevrey of
order one.
\end{theorem}

\begin{proof}
Suppose first that \(\alpha\notin\Q\).  There is a \(\Q\)-linear functional
\(\ell:K\to\Q\), a positive integer \(M\), and a nonzero \(r\in\Q\) such that
\[
 \ell(1)=0,\qquad \ell(\alpha)=r,\qquad
 M\ell(\OO_K)\subset\Z.
\]
Applying \(M\ell\) to \(\Delta_Nb_N\in\OO_K\) gives
\(M\Delta_Nr/N!\in\Z\).  After clearing the fixed numerator and denominator
of \(Mr\), this implies
\[
 \Delta_N\ge \frac{N!}{C}
\]
for a fixed \(C>0\).  Hence \(\log\Delta_N/N\to\infty\).

Now let \(\alpha=a/b\in\Q\) in lowest terms.  For every prime \(p\nmid b\),
the series
\[
 W_{p,\xi}=\sum_{k\ge0}(-1)^kk!\xi^k
\]
converges in \(\Q_p\).  Call \(p\) \emph{bad} if
\(W_{p,\xi}\ne a/b\).  At a bad prime,
\(v_p(a/b-S_N(\xi))\) is eventually constant, and therefore
\begin{equation}\label{eq:bad-local}
 v_p(\Delta_N)\ge v_p(N!)-O_p(1).
\end{equation}
We claim that
\begin{equation}\label{eq:bad-diverge}
 \sum_{p\ \mathrm{bad}}\frac{\log p}{p-1}=\infty.
\end{equation}
Indeed, otherwise let \(\mathcal P\) be the set of good primes, after deleting
the finitely many primes dividing \(b\).  The product formula and
\(v_p(N!)\le N/(p-1)\) give, for a fixed \(C_\xi>0\),
\[
 C_\xi^N N!\prod_{p\in\mathcal P}|N!|_p^2
 \le \frac{C^N}{N!}\longrightarrow0.
\]
Theorem~1 of Matala-aho--Zudilin \cite{MZ2018}, applied at their variable
\(\xi\), says that either one \(W_{p,\xi}\), \(p\in\mathcal P\), is
irrational or two such values are distinct.  Both alternatives contradict
\(W_{p,\xi}=a/b\) for every \(p\in\mathcal P\).  Thus
\eqref{eq:bad-diverge} holds.

For any finite set \(B\) of bad primes, \eqref{eq:bad-local} yields
\[
 \liminf_{N\to\infty}\frac{\log\Delta_N}{N}
 \ge \sum_{p\in B}\frac{\log p}{p-1}.
\]
Letting \(B\) increase proves the theorem.  Finally,
\(\Eul_\xi\) is holonomic and division by \(z-1\) preserves holonomy; the
conjugates of the normalized coefficients \(b_n\) have exponential growth.
The superexponential denominator is therefore the precise failed
arithmetic-Gevrey condition.
\end{proof}

\begin{remark}
The theorem is a global statement.  At a good prime the factorial tail makes
the local denominator disappear, while every fixed bad prime contributes the
asymptotic rate \(\log p/(p-1)\).  The global-relation theorem forces the sum
of these bad rates to diverge; compare the stronger residue-class results in
\cite{EMHS2019}.
\end{remark}

\section{A sharp p-adic Pad\'e radius transition}

Matala-aho and Zudilin use \cite{MZ2018}
\begin{align*}
 Q_n(z)&=\sum_{k=0}^n k!\binom nk^2z^k,\\
 Q_n(z)\Eul_1(z)-P_n(z)&=
 n!^2z^{2n}\sum_{k\ge0}(-1)^kk!\binom{n+k}{k}^2z^k.
\end{align*}
Set \(q_n=Q_n(1)\), \(s_n=P_n(1)\), with \(q_0=1,s_0=0\), and
\(r_n(\alpha)=q_n\alpha-s_n\).  Both Pad\'e coordinates satisfy
\[
 u_{n+1}=2(n+1)u_n-n^2u_{n-1}\qquad(n\ge1),
\]
and \(r_1(\alpha)=2\alpha-1\).  Hence
\[
 F_\alpha(z)=\sum_{n\ge0}r_n(\alpha)\frac{z^n}{n!}
\]
satisfies
\begin{equation}\label{eq:pade-ode}
 (1-z)^2F_\alpha'(z)+(z-2)F_\alpha(z)+1=0.
\end{equation}

\begin{theorem}\label{thm:radius}
For a prime \(p\), put
\(\delta_p=\sum_{m\ge0}(-1)^mm!\in\Z_p\).  For every
\(\alpha\in\Q_p\),
\[
 R_p(F_\alpha)=
 \begin{cases}
 p^{1/(p-1)},&\alpha=\delta_p,\\
 p^{-1/(p-1)},&\alpha\ne\delta_p.
 \end{cases}
\]
Furthermore \(F_\alpha\in\Z_p[[z]]\) exactly in the matched case
\(\alpha=\delta_p\).
\end{theorem}

\begin{proof}
In the matched case the Pad\'e remainder identity gives
\[
 r_n(\delta_p)
 =n!^2\sum_{k\ge0}(-1)^kk!\binom{n+k}{k}^2,
\]
so \(r_n(\delta_p)/n!\in n!\Z_p\).  Thus the radius is at least
\(\rho_p=p^{1/(p-1)}\), and the coefficients are integral.

To see exactness, expand a hypothetical analytic solution of
\eqref{eq:pade-ode} at \(z=1\).  Writing \(t=z-1\), the unique formal
solution regular there is
\[
 \sum_{m\ge0}m!t^m,
\]
whose \(p\)-adic radius is exactly \(\rho_p\).  Since nonarchimedean discs of
radius greater than one centered at \(0\) and \(1\) coincide, a radius larger
than \(\rho_p\) for \(F_{\delta_p}\) would contradict this local expansion.

The normalized homogeneous solution of \eqref{eq:pade-ode} is
\begin{equation}\label{eq:homogeneous}
 H(z)=\frac{1}{1-z}\exp\left(\frac{z}{1-z}\right)
 =\sum_{n\ge0}q_n\frac{z^n}{n!}.
\end{equation}
Therefore
\[
 F_\alpha-F_{\delta_p}=(\alpha-\delta_p)H.
\]
For \(n\equiv0\pmod p\), the formula for \(q_n\) gives
\(q_n\equiv1\pmod p\): terms \(1\le k<p\) have
\(p\mid\binom nk\), while terms \(k\ge p\) contain \(p\mid k!\).
It follows that \(H\) has exact radius \(p^{-1/(p-1)}\).  If
\(\alpha\ne\delta_p\), this smaller-radius term dominates
\(F_{\delta_p}\), proving the second case.  Along the same subsequence the
coefficient valuations are unbounded below; in particular no fixed scalar
makes the full series integral.  Conversely the matched coefficients were
already shown integral.
\end{proof}

The preceding theorem is the \(\xi=1\) instance of the all-prime,
all-parameter result below.  We retain it as a transparent model calculation;
the next section supplies the general Laguerre recurrence proof and removes
every coprimality restriction.

% Standalone section fragment.  It assumes amsmath, amssymb, and amsthm,
% together with theorem, proposition, corollary, and remark environments.

\section{The all-prime Pad\'e transition and its arithmetic obstructions}
\label{sec:all-prime-pade}

Fix $0\ne\xi\in\mathbb Z$ and write
\[
 \mathcal E(u)=\sum_{m\geq0}(-1)^m m!u^m,
 \qquad
 \delta_p(\xi)=\sum_{m\geq0}(-1)^m m!\xi^m\in\mathbb Z_p.
\]
We use the classical Pad\'e pair~\cite{MZ2018}
\begin{align}
 Q_n(u)&=\sum_{k=0}^n k!\binom nk^2u^k,\label{eq:Qn-def}\\
 Q_n(u)\mathcal E(u)-P_n(u)
 &=n!^2u^{2n}\sum_{k\geq0}(-1)^k k!
       \binom{n+k}{k}^2u^k.\label{eq:pade-rem}
\end{align}
Here $Q_0=1$ and $P_0=0$.  Put
\[
 q_n(\xi)=Q_n(\xi),\qquad s_n(\xi)=P_n(\xi),\qquad
 F_{\alpha}^{(\xi)}(z)
 =\sum_{n\geq0}\bigl(q_n(\xi)\alpha-s_n(\xi)\bigr)\frac{z^n}{n!}.
\]

\begin{lemma}[Pad\'e recurrence and affine equation]
\label{lem:pade-recurrence-all-prime}
For $n\geq1$, both $q_n(\xi)$ and $s_n(\xi)$ satisfy
\begin{equation}
 u_{n+1}=\bigl(1+(2n+1)\xi\bigr)u_n-n^2\xi^2u_{n-1}.
\label{eq:pade-three-term-xi}
\end{equation}
Consequently $F=F_\alpha^{(\xi)}$ satisfies
\begin{equation}
 (1-\xi z)^2F'(z)-(1+\xi-\xi^2z)F(z)+1=0.
\label{eq:affine-pade-ode}
\end{equation}
Its normalized homogeneous solution is
\begin{equation}
 H_\xi(z)=\frac{1}{1-\xi z}
 \exp\!\left(\frac{z}{1-\xi z}\right)
 =\sum_{n\geq0}q_n(\xi)\frac{z^n}{n!}.
\label{eq:H-xi}
\end{equation}
\end{lemma}

\begin{proof}
The polynomial identity
\[
 Q_n(u)=n!u^nL_n(-1/u)
\]
and the Laguerre recurrence give
\[
 Q_{n+1}(u)=\bigl(1+(2n+1)u\bigr)Q_n(u)-n^2u^2Q_{n-1}(u).
\]
The Pad\'e determinant identity
\begin{equation}
 D_n(u):=Q_n(u)P_{n+1}(u)-Q_{n+1}(u)P_n(u)=n!^2u^{2n}
\label{eq:pade-casoratian}
\end{equation}
then yields, as a polynomial identity,
\begin{align*}
 Q_n(u)&\bigl(P_{n+1}(u)-(1+(2n+1)u)P_n(u)
                    +n^2u^2P_{n-1}(u)\bigr)\\
 &=D_n(u)-n^2u^2D_{n-1}(u)=0.
\end{align*}
Thus $P_n$ obeys the same recurrence.  Evaluation at $u=\xi$ proves
\eqref{eq:pade-three-term-xi}.

For $r_n=q_n(\xi)\alpha-s_n(\xi)$ one has
$r_0=\alpha$ and $r_1=(1+\xi)\alpha-1$.  Multiplying
\eqref{eq:pade-three-term-xi} by $z^n/n!$ and summing gives
\eqref{eq:affine-pade-ode}.  Solving its homogeneous part with value $1$
at $z=0$ gives \eqref{eq:H-xi}.
\end{proof}

\begin{theorem}[All-prime radius transition]
\label{thm:all-prime-radius}
Let $p$ be any prime, set $s=v_p(\xi)$, and let $\alpha\in\mathbb Q_p$.
Then
\begin{equation}
 R_p\!\left(F_\alpha^{(\xi)}\right)=
 \begin{cases}
 p^{\frac1{p-1}+2s}
   =p^{1/(p-1)}|\xi|_p^{-2},&\alpha=\delta_p(\xi),\\[2mm]
 p^{-1/(p-1)},&\alpha\ne\delta_p(\xi).
 \end{cases}
\label{eq:all-prime-radius}
\end{equation}
Moreover
\begin{equation}
 F_\alpha^{(\xi)}\in\mathbb Z_p[[z]]
 \quad\Longleftrightarrow\quad
 \alpha=\delta_p(\xi).
\label{eq:integrality-matching}
\end{equation}
If the right-hand condition fails, the valuations of the coefficients of
$F_\alpha^{(\xi)}$ are unbounded below; in particular no fixed nonzero
scalar makes the whole series integral.
\end{theorem}

\begin{proof}
Specializing \eqref{eq:pade-rem} at $u=\xi$ gives
\begin{equation}
 \frac{q_n(\xi)\delta_p(\xi)-s_n(\xi)}{n!}
 =n!\xi^{2n}\sum_{k\geq0}(-1)^k k!
          \binom{n+k}{k}^2\xi^k
 \in n!\xi^{2n}\mathbb Z_p.
\label{eq:matched-coefficient-divisibility}
\end{equation}
Hence the matched radius is at least
$p^{1/(p-1)+2s}$.

Let $c=1/\xi$ and $t=z-c$.  Equation
\eqref{eq:affine-pade-ode} becomes
\begin{equation}
 \xi^2t^2F'(t)+(\xi^2t-1)F(t)+1=0.
\label{eq:local-affine-ode}
\end{equation}
Its unique formal solution regular at $t=0$ is
\begin{equation}
 F_{\mathrm{reg}}(t)=\sum_{m\geq0}m!\xi^{2m}t^m,
\label{eq:local-regular-solution}
\end{equation}
whose radius is exactly $p^{1/(p-1)+2s}$.  This number is strictly larger
than $|c|_p=p^s$.  The lower bound from
\eqref{eq:matched-coefficient-divisibility} therefore places $c$ inside the
disc of the matched solution, and its re-expansion at $c$ must be
\eqref{eq:local-regular-solution}.  If its radius at $0$ were larger than
$p^{1/(p-1)+2s}$, the ultrametric discs with centers $0$ and $c$ and that
larger radius would coincide, contradicting the exact radius of
\eqref{eq:local-regular-solution}.  This proves the first case of
\eqref{eq:all-prime-radius}.

The coefficients of $H_\xi$ have radius at least $p^{-1/(p-1)}$, since
$q_n(\xi)\in\mathbb Z$.  If $p\mid\xi$, then
$q_n(\xi)\equiv1\pmod p$ for every $n$.  If $p\nmid\xi$ and
$n\equiv0\pmod p$, then the terms in \eqref{eq:Qn-def} with
$1\leq k<p$ contain $p\mid\binom nk$, while those with $k\geq p$ contain
$p\mid k!$; hence again $q_n(\xi)\equiv1\pmod p$.  Along the indicated
subsequence, Legendre's formula gives
\[
 \frac1n v_p\!\left(\frac{q_n(\xi)}{n!}\right)
 \longrightarrow-\frac1{p-1}.
\]
Thus $R_p(H_\xi)=p^{-1/(p-1)}$.  Finally,
\[
 F_\alpha^{(\xi)}-F_{\delta_p(\xi)}^{(\xi)}
   =(\alpha-\delta_p(\xi))H_\xi.
\]
When the scalar is nonzero, the term on the right has strictly smaller
radius than the matched series, proving the second case of
\eqref{eq:all-prime-radius}.  Formula
\eqref{eq:matched-coefficient-divisibility} proves integrality in the
matched case.  In the unmatched case, on the same unit subsequence,
\[
 v_p\!\left((\alpha-\delta_p(\xi))\frac{q_n(\xi)}{n!}\right)
 =v_p(\alpha-\delta_p(\xi))-v_p(n!)\longrightarrow-\infty,
\]
whereas the matched coefficient is integral.  This proves
\eqref{eq:integrality-matching} and the final assertion.
\end{proof}

\begin{proposition}[Local denominator--radius equivalence]
\label{prop:local-denominator-radius}
Put
\[
 S_n(\xi)=\sum_{k=0}^n(-1)^kk!\xi^k,
 \qquad b_n=\frac{\alpha-S_n(\xi)}{n!},
\]
and, for $\alpha\in\mathbb Q_p$, define
\[
 d_{p,N}=\max_{0\leq n\leq N}\max\{0,-v_p(b_n)\}.
\]
Then the following are equivalent:
\begin{equation}
 \alpha=\delta_p(\xi),\qquad
 b_n\in\mathbb Z_p\ (n\geq0),\qquad
 d_{p,N}=0\ (N\geq0),\qquad
 R_p(F_\alpha^{(\xi)})>1.
\label{eq:four-local-equivalences}
\end{equation}
If $\alpha\ne\delta_p(\xi)$ and
$c_p=v_p(\alpha-\delta_p(\xi))$, then, for all sufficiently large $N$,
\begin{equation}
 d_{p,N}=v_p(N!)-c_p,
 \qquad
 \lim_{N\to\infty}\frac{d_{p,N}}N=\frac1{p-1}.
\label{eq:exact-local-denominator}
\end{equation}
\end{proposition}

\begin{proof}
In the matched case,
\[
 b_n=\sum_{k\geq n+1}(-1)^k\frac{k!}{n!}\xi^k\in\mathbb Z_p.
\]
In the unmatched case, $S_n(\xi)\to\delta_p(\xi)$ and ultrametric
stability gives
$v_p(\alpha-S_n(\xi))=c_p$ for all sufficiently large $n$.  Hence
$-v_p(b_n)=v_p(n!)-c_p$ eventually.  This is nondecreasing apart from
plateaux and ultimately exceeds the maximum contributed by the finitely
many earlier terms, proving \eqref{eq:exact-local-denominator}.  The radius
equivalence is Theorem~\ref{thm:all-prime-radius}.
\end{proof}

\begin{corollary}[Relative adelic radius defect]
\label{cor:adelic-radius-defect}
Let $\alpha=a/b\in\mathbb Q$ in lowest terms and define
\[
 \mathcal B(\alpha,\xi)
 =\{p:p\nmid b\xi,\ \alpha\ne\delta_p(\xi)\}.
\]
Then
\begin{equation}
 \sum_{p\in\mathcal B(\alpha,\xi)}\frac{\log p}{p-1}=+\infty.
\label{eq:bad-prime-divergence}
\end{equation}
Consequently, with
$\rho_{p,\xi}=p^{1/(p-1)}|\xi|_p^{-2}$,
\begin{equation}
 \prod_{p\nmid b\xi}
 \frac{R_p(F_\alpha^{(\xi)})}{\rho_{p,\xi}}=0,
\label{eq:relative-adelic-product-zero}
\end{equation}
where the product is the limit over finite sets of primes.
\end{corollary}

\begin{proof}
Suppose that the sum in \eqref{eq:bad-prime-divergence} were finite, and let
$\mathcal P$ be the primes $p\nmid b\xi$ for which
$\delta_p(\xi)=\alpha$.  The complement of $\mathcal P$ then has finite
$\sum_p\log p/(p-1)$.  Since
\[
 \prod_{p\in\mathcal P}|n!|_p^2
 =\frac1{(n!)^2}\prod_{p\notin\mathcal P}p^{2v_p(n!)}
 \leq\frac{C^n}{(n!)^2}
\]
for a fixed $C>0$, the hypothesis of the global-relation theorem of
Matala-aho and Zudilin~\cite{MZ2018} is satisfied.  That
theorem forbids all values $\delta_p(\xi)$, $p\in\mathcal P$, from being the
same rational number, a contradiction.  Thus
\eqref{eq:bad-prime-divergence} holds.  For $p\nmid b\xi$, Theorem
\ref{thm:all-prime-radius} gives a local ratio $1$ at a matched prime and
$p^{-2/(p-1)}$ at an unmatched prime.  Taking logarithms proves
\eqref{eq:relative-adelic-product-zero}.
\end{proof}

\begin{remark}[What the adelic product does and does not say]
The product in \eqref{eq:relative-adelic-product-zero} is normalized by the
matched local benchmark.  It is not itself an application of the
Borel--Dwork--P\'olya--Bertrandias rationality criterion.  Rather, it records
why the direct application fails here: for rational $\alpha$ there are
infinitely many bad integral lattices, with divergent total radius defect.
It does not exclude a different auxiliary series or different adelic
domains; compare the original Dwork criterion~\cite{Dwork1960} and its
capacity-theoretic extension~\cite{BostChambertLoir2009}.
\end{remark}

\subsection{Frobenius and Borel obstructions}

\begin{proposition}[Exact $p$-curvature]
\label{prop:exact-p-curvature}
Let $u=1-\xi z$.  The homogeneous equation in
Lemma~\ref{lem:pade-recurrence-all-prime} is the rank-one connection
\[
 \nabla_\partial=\partial-a_\xi(z),
 \qquad
 a_\xi(z)=\frac{\xi}{u}+\frac1{u^2}.
\]
At every prime $p$, its reduction has $p$-curvature
\begin{equation}
 \psi_p(\partial)=-\frac1{(1-\bar\xi z)^{2p}},
\label{eq:exact-p-curvature}
\end{equation}
up to the conventional overall sign.  In particular it is nonzero, hence
not nilpotent, at every prime.
\end{proposition}

\begin{proof}
In characteristic $p$, the rank-one Jacobson formula gives
\[
 (\partial-a_\xi)^p
 =-\bigl(\partial^{p-1}a_\xi+a_\xi^p\bigr),
\]
because $\partial^{[p]}=0$ on $\mathbb F_p(z)$.  Wilson's congruence gives
\[
 \partial^{p-1}(\xi u^{-1})=-\xi^pu^{-p},
 \qquad
 \partial^{p-1}(u^{-2})=p!\xi^{p-1}u^{-p-1}=0,
\]
whereas Frobenius gives
\[
 a_\xi^p=\xi^pu^{-p}+u^{-2p}.
\]
Substitution proves \eqref{eq:exact-p-curvature}.
\end{proof}

\begin{corollary}[Global-nilpotence no-go]
\label{cor:global-nilpotence-no-go}
The minimal rank-one operator annihilating $H_\xi$ is not globally
nilpotent.  Consequently it is not a subquotient of an ordinary smooth
proper Gauss--Manin connection and is not a $G$-operator in the usual
globally nilpotent category.
\end{corollary}

\begin{proof}
For a rank-one connection, nilpotence of $p$-curvature is equivalent to its
vanishing.  Proposition~\ref{prop:exact-p-curvature} therefore rules out
global nilpotence.  The stated consequences are the standard
Gauss--Manin/global-nilpotence theorems of Katz~\cite{Katz1970} and the
Andr\'e--Chudnovsky--Katz theory of $G$-operators~\cite{Andre1989}.
\end{proof}

\begin{remark}
This is deliberately a statement about the ordinary globally nilpotent
category.  It does not rule out every notion of irregular exponential
geometric origin.  Moreover the connection is independent of $\alpha$:
$p$-curvature cannot distinguish the matched initial value.  Matching is a
property of an overconvergent horizontal lift, not of the differential
module alone.
\end{remark}

\begin{proposition}[Dieudonn\'e--Dwork defect]
\label{prop:dwork-defect}
Put $w(z)=z/(1-\xi z)$.  For every prime $p$,
\begin{equation}
 [z^p]\bigl(w(z^p)-pw(z)\bigr)=1-p\xi^{p-1}\notin p\mathbb Z_p.
\label{eq:dwork-defect-coefficient}
\end{equation}
Consequently neither $\exp(w(z))$ nor $H_\xi(z)$ belongs to
$\mathbb Z_p[[z]]$.
\end{proposition}

\begin{proof}
The Dieudonn\'e--Dwork criterion~\cite{Dieudonne1957} says that, for
$S\in z\mathbb Q_p[[z]]$,
\[
 \exp(S)\in1+z\mathbb Z_p[[z]]
 \quad\Longleftrightarrow\quad
 S(z^p)-pS(z)\in pz\mathbb Z_p[[z]].
\]
Since $w(z)=\sum_{m\geq1}\xi^{m-1}z^m$, equation
\eqref{eq:dwork-defect-coefficient} follows immediately and violates the
criterion.  Finally $(1-\xi z)^{-1}$ and its inverse both lie in
$\mathbb Z_p[[z]]$, so multiplication by this factor cannot repair
integrality.
\end{proof}

\begin{proposition}[Fixed Borel logarithm]
\label{prop:fixed-borel-logarithm}
Let
\[
 \frac{E_\xi(z)-\alpha}{z-1}=\sum_{n\geq0}n!b_nz^n,
 \qquad
 E_\xi(z)=\sum_{n\geq0}(-1)^nn!\xi^nz^n,
\]
and let $\mathcal B_1$ denote the order-one formal Borel transform.  Then
\[
 B_{\alpha,\xi}(t):=
 \mathcal B_1\!\left(\frac{E_\xi-\alpha}{z-1}\right)(t)
 =\sum_{n\geq0}b_nt^n
\]
satisfies
\begin{equation}
 B_{\alpha,\xi}'-B_{\alpha,\xi}
 =\frac{\xi}{(1+\xi t)^2},
 \qquad B_{\alpha,\xi}(0)=\alpha-1.
\label{eq:borel-affine-ode}
\end{equation}
If $t_0=-1/\xi$ and $x=t-t_0$, then on every local complex branch
\begin{equation}
 B_{\alpha,\xi}(t)
 =-\frac{e^x}{\xi x}-\frac{e^x}{\xi}\log x+e^xh_{\alpha,\xi}(x),
 \qquad h_{\alpha,\xi}\in\mathbb C\{x\}.
\label{eq:fixed-borel-log}
\end{equation}
In particular no choice of $\alpha$ makes this Borel transform meromorphic
at $t_0$, and none makes it rational.
\end{proposition}

\begin{proof}
Coefficient comparison gives
$b_n=(\alpha-S_n(\xi))/n!$ and hence
\[
 (n+1)b_{n+1}=b_n+(n+1)(-1)^n\xi^{n+1}.
\]
Summation proves \eqref{eq:borel-affine-ode}, equivalently
\[
 B_{\alpha,\xi}(t)=e^t\left(
 \alpha-1+\int_0^t\frac{\xi e^{-v}}{(1+\xi v)^2}\,dv\right).
\]
At $t=t_0+x$, equation \eqref{eq:borel-affine-ode} reads
$B'-B=1/(\xi x^2)$.  Equivalently,
\[
 (e^{-x}B)'=\frac{e^{-x}}{\xi x^2}.
\]
Since
$e^{-x}x^{-2}=x^{-2}-x^{-1}+\text{a holomorphic germ}$, integration and
multiplication by $e^x$ give \eqref{eq:fixed-borel-log}.  The homogeneous
correction $Ce^t$, which is the only dependence on $\alpha$, is holomorphic
and cannot cancel either singular term.
\end{proof}

\begin{corollary}[Polynomial reduction in the matched case]
\label{cor:matched-polynomial-reduction}
If $p\nmid\xi$, then
\begin{equation}
 \overline{F_{\delta_p(\xi)}^{(\xi)}}(z)
 =\sum_{m=0}^{p-1}m!\bar\xi^{2m}
       (z-\bar\xi^{-1})^m\in\mathbb F_p[z].
\label{eq:matched-polynomial-mod-p}
\end{equation}
If $p\mid\xi$, then
$\overline{F_{\delta_p(\xi)}^{(\xi)}}(z)=1$.  Thus the matched reduction
has a finite-support, hence $p$-automatic, coefficient sequence.  In the
unmatched case no fixed scalar gives an integral series to reduce.
\end{corollary}

\begin{proof}
For $p\nmid\xi$, reduce the local expansion
\eqref{eq:local-regular-solution}; all terms with $m\geq p$ vanish modulo
$p$, giving \eqref{eq:matched-polynomial-mod-p}.  If $p\mid\xi$, equation
\eqref{eq:matched-coefficient-divisibility} shows that every positive-degree
coefficient is divisible by $p$, while
$F_{\delta_p(\xi)}^{(\xi)}(0)=\delta_p(\xi)\equiv1\pmod p$.  The last two
claims follow from Theorem~\ref{thm:all-prime-radius}; the automaticity
statement is also the elementary polynomial case of Christol's
theorem~\cite{ChristolEtAl1980}.
\end{proof}

\begin{remark}[Kurepa specialization]
For $\xi=-1$,
\[
 \delta_p(-1)\equiv\sum_{m=0}^{p-1}m!=!p\pmod p,
 \qquad
 \overline{F_{\delta_p(-1)}^{(-1)}}(z)
 =\sum_{m=0}^{p-1}m!(z+1)^m.
\]
Thus Kurepa's conjecture is equivalent to the assertion that the constant
term of this matched polynomial is nonzero for every odd prime.  This is an
exact reformulation, not a proof of the conjecture.
\end{remark}

% Bibliography keys required by this fragment:
% MatalaAhoZudilin2018 -- T. Matala-aho and W. Zudilin,
%   Euler's factorial series and global relations, JNT 186 (2018), 202--210.
% Dieudonne1957 -- J. Dieudonne, On the Artin--Hasse exponential series,
%   Proc. AMS 8 (1957), 210--214.
% Katz1970 -- N. Katz, Nilpotent connections and the monodromy theorem,
%   Publ. Math. IHES 39 (1970), 175--232.
% ChristolEtAl1980 -- G. Christol, T. Kamae, M. Mendes France, G. Rauzy,
%   Suites algebriques, automates et substitutions, BSMF 108 (1980), 401--419.
% Dwork1960 -- B. Dwork, On the rationality of the zeta function of an
%   algebraic variety, Amer. J. Math. 82 (1960), 631--648.
% BostChambertLoir2009 -- J.-B. Bost and A. Chambert-Loir, Analytic curves
%   in algebraic varieties over number fields, in Algebra, Arithmetic, and
%   Geometry, Progress in Mathematics 269 (2009), 69--124.
% Andre1989 -- Y. Andre, G-functions and Geometry, Aspects of Mathematics
%   E13, Vieweg, 1989.


\section{Bell--Gould prime-power filters and mixed integrals}
\label{sec:touchard-mixed}

The purpose of this section is twofold.  First we isolate an exact
integer-valued-polynomial mechanism which annihilates the Bell and Gould
functionals simultaneously modulo an arbitrary integer.  The boundary term
in the Gould functional is retained throughout; omitting it makes the
prime-power lift invalid.  Second, we match the resulting reciprocal
integrals with logarithmic moments.  The matching places the answer in the
pure \(E\)-value direction \(e^\lambda\Ein(\lambda)\), where
Theorem~\ref{thm:pure-exp-base} turns the construction into an unconditional
algebraic-independence theorem.

\subsection{The Bell--Gould adjoint calculus}

Let \((\mathsf B_n)_{n\geq0}\) be the Bell sequence and let
\((\mathsf G_n)_{n\geq0}\) be the Gould sequence, normalized by
\[
 \mathsf B_0=1,
 \qquad
 \mathsf B_{n+1}=\sum_{k=0}^n\binom nk\mathsf B_k,
\]
and
\[
 \mathsf G_0=0,
 \quad \mathsf G_1=1,
 \qquad
 \mathsf G_{n+1}=\sum_{k=0}^n\binom nk\mathsf G_k
 \quad(n\geq1).
\]
Define the two integral linear functionals on \(\Z[x]\) by
\[
 \mathfrak B(x^n)=\mathsf B_n,
 \qquad
 \mathfrak G(x^n)=\mathsf G_n.
\]
For \(r\geq1\), write
\((x)_r=x(x-1)\cdots(x-r+1)\), and put
\begin{equation}\label{eq:touchard-theta}
 \Theta_r(f)=
 \sum_{j=0}^{r-1}(-1)^{r-1-j}(r-1-j)!f(j).
\end{equation}

\begin{lemma}[Bell--Gould adjoints]
\label{lem:bell-gould-adjoints}
For every \(r\geq1\) and \(f\in\Z[x]\),
\begin{align}
 \mathfrak B\bigl((x)_rf(x)\bigr)
   &=\mathfrak B\bigl(f(x+r)\bigr),
   \label{eq:bell-adjoint}\\
 \mathfrak G\bigl((x)_rf(x)\bigr)
   &=\mathfrak G\bigl(f(x+r)\bigr)+\Theta_r(f).
   \label{eq:gould-adjoint}
\end{align}
\end{lemma}

\begin{proof}
For \(r=1\), the two identities are precisely the defining binomial
recurrences, with the exceptional initial value \(\mathsf G_1=1\) producing
the term \(f(0)\) in the second identity.  Iteration proves
\eqref{eq:bell-adjoint}.  For the Gould identity, suppose it holds for
\(r\).  Apply the case \(r=1\) to
\(h(x)=(x-1)_rf(x)\), and then apply the induction hypothesis to
\(f(x+1)\).  The new boundary contribution is
\[
 (-1)^rr!f(0)+
 \sum_{j=0}^{r-1}(-1)^{r-1-j}(r-1-j)!f(j+1),
\]
which is \(\Theta_{r+1}(f)\).  This completes the induction.
\end{proof}

Let
\begin{equation}\label{eq:signed-touchard}
 \tau_n(x)=\sum_{k=0}^n
 \left\{\begin{matrix}n\\k\end{matrix}\right\}(-x)^k
\end{equation}
be the signed Touchard polynomial.  The elementary reciprocal integral is
\begin{equation}\label{eq:touchard-remainder}
 \mathscr R_n
 :=\int_0^\infty\frac{e^{-x}\tau_n(x)}{1+x}\,dx
 =\mathsf B_n\delta-\mathsf G_n.
\end{equation}
Indeed, the exponential generating function of the left side is
\[
 \int_0^\infty\frac{e^{-e^z x}}{1+x}\,dx
 =e^{e^z}E_1(e^z).
\]
The function \(\delta\exp(e^z-1)\) minus this generating function has
derivative equal to itself multiplied by \(e^z\), plus \(1\), and hence has coefficients
\((\mathsf G_n)\) with the displayed normalization.

\subsection{A boundary-aware prime-power lift}

Fix a prime \(p\).  Put
\[
 D_p(x)=x^p-x.
\]
There is a \(G_p\in\Z[x]\) such that
\begin{equation}\label{eq:dp-falling}
 D_p(x)=(x)_p+pG_p(x).
\end{equation}
To see this, subtract the two monic degree-\(p\) polynomials.  The difference
has degree at most \(p-1\) and vanishes at all \(p\) elements of
\(\mathbf F_p\), so every coefficient is divisible by \(p\).

For \(r\geq1\), set
\[
 m_{p,r}=\min\{m\geq0:v_p(m!)\geq r\}.
\]
Choose \(a\in\Z\), put \(W_{p,r,a}(x)=(x-a)_{m_{p,r}}\), and define
\begin{equation}\label{eq:touchard-lift-operator}
 \mathcal L_p[P](x)=D_p(x)P(x)-P(x+p).
\end{equation}

For a polynomial \(P\in\Z[x]\), use the abbreviations
\begin{align*}
 \mathrm{FD}_r(P)&:\Longleftrightarrow
      P(k)\in p^r\Z\quad(k\in\Z),\\
 \mathrm{AN}_s(P)&:\Longleftrightarrow
      \mathfrak B(Pf),\mathfrak G(Pf)\in p^s\Z
      \quad(f\in\Z[x]).
\end{align*}

\begin{lemma}[One-step lift]
\label{lem:touchard-one-step}
Suppose that \(1\leq s<r\) and that \(P\) satisfies
\(\mathrm{FD}_r(P)\) and \(\mathrm{AN}_s(P)\).  Then
\(Q=\mathcal L_p[P]\) satisfies
\(\mathrm{FD}_r(Q)\) and \(\mathrm{AN}_{s+1}(Q)\).
\end{lemma}

\begin{proof}
First note the shifted-annihilator consequence of
Lemma~\ref{lem:bell-gould-adjoints}.  If \(r\geq s\), then for every
\(g\in\Z[x]\),
\begin{align*}
 \mathfrak B(P(x+p)g(x))
 &=\mathfrak B(P(x)(x)_pg(x-p)),\\
 \mathfrak G(P(x+p)g(x))
 &=\mathfrak G(P(x)(x)_pg(x-p))
   -\Theta_p(P(x)g(x-p)).
\end{align*}
The first terms on the right are divisible by \(p^s\), and every summand in
the boundary term is divisible by \(p^r\).  Thus both shifted functionals
are divisible by \(p^s\).

The fixed-divisor assertion for \(Q\) follows immediately from
\eqref{eq:touchard-lift-operator}.  Write
\[
 f(x+p)-f(x)=p\Delta_pf(x),
 \qquad \Delta_pf\in\Z[x].
\]
Using \eqref{eq:dp-falling} and
Lemma~\ref{lem:bell-gould-adjoints} gives
\begin{align*}
 \mathfrak B(Qf)
 &=p\mathfrak B(P(x+p)\Delta_pf)
   +p\mathfrak B(PG_pf),\\
 \mathfrak G(Qf)
 &=p\mathfrak G(P(x+p)\Delta_pf)
   +p\mathfrak G(PG_pf)+\Theta_p(Pf).
\end{align*}
The first two terms in each line are divisible by \(p^{s+1}\).  The last
term in the second line is divisible by \(p^r\), hence by \(p^{s+1}\).
\end{proof}

\begin{theorem}[Common Bell--Gould prime-power annihilator]
\label{thm:bell-gould-prime-power}
For every prime \(p\), \(r\geq1\), and \(a\in\Z\), the polynomial
\begin{equation}\label{eq:prime-power-annihilator}
 S_{p,r}^{(a)}(x)=
 \mathcal L_p^{\,r-1}
 \left[W_{p,r,a}(x)(D_p(x)-1)\right]
\end{equation}
is monic and primitive, has degree \(m_{p,r}+rp\), and satisfies
\begin{equation}\label{eq:prime-power-universal}
 p^r\mid\mathfrak B(S_{p,r}^{(a)}f),
 \qquad
 p^r\mid\mathfrak G(S_{p,r}^{(a)}f)
 \quad(f\in\Z[x]).
\end{equation}
\end{theorem}

\begin{proof}
Put \(P_0=W_{p,r,a}(D_p-1)\).  Since
\[
 W_{p,r,a}(k)=m_{p,r}!\binom{k-a}{m_{p,r}}
 \quad(k\in\Z),
\]
the polynomial \(P_0\) satisfies \(\mathrm{FD}_r(P_0)\).  For
\(h=W_{p,r,a}f\), equations
\eqref{eq:bell-adjoint}--\eqref{eq:dp-falling} give
\begin{align*}
 \mathfrak B(P_0f)
 &=\mathfrak B(h(x+p)-h(x))+p\mathfrak B(G_ph),\\
 \mathfrak G(P_0f)
 &=\mathfrak G(h(x+p)-h(x))+p\mathfrak G(G_ph)
   +\Theta_p(h).
\end{align*}
The difference \(h(x+p)-h(x)\) belongs to \(p\Z[x]\), while every value
\(h(j)\) occurring in \(\Theta_p(h)\) is divisible by \(p^r\).  Hence
\(P_0\) satisfies \(\mathrm{AN}_1(P_0)\).  Iterating
Lemma~\ref{lem:touchard-one-step} \(r-1\) times proves
\eqref{eq:prime-power-universal}.

The starting polynomial is monic of degree \(m_{p,r}+p\).
The operator \(\mathcal L_p\) preserves monicity and raises the degree by
\(p\), since its first summand has degree \(p+\deg P\) and its second has
degree \(\deg P\).  This proves the degree formula.  A monic integral
polynomial is primitive.
\end{proof}

The carrier \(W_{p,r,a}\) has the smallest possible degree among primitive
integer polynomials whose fixed divisor is divisible by \(p^r\): the fixed
divisor of a primitive degree-\(d\) polynomial divides \(d!\), while
\(W_{p,r,a}\) attains \(d=m_{p,r}\); see
\cite{PrasadRajkumarReddy2019}.  This is an optimality statement for the
fixed-divisor carrier, not for all possible simultaneous Bell--Gould
annihilators.

\subsection{Coefficientwise CRT for an arbitrary modulus}

\begin{theorem}[Arbitrary-modulus Bell--Gould filter]
\label{thm:bell-gould-arbitrary-modulus}
Let
\[
 M=\prod_{p\mid M}p^{r_p}\geq2,
 \qquad
 d_p=m_{p,r_p}+r_pp,
 \qquad
 D_M=\max_{p\mid M}d_p.
\]
There is a monic primitive \(C_M\in\Z[x]\), of degree \(D_M\), such that
\begin{equation}\label{eq:arbitrary-modulus-universal}
 M\mid\mathfrak B(C_Mf),
 \qquad
 M\mid\mathfrak G(C_Mf)
 \quad(f\in\Z[x]).
\end{equation}
Every nonleading coefficient of \(C_M\) may be chosen to have absolute
value at most \(M/2\).
\end{theorem}

\begin{proof}
For every \(p\mid M\), choose an integer \(a_p\) and put
\[
 T_p(x)=x^{D_M-d_p}S_{p,r_p}^{(a_p)}(x).
\]
All \(T_p\) are monic of degree \(D_M\).  Apply the Chinese remainder
theorem independently to their nonleading coefficients and set the leading
coefficient equal to \(1\).  This gives
\[
 C_M\equiv T_p\pmod{p^{r_p}}\qquad(p\mid M),
\]
with centered representatives bounded by \(M/2\).  Multiplication by a
power of \(x\) preserves \eqref{eq:prime-power-universal}; moreover
\(C_M-T_p\in p^{r_p}\Z[x]\).  Since the two functionals take integral
values on \(\Z[x]\), both values in
\eqref{eq:arbitrary-modulus-universal} are divisible by every
\(p^{r_p}\), and hence by \(M\).
\end{proof}

Let (E) denote the forward shift in (n),
\[
 (Ea)_n=a_{n+1},
\]
acting coefficientwise on polynomial-valued sequences.

\begin{corollary}[Factorial specialization]
\label{cor:bell-gould-factorial}
For \(M=N!\), \(N\geq2\), one may take a monic common filter of degree
\(D_N\) satisfying
\begin{equation}\label{eq:factorial-degree-bound}
 3N-2\lfloor\log_2N\rfloor-3
 \leq D_N\leq3N.
\end{equation}
Consequently, for every \(n\geq0\), there are integers \(P_{N,n},Q_{N,n}\)
such that
\begin{equation}\label{eq:factorial-touchard-integral}
 \frac1{N!}\int_0^\infty
 \frac{e^{-x}[C_{N!}(E)\tau_n](x)}{1+x}\,dx
 =Q_{N,n}\delta-P_{N,n}.
\end{equation}
\end{corollary}

\begin{proof}
If \(r_p=v_p(N!)\), then
\[
 m_{p,r_p}=p\lfloor N/p\rfloor,
 \qquad
 d_p=p\lfloor N/p\rfloor+pv_p(N!).
\]
The quotient of \(N!\) by
\((p\lfloor N/p\rfloor)!\) is a \(p\)-adic unit, and deleting the last
multiple of \(p\) lowers the valuation, proving the first equality.  Also
\[
 d_p\leq N+p\sum_{j\geq1}\frac{N}{p^j}\leq3N.
\]
For \(p=2\), Legendre's formula
\(v_2(N!)=N-s_2(N)\), together with
\(s_2(N)\leq\lfloor\log_2N\rfloor+1\), gives the lower bound in
\eqref{eq:factorial-degree-bound}.  Finally apply
\eqref{eq:touchard-remainder} and
Theorem~\ref{thm:bell-gould-arbitrary-modulus} with \(f=x^n\).
\end{proof}

The congruence tradition begins with Touchard's classical congruence for
Bell numbers and its polynomial variants
\cite{Touchard1933,SunZagier2011}; the adjoint language belongs to the
classical umbral calculus \cite{Roman1984}, and the companion Gould
recurrence was studied systematically by Gould and Quaintance
\cite{GouldQuaintance2007}.  The preceding proof records a particular
boundary-aware simultaneous lift and its arbitrary-modulus CRT packaging.
No claim of historical priority for that exact packaging is made without a
separate search through the older congruence and umbral literature.

\subsection{Matched mixed moments}

Retain the pure exponential field
\[
 \Eexp=\Qbar(e^\alpha:\alpha\in\Qbar).
\]
For \(U(t)=\sum_ku_kt^k\in\Qbar[t]\), define
\begin{equation}\label{eq:matched-beta-r}
 \beta(U)=\sum_ku_kk!,
 \qquad
 R(U)=\sum_ku_kk!H_k.
\end{equation}
For \(\lambda\in\Qbar_{>0}\) and
\(V(t)=\sum_kv_kt^k\in\Qbar[t]\), put
\begin{equation}\label{eq:matched-s}
 S_\lambda(V)=
 \sum_{k\geq1}v_k\sum_{j=0}^{k-1}
 (-\lambda)^j(k-1-j)!.
\end{equation}

\begin{theorem}[Matched mixed-integral independence]
\label{thm:matched-mixed-integrals}
Let \(\lambda_1,\ldots,\lambda_s\) be distinct positive algebraic numbers.
For each \(i\), choose \(U_i,V_i\in\Qbar[t]\) such that
\begin{equation}\label{eq:matched-condition}
 V_i(-\lambda_i)=-\beta(U_i)\ne0,
\end{equation}
and define
\begin{equation}\label{eq:matched-integral}
 \mathscr I_i=
 e^{\lambda_i}\int_0^\infty e^{-t}U_i(t)
       \log\!\left(\frac{t}{\lambda_i}\right)dt
 +\int_0^\infty\frac{e^{-t}V_i(t)}{t+\lambda_i}\,dt.
\end{equation}
Then \(\mathscr I_1,\ldots,\mathscr I_s\) are algebraically independent
over \(\Eexp\).
\end{theorem}

\begin{proof}
The logarithmic moments are
\begin{equation}\label{eq:matched-log-moment}
 \int_0^\infty e^{-t}t^k\log(t/\lambda)\,dt
 =k!(H_k-\gamma-\log\lambda).
\end{equation}
If
\(J_k(\lambda)=\int_0^\infty e^{-t}t^k/(t+\lambda)\,dt\), then
\[
 J_0(\lambda)=e^\lambda E_1(\lambda),
 \qquad
 J_k(\lambda)=(k-1)!-\lambda J_{k-1}(\lambda).
\]
Therefore
\begin{equation}\label{eq:matched-rational-moment}
 \int_0^\infty\frac{e^{-t}V(t)}{t+\lambda}\,dt
 =S_\lambda(V)+V(-\lambda)e^\lambda E_1(\lambda).
\end{equation}
Equations \eqref{eq:matched-condition}--\eqref{eq:matched-rational-moment}
and
\(\Ein(\lambda)=\gamma+\log\lambda+E_1(\lambda)\) give the exact identity
\begin{equation}\label{eq:matched-affine-ein}
 \boxed{
 \mathscr I_i=e^{\lambda_i}R(U_i)+S_{\lambda_i}(V_i)
 -e^{\lambda_i}\beta(U_i)\Ein(\lambda_i).}
\end{equation}
The constant term and the nonzero coefficient of \(\Ein(\lambda_i)\) lie
in \(\Eexp\).  Thus the displayed values are obtained from the distinct
\(\Ein(\lambda_i)\) by an invertible diagonal affine transformation over
\(\Eexp\).  Theorem~\ref{thm:pure-exp-base} proves the claim.
\end{proof}

There is also an exact relation-ideal version.  For a fixed \(\lambda_i\),
let \((U_{ij},V_{ij})\), \(1\leq j\leq r_i\), satisfy
\eqref{eq:matched-condition}, and put
\[
 a_{ij}=e^{\lambda_i}\beta(U_{ij}),
 \qquad
 c_{ij}=e^{\lambda_i}R(U_{ij})+S_{\lambda_i}(V_{ij}).
\]
If \(T_{ij}\) is the coordinate corresponding to
\(\mathscr I_{ij}\), then the complete relation ideal over \(\Eexp\) is
generated by
\begin{equation}\label{eq:matched-full-ideal}
 a_{i1}(T_{ij}-c_{ij})
 -a_{ij}(T_{i1}-c_{i1}),
 \qquad
 1\leq i\leq s,\quad 2\leq j\leq r_i.
\end{equation}
Indeed, \(\mathscr I_{ij}-c_{ij}=-a_{ij}\Ein(\lambda_i)\), and the
\(s\) remaining \(\Ein\)-coordinates are algebraically independent.

The smallest specialization already gives a countable independent family.
Taking \(U=1\) and \(V=-1\) yields
\begin{equation}\label{eq:minimal-matched-integral}
 \boxed{
 \mathscr I(\lambda)
 =e^\lambda\int_0^\infty e^{-t}\log(t/\lambda)\,dt
 -\int_0^\infty\frac{e^{-t}}{t+\lambda}\,dt
 =-e^\lambda\Ein(\lambda).}
\end{equation}
Hence
\(\{\mathscr I(\lambda):\lambda\in\Qbar_{>0}\}\) is algebraically
independent over \(\Eexp\), in the finite-subset sense.

\begin{corollary}[Filtered Touchard integrals]
\label{cor:filtered-touchard-transcendence}
For every \(n\geq0\),
\begin{equation}\label{eq:paired-touchard}
 \mathscr R_n-e\mathsf B_n
 \int_0^\infty e^{-t}\log t\,dt
 =e\mathsf B_n\Ein(1)-\mathsf G_n
\end{equation}
is transcendental over \(\Eexp\).

More generally, let \(P\in\Z[X]\), and put
\[
 q_n=P(E)\mathsf B_n,
 \qquad p_n=P(E)\mathsf G_n.
\]
If \(q_n\ne0\), then
\begin{equation}\label{eq:paired-filtered-touchard}
 \int_0^\infty
 \frac{e^{-x}[P(E)\tau_n](x)}{1+x}\,dx
 -eq_n\int_0^\infty e^{-t}\log t\,dt
 =q_ne\Ein(1)-p_n
\end{equation}
is transcendental over \(\Eexp\).
If \(P\) is monic, then \(q_n>0\) and \(q_n\to\infty\) for all
sufficiently large \(n\).
\end{corollary}

\begin{proof}
The identities follow from \eqref{eq:touchard-remainder} and
\(\int_0^\infty e^{-t}\log t\,dt=-\gamma\).  Their right sides are
nonconstant affine expressions in \(\Ein(1)\) over \(\Eexp\), so they are
transcendental.

For the final assertion, Dobi\'nski's formula gives
\[
 q_n=e^{-1}\sum_{k\geq0}\frac{k^nP(k)}{k!}.
\]
A monic real polynomial is positive at all sufficiently large nonnegative
integers.  Its negative terms in the displayed sum therefore have finite
support.  Choose a positive integer \(L\) beyond that support with
\(P(L)>0\).  The single positive term at \(L\) eventually dominates the
finite negative sum, whose bases are all smaller than \(L\).  It follows
that \(q_n>0\) eventually; choosing \(L\geq2\) also gives \(q_n\to\infty\).
\end{proof}

In particular, for the polynomial \(C_M\) of
Theorem~\ref{thm:bell-gould-arbitrary-modulus}, the two integers \(q_n,p_n\)
are divisible by \(M\).  Dividing
\eqref{eq:paired-filtered-touchard} by \(M\) therefore gives, for all
sufficiently large \(n\), an explicit normalized mixed integral which is
transcendental over \(\Eexp\) and whose \(\Ein(1)\)-coefficient tends to
infinity.  This closes nonvanishing and coefficient growth in the matched
\(e\gamma+\delta=e\Ein(1)\) direction.  It does not prove nonvanishing or
smallness for the unpaired \(\delta\)-linear form and therefore gives no
irrationality or transcendence theorem for \(\delta\) itself.


\section{Hankel determinants and critical capacity}

Let
\[
 D_n(x)=\det(H_{i+j}-x)_{0\le i,j\le n}.
\]

\begin{theorem}\label{thm:hankel}
For every \(n\ge0\) and \(x\in\C\),
\[
 D_n(x)=(-1)^n\frac{2H_n-x}
 {n!\displaystyle\prod_{j=1}^n\binom{2j}{j}\binom{2j-1}{j}}.
\]
For fixed \(x\), \( |D_n(x)|^{1/n^2}\to1/4\).
\end{theorem}

\begin{proof}
Starting from the bottom, replace row \(i\) by row \(i\) minus row \(i-1\)
for \(1\le i\le n\).  The bottom \(n\) rows become
\[
 \left(\frac1{i+j}\right)_{1\le i\le n,\ 0\le j\le n}.
\]
Its null vector is the coefficient vector \((c_0,\ldots,c_n)\) of the
shifted Legendre polynomial
\[
 P_n(2t-1)=\sum_{j=0}^nc_jt^j,
\]
because \(\int_0^1t^{i-1}P_n(2t-1)\,dt=0\) for \(1\le i\le n\).
The normalizations give
\[
 \sum_{j=0}^nc_j=1,\qquad
 \sum_{j=0}^nc_jH_j=2H_n.
\]
Laplace expansion along the first row is consequently a constant multiple of
\(2H_n-x\).  The cofactor of its last entry is the \(n\)-by-\(n\) Hilbert
determinant, while \(c_n=\binom{2n}{n}\).  The Cauchy determinant formula
then gives the displayed constant
\[
 \frac{(-1)^n}{n!\prod_{j=1}^n\binom{2j}{j}\binom{2j-1}{j}}.
\]
Finally Stirling's formula yields
\[
 \log\left(n!\prod_{j=1}^n\binom{2j}{j}\binom{2j-1}{j}\right)
 = (\log4)n^2+O(n\log n).
\]
Since \(2H_n-x\) has subexponential size and is nonzero for all sufficiently
large \(n\), the limit follows.
\end{proof}

The generating function is
\begin{equation}\label{eq:harmonic-gen}
 B_x(z)=\sum_{n\ge0}(H_n-x)z^n
 =\frac{-x-\log(1-z)}{1-z}.
\end{equation}
After the moment transformation implicit in the row differences, the
archimedean singular compact is \([0,1]\), whose logarithmic capacity is
\(1/4\).  Thus the determinant reaches, rather than beats, the extremal
capacity scale.  The parameter \(x\) enters the Hankel matrix through the
rank-one perturbation \(-x\mathbf1\mathbf1^{\mathsf T}\) and changes only
the \(O(n\log n)\) part.  Consequently this particular
P\'olya--Bertrandias/Hankel normalization has no strict inequality left with
which to distinguish \(x=\gamma\).  This is a statement about this
normalization, not a no-go theorem for every possible capacity construction.

\section{Positive mixed Hankel determinants and finite arithmetic reconstruction}
\label{sec:positive-mixed-hankel}

Throughout this section put $K=\Eexp$.

Put $H_0=0$, $H_m=\sum_{j=1}^m1/j$, and define
\begin{align}
 G_m
 &=\int_0^\infty x^me^{-x}(x-\log x)\,dx
   =m!(\gamma+m+1-H_m),
 \label{eq:mixed-G-moment}\\
 I_{k+1}
 &=\int_0^1t^k\frac{e^{1-1/t}}t\,dt
   =\int_0^\infty\frac{e^{-x}}{(1+x)^{k+1}}\,dx.
 \label{eq:mixed-I-moment}
\end{align}
Integration by parts gives $I_{m+1}=(1-I_m)/m$.  Hence
\begin{equation}\label{eq:mixed-I-arithmetic}
 k!I_{k+1}=A_k+(-1)^k\delta,
 \qquad A_0=0,
 \qquad A_k=(k-1)!-A_{k-1}\in\mathbb Z.
\end{equation}
For $d\geq1$ set
\begin{equation}\label{eq:mixed-Hankel-definition}
 M_k^{(d)}=G_{k+d-1}I_{k+1},
 \qquad
 \mathcal P_d(\gamma,\delta)
 =\det\bigl(M_{i+j}^{(d)}\bigr)_{0\leq i,j<d}.
\end{equation}

\begin{theorem}[Positive full-degree mixed determinant]
\label{thm:positive-full-degree-mixed-determinant}
For every $d\geq1$,
\[
 \mathcal P_d(\gamma,\delta)\in\mathbb Q[\gamma,\delta],
 \qquad \mathcal P_d(\gamma,\delta)>0.
\]
Its total degree is exactly $2d$, and its top mixed coefficient is
\begin{equation}\label{eq:mixed-Hankel-top-coefficient}
 [\gamma^d\delta^d]\mathcal P_d
 =(-1)^{d(d-1)/2}\bigl((d-1)!\bigr)^d\ne0.
\end{equation}
More precisely, it has the Heine representation
\begin{align}
 \mathcal P_d
 &=\frac1{d!}
 \int_{((0,\infty)\times(0,1))^d}
 \prod_{1\leq i<j\leq d}(x_it_i-x_jt_j)^2
 \notag\\[-1mm]
 &\qquad\times
 \prod_{\ell=1}^d
 x_\ell^{d-1}e^{-x_\ell}(x_\ell-\log x_\ell)
 \frac{e^{1-1/t_\ell}}{t_\ell}
 \,dx_\ell\,dt_\ell.
 \label{eq:mixed-Hankel-Heine}
\end{align}
\end{theorem}

\begin{proof}
The product variable $y=xt$ under the positive measure
\[
 x^{d-1}e^{-x}(x-\log x)
 \frac{e^{1-1/t}}t\,dx\,dt
\]
on $(0,\infty)\times(0,1)$ has kth moment $M_k^{(d)}$.
Andreief's identity gives \eqref{eq:mixed-Hankel-Heine}.  Since
$x-\log x\geq1$ and the product measure is non-atomic with infinite
support, the Vandermonde is nonzero almost everywhere off a null set;
the integral is therefore strictly positive.

Equations \eqref{eq:mixed-G-moment} and
\eqref{eq:mixed-I-arithmetic} prove polynomial membership and the upper
bound $2d$ for the total degree.  The coefficient of $\gamma\delta$ in
$M_k^{(d)}$ is
\[
 (-1)^k\frac{(k+d-1)!}{k!}
 =(-1)^k(k+1)_{d-1},
\]
where $(x)_n$ is the rising factorial.  Consequently the coefficient of
$\gamma^d\delta^d$ in the determinant equals
\[
 \det\bigl((-1)^{i+j}(i+j+1)_{d-1}\bigr)_{0\leq i,j<d}.
\]
Factoring the row and column signs has no net effect.  For
$p(k)=(k+1)_{d-1}$, apply the unitriangular forward-difference transform
to both rows and columns.  The resulting matrix has entries
$\mathsf D^{i+j}p(0)$; it vanishes past the anti-diagonal, while every
anti-diagonal entry is
$\mathsf D^{d-1}p(0)=(d-1)!$.  Its determinant is therefore the right
side of \eqref{eq:mixed-Hankel-top-coefficient}.  This coefficient is
nonzero, so the total degree is exactly $2d$.
\end{proof}

\begin{proposition}[Finite reconstruction from the first two mixed determinants]
\label{prop:mixed-hankel-reconstruction}
Put $U=\mathcal P_1$ and $V=\mathcal P_2$ for the determinants in
Theorem~\ref{thm:positive-full-degree-mixed-determinant}.  Then the extension
\[
 K(\gamma,\delta)/K(U,V)
\]
is algebraic of degree at most four.  More explicitly,
\[
 U=(\gamma+1)\delta,
\]
and $\delta$ is a root of
\begin{equation}\label{eq:mixed-hankel-quartic}
 2T^4+(U-4)T^3
 +(2U^2-16U+2V+2)T^2
 +(-16U^2+8U)T+8U^2=0,
\end{equation}
after which $\gamma=U/\delta-1$.
Consequently
\[
 \operatorname{trdeg}_K K(U,V)
 =\operatorname{trdeg}_K K(\gamma,\delta).
\]
In particular, at least one of $\mathcal P_1,\mathcal P_2$ is
transcendental over $K$, and the explicit number
\[
 \mathcal P_1+i\mathcal P_2
\]
is transcendental over $K$.  Moreover $\mathcal P_1,\mathcal P_2$ are
algebraically independent over $K$ if and only if $\gamma,\delta$ are.
\end{proposition}

\begin{proof}
The formulas for $G_m$ and $I_m$ give
$\mathcal P_1=(\gamma+1)\delta$.  The explicit second determinant is
\[
 2\mathcal P_2=-18+36\delta-5\delta^2-24\gamma
 +48\gamma\delta-5\gamma\delta^2-8\gamma^2
 +16\gamma^2\delta-2\gamma^2\delta^2.
\]
Substituting $\gamma=U/\delta-1$ and clearing $\delta^2$ gives
\eqref{eq:mixed-hankel-quartic}.  Since $\delta>0$, the displayed
recovery of $\gamma$ is valid.  This proves finiteness and equality of
transcendence degrees over $K$.

The pure exponential-base theorem makes
$\eta=\Ein(1)=\gamma+\delta/e$ transcendental over $K$.  Since $e\in K$,
the two numbers $\gamma,\delta$ cannot both be algebraic over $K$.
Therefore $U,V$ cannot both be algebraic over $K$.  The field $K$ is
stable under complex conjugation.  If $U+iV$ were algebraic over $K$,
then its conjugate would be algebraic over $K$, and hence so would its
real and imaginary parts $U,V$, a contradiction.  The final equivalence
is the case of transcendence degree two.
\end{proof}


% Standalone insert for the revised gamma paper.
% Expected packages: amsmath, amssymb, amsthm, mathtools.
% Suggested bibliography keys are listed at the end of this file.

\section{Geronimus--Markov reduction and the Laguerre--Kummer tower}
\label{sec:geronimus-kummer}

This section replaces finite symbolic verification by an all-degree argument.
It also separates three different statements that should not be conflated:
the exact determinant identity, the classical Painlev\'e specialization, and
the arithmetic usefulness of the resulting linear forms.

Let \(d\nu\) be a positive measure of infinite support on
\([0,\infty)\), with finite moments
\[
  m_k=\int_0^\infty x^k\,d\nu(x) \qquad(k\geq 0).
\]
Write \(p_j\) for its monic orthogonal polynomials and
\[
  \Delta_j=\det(m_{r+c})_{0\leq r,c<j},\qquad \Delta_0=1.
\]
For \(s>0\), set
\[
  \mu_k(s)=\int_0^\infty\frac{x^k}{x+s}\,d\nu(x),
  \qquad
  \mathcal D_n(s)=\det(\mu_{r+c}(s))_{0\leq r,c<n}.
\]

\begin{theorem}[Geronimus--Markov determinant reduction]
\label{thm:GM-reduction}
For every \(n\geq1\),
\begin{equation}
 \boxed{
 \mathcal D_n(s)=(-1)^{n-1}\Delta_{n-1}
 \int_0^\infty\frac{p_{n-1}(x)}{x+s}\,d\nu(x).}
 \label{eq:GM-reduction}
\end{equation}
If \(F(s)=\mu_0(s)\), then equivalently
\begin{equation}
 \begin{aligned}
 \mathcal D_n(s)
   &=\Delta_{n-1}\bigl(Q_{n-1}(s)F(s)-P_{n-2}(s)\bigr),\\
 Q_{n-1}(s)&=(-1)^{n-1}p_{n-1}(-s)>0,
 \end{aligned}
 \label{eq:Markov-Pade}
\end{equation}
where \(P_{n-2}\) is a polynomial of degree at most \(n-2\), with the
convention \(P_{-1}=0\).
Thus the Cauchy-transform Hankel determinant is precisely the classical
Markov--Pad\'e remainder multiplied by the preceding moment determinant.
\end{theorem}

\begin{proof}
In the defining matrix for \(\mathcal D_n\), perform the column operations
\[
 C_j\longleftarrow C_j+sC_{j-1},
 \qquad j=n-1,n-2,\ldots,1,
\]
in this order.  Since
\(\mu_k+s\mu_{k-1}=m_{k-1}\) for \(k\geq1\), the resulting columns are
\[
  (\mu_i)_{i=0}^{n-1},
  (m_i)_{i=0}^{n-1},
  \ldots,
  (m_{i+n-2})_{i=0}^{n-1}.
\]
Expansion in the first column is the functional
\(f\mapsto\int f(x)(x+s)^{-1}d\nu(x)\) applied to a polynomial orthogonal,
with respect to \(d\nu\), to \(1,x,\ldots,x^{n-2}\).  Its leading coefficient
is the cofactor \((-1)^{n-1}\Delta_{n-1}\).  The polynomial is therefore
\((-1)^{n-1}\Delta_{n-1}p_{n-1}\), proving
\eqref{eq:GM-reduction}.

Now write
\[
 \frac{p_{n-1}(x)}{x+s}
 =\frac{p_{n-1}(-s)}{x+s}
  +\frac{p_{n-1}(x)-p_{n-1}(-s)}{x+s}.
\]
The second quotient is a polynomial of degree at most \(n-2\), which gives
\eqref{eq:Markov-Pade}.  Finally, all zeros of \(p_{n-1}\) lie in
\((0,\infty)\), so \((-1)^{n-1}p_{n-1}(-s)>0\).
\end{proof}

\begin{remark}[Necessary scope]
For an arbitrary signed or finitely supported weight, the rank-one argument
only proves degree at most one in \(F\).  Exact degree one follows under the
positive infinite-support assumptions above.  Moreover,
\eqref{eq:Markov-Pade} is a structural Pad\'e identity; it does not by itself
compare arithmetic heights.  Any assertion that the determinant is
arithmetically weaker than Pad\'e requires a separate estimate for
\(\Delta_{n-1}\).
\end{remark}

We now specialize to the Laguerre measure.  For \(a>-1\), put
\begin{align*}
 d\nu_a(x)&=x^ae^{-x}\,dx,\\
 D_n^{(a)}(s)&=
 \det\left(\int_0^\infty
 \frac{x^{i+j+a}e^{-x}}{x+s}\,dx\right)_{0\leq i,j<n},\\
 C_n(a)&=\prod_{k=0}^{n-1}k!\,\Gamma(k+a+1)
 =\frac{G(n+1)G(n+a+1)}{G(a+1)}.
\end{align*}
Here and below \(G\) denotes the Barnes \(G\)-function.

\begin{theorem}[Exact Laguerre--Kummer collapse]
\label{thm:Laguerre-Kummer-exact}
For every \(n\geq1\), \(a>-1\), and \(s>0\),
\begin{equation}
 \boxed{D_n^{(a)}(s)=C_n(a)U(n,1-a,s).}
 \label{eq:Laguerre-Kummer-exact}
\end{equation}
Consequently \(y=D_n^{(a)}\) satisfies
\begin{equation}
 sy''+(1-a-s)y'-ny=0.
 \label{eq:Kummer-Dn}
\end{equation}
If
\(\Sigma_n=s(\log D_n^{(a)})'\), then
\begin{equation}
 s\Sigma_n'=(s+a)\Sigma_n-\Sigma_n^2+ns.
 \label{eq:Riccati-Sigma}
\end{equation}
In addition, if \(u_n=U(n,1-a,s)\) and \(u_0=1\), then
\begin{equation}
 n(n+a)u_{n+1}=(s+2n+a-1)u_n-u_{n-1}
 \qquad(n\geq1).
 \label{eq:Kummer-discrete}
\end{equation}
\end{theorem}

\begin{proof}
The monic Laguerre polynomial is
\(p_{n-1}(x)=(-1)^{n-1}(n-1)!L_{n-1}^{(a)}(x)\), and
\(\Delta_{n-1}=C_{n-1}(a)\).  The Laplace representation of
\((x+s)^{-1}\) and the Laguerre transform give
\begin{equation}
 \int_0^\infty x^ae^{-(1+t)x}L_{n-1}^{(a)}(x)\,dx
 =\frac{\Gamma(n+a)}{(n-1)!}
   \frac{t^{n-1}}{(1+t)^{n+a}}.
 \label{eq:Laguerre-Laplace}
\end{equation}
Substitution in \eqref{eq:GM-reduction}, followed by
\begin{equation}
 U(n,1-a,s)=\frac1{\Gamma(n)}\int_0^\infty
 e^{-st}\frac{t^{n-1}}{(1+t)^{n+a}}\,dt,
 \label{eq:U-integral}
\end{equation}
yields \eqref{eq:Laguerre-Kummer-exact}, because
\[
 C_{n-1}(a)\Gamma(n)\Gamma(n+a)=C_n(a).
\]
Equations \eqref{eq:Kummer-Dn} and \eqref{eq:Riccati-Sigma} follow from
Kummer's equation.  Equation \eqref{eq:Kummer-discrete} is the monic
Laguerre three-term recurrence for its function of the second kind.  It can
also be checked directly from \eqref{eq:U-integral}: integrate the derivative
of
\[
 e^{-st}t^{n-1}(1+t)^{-n-a}.
\]
For \(n\geq2\) both endpoint terms vanish and elementary rearrangement gives
\eqref{eq:Kummer-discrete}; for \(n=1\), the endpoint at zero is \(-1\), which
is precisely the term \(-u_0\).
\end{proof}

\begin{remark}[Literature status]
The mathematical content of \eqref{eq:Laguerre-Kummer-exact} is classical:
the Stieltjes transform of a monic Laguerre polynomial is its function of the
second kind and is a Kummer \(U\)-function.  See, for example,
Dea\~no--Huertas--Rom\'an~\cite{DeanoHuertasRoman2018}, especially their
equations~(4)--(5), and DLMF~\cite{DLMF}, \S13.4.  The point here is to provide
a complete determinant proof and to fix the normalization, not to claim a new
special-function identity.
\end{remark}

\subsection{Rigorous growth and the classical Pad\'e remainder}

\begin{theorem}[Superfactorial divergence]
\label{thm:Laguerre-divergence}
For fixed \(a>-1\) and \(s>0\),
\begin{equation}
 \frac{e^{-2s}}{3^{a+1}2^{n-1}\Gamma(n)}
 \leq U(n,1-a,s)\leq\frac1{s\Gamma(n)}.
 \label{eq:U-bounds}
\end{equation}
Consequently
\begin{equation}
 \boxed{
 \log D_n^{(a)}(s)=n^2\log n-\frac32n^2+O(n\log n),}
 \qquad
 D_n^{(a)}(s)\longrightarrow+\infty.
 \label{eq:Dn-growth}
\end{equation}
\end{theorem}

\begin{proof}
For the lower bound in \eqref{eq:U-bounds}, restrict
\eqref{eq:U-integral} to \(1\leq t\leq2\) and write its integrand as
\[
 e^{-st}\left(\frac{t}{1+t}\right)^{n-1}(1+t)^{-a-1}.
\]
On this interval the three factors are bounded below by
\(e^{-2s}\), \(2^{-(n-1)}\), and \(3^{-a-1}\), respectively.  For the upper
bound, the last two factors have product at most one, and
\(\int_0^\infty e^{-st}dt=s^{-1}\).
Using
\[
 C_n(a)=\frac{G(n+1)G(n+a+1)}{G(a+1)}
\]
and the standard Barnes-\(G\) asymptotic now gives
\eqref{eq:Dn-growth}; see DLMF~\cite[\S5.17]{DLMF}.
\end{proof}

For \(a=0\), let \(B(s)=e^sE_1(s)=U(1,1,s)\), so that
\(B(1)=\delta\), and put
\[
 Q_{n-1}(s)=(n-1)!L_{n-1}(-s).
\]
Theorem~\ref{thm:GM-reduction} gives a polynomial
\(P_{n-2}(s)\in\mathbb Z[s]\) such that
\begin{equation}
 \boxed{
 Q_{n-1}(s)B(s)-P_{n-2}(s)=(n-1)!^2U(n,1,s).}
 \label{eq:classical-Pade-remainder}
\end{equation}
Indeed, \((n-1)!L_{n-1}(x)\in\mathbb Z[x]\), polynomial division by
\(x+s\) stays in \(\mathbb Z[x,s]\), and the moments of \(e^{-x}dx\) are
factorials.  At \(s=1\), \eqref{eq:classical-Pade-remainder} is the classical
Laguerre--Stieltjes Pad\'e remainder for the Euler--Gompertz constant.  In
particular,
\begin{equation}
 Q_{n-1}(1)\delta-P_{n-2}(1)
 \geq \frac{e^{-2}}{3}\frac{(n-1)!}{2^{n-1}}
 \longrightarrow+\infty.
 \label{eq:Pade-divergence}
\end{equation}
Thus both the determinant and its underlying classical Pad\'e linear form
diverge; this is a theorem, not numerical evidence.  The continued-fraction
identification is classical; see
Lagarias~\cite[\S2.5]{Lagarias2013} and
Ernvall-Hyt\"onen--Matala-aho--Sepp\"al\"a~\cite{EMHS2023}.

\subsection{Hausdorff moments, total positivity, and random matrices}

\begin{proposition}[Complete monotonicity and Tur\'an inequality]
\label{prop:U-Hausdorff}
Define
\[
 h_m(a,s)=m!U(m+1,1-a,s)\qquad(m\geq0).
\]
Then
\begin{equation}
 \boxed{
 h_m(a,s)=\int_0^1 r^m(1-r)^{a-1}
 \exp\left(-\frac{sr}{1-r}\right)dr.}
 \label{eq:U-Hausdorff}
\end{equation}
Consequently \((h_m)_{m\geq0}\) is a strict Hausdorff moment sequence.  With
\(\Delta f_m=f_{m+1}-f_m\),
\begin{equation}
 (-1)^k\Delta^kh_m>0\qquad(k,m\geq0).
 \label{eq:complete-monotone}
\end{equation}
It is strictly log-convex, and therefore
\begin{equation}
0<n\frac{U(n+1,1-a,s)}{U(n,1-a,s)}<1,
\qquad
n\frac{U(n+1,1-a,s)}{U(n,1-a,s)}\nearrow1,
\qquad(n\geq1).
\label{eq:U-ratio}
\end{equation}
and
\begin{equation}
 \frac{n}{n+1}U(n+1,1-a,s)^2
 <U(n,1-a,s)U(n+2,1-a,s)
 \qquad(n\geq1).
 \label{eq:U-Turan}
\end{equation}
\end{proposition}

\begin{proof}
In \eqref{eq:U-integral}, set \(r=t/(1+t)\) to obtain
\eqref{eq:U-Hausdorff}.  Formula \eqref{eq:complete-monotone} follows by
inserting the factor \((1-r)^k\) under the integral.  Strict log-convexity is
Cauchy--Schwarz, since the representing measure is not a point mass.  The
successive moment ratios are strictly increasing and tend to the essential
supremum of the support, which is one; this gives \eqref{eq:U-ratio}.
Rewriting strict log-convexity for \(h_n=n!U(n+1,1-a,s)\) gives
\eqref{eq:U-Turan}.
\end{proof}

\begin{proposition}[Total positivity and determinacy]
\label{prop:Laguerre-total-positive}
For \(a>-1\) and \(s>0\), the infinite Hankel matrix
\[
 \left(\int_0^\infty\frac{x^{i+j+a}e^{-x}}{x+s}\,dx\right)_{i,j\geq0}
\]
is strictly totally positive.  Its Stieltjes moment problem is determinate.
\end{proposition}

\begin{proof}
Andreief's identity expresses a minor with increasing row exponents
\(r_1<\cdots<r_k\) and column exponents \(c_1<\cdots<c_k\) as the integral of
\[
 \det(x_j^{r_i})_{i,j=1}^k\det(x_j^{c_i})_{i,j=1}^k
 \prod_{j=1}^k\frac{x_j^ae^{-x_j}}{x_j+s}\,dx_j.
\]
On \(0<x_1<\cdots<x_k\), each alternant is a Vandermonde determinant times
a Schur polynomial and is strictly positive.  Hence every minor is positive.
For determinacy, if \(\mu_k\) denotes the moment sequence, then
\[
 \mu_{2k}\leq s^{-1}\Gamma(2k+a+1),
\]
so \(\sum_k\mu_{2k}^{-1/(2k)}=\infty\); Carleman's criterion applies.
\end{proof}

\begin{corollary}[Laguerre-unitary-ensemble interpretation]
\label{cor:LUE-inverse-characteristic}
Let \(X\) have the \(n\)-dimensional Laguerre unitary ensemble with parameter
\(a>-1\), normalized so that its eigenvalue density is proportional to
\[
 \prod_{i<j}(x_i-x_j)^2\prod_{i=1}^n x_i^ae^{-x_i}.
\]
Then
\begin{equation}
 \boxed{\mathbb E\det(X+sI)^{-1}=U(n,1-a,s).}
 \label{eq:LUE-inverse-characteristic}
\end{equation}
\end{corollary}

\begin{proof}
Andreief's identity identifies the ratio of the pole-deformed and undeformed
partition functions with \(D_n^{(a)}(s)/C_n(a)\).  Apply
\eqref{eq:Laguerre-Kummer-exact}.
\end{proof}

Equation \eqref{eq:LUE-inverse-characteristic} is the standard average inverse
characteristic polynomial, or Laguerre function of the second kind.  It is a
useful interpretation, not a new random-matrix problem solved here.

\subsection{The exact Painlev\'e V specialization and its boundary condition}

The Riccati equation \eqref{eq:Riccati-Sigma} is not itself the standard
Jimbo--Miwa--Okamoto sigma form.  Put
\[
 H_n(s)=\Sigma_n(s)+n.
\]
We use the JMO convention~\cite[Appendix~C, Eq.~(C.45)]{JimboMiwa1981}
\[
 (sH'')^2=
 \bigl[H-sH'+2(H')^2+(\nu_0+\nu_1+\nu_2+\nu_3)H'\bigr]^2
 -4\prod_{j=0}^3(H'+\nu_j).
\]
Direct substitution gives
\begin{equation}
\begin{split}
 (sH_n'')^2={}&
 \bigl[H_n-sH_n'+2(H_n')^2-(2n+a+1)H_n'\bigr]^2\\
 &-4H_n'(H_n'-1)(H_n'-n)(H_n'-n-a).
\end{split}
\label{eq:JMO-PV}
\end{equation}
This is the JMO sigma form of \(P_V\) with
\[
 (\nu_0,\nu_1,\nu_2,\nu_3)=(0,-1,-n,-n-a).
\]
For \(a>0\), it is the \(\lambda=-1\) specialization of the deformed Laguerre
theorem of Mu--Lyu~\cite{MuLyu2025}; for \(-1<a\leq0\), the identification
follows directly from \eqref{eq:Riccati-Sigma}.  Equation
\eqref{eq:Riccati-Sigma} also shows that this specialization lies on a
classical Riccati locus.

For \(a=0\), the identity \(B'=B-s^{-1}\) is invariant under
\(B\mapsto B+ce^s\).  At the level of moments this is
\begin{equation}
 \mu_k(s)\longmapsto\mu_k(s)+ce^s(-s)^k,
 \label{eq:Uvarov-shift}
\end{equation}
namely the Uvarov mass modification
\[
 \frac{e^{-x}}{x+s}\,dx
 \longmapsto
 \frac{e^{-x}}{x+s}\,dx+ce^s\delta_{-s}.
\]
The transformed determinant is another linear combination of the two Kummer
solutions and remains on the same Riccati/Painlev\'e locus.  However, it does
not preserve the distinguished Stieltjes boundary when \(c\neq0\).  For every
\(a>-1\), that boundary is
\begin{equation}
 D_n^{(a)}(s)\sim C_n(a)s^{-n}\qquad(s\to+\infty),
 \label{eq:recessive-boundary}
\end{equation}
or equivalently
\[
 H_n(s)=\frac{n(n+a)}s+O(s^{-2}).
\]
If \(c<0\), the added atom has negative mass; if \(c\notin\mathbb R\), the
functional is not a positive real measure; and if \(c>0\), positive mass is
placed at \(-s\notin[0,\infty)\).  Thus no \(c\neq0\) preserves the Stieltjes
class under discussion.
Thus the differential equation without boundary data has the translation
symmetry, whereas the actual Stieltjes and random-matrix solution is
canonically selected.  No arithmetic ``no-go'' statement follows from the
formal shift alone.  The direct \(B\)-shift applies to \(a=0\), and after
reduction to \(B\) to nonnegative integral \(a\); it is not an action on the
arbitrary-real-\(a\) family.

\section{A confluent reciprocal-Gamma matching determinant}
\label{sec:reciprocal-gamma-matching}

Put \(H_0=0\), \(H_m=\sum_{r=1}^m r^{-1}\), and
\[
 K_N(x)=\det\left(\frac{H_{i+j}-x}{(i+j)!}\right)_{0\leq i,j\leq N},
 \qquad
 M_N=\binom{N+1}{2}.
\]
Let \(\mathfrak M_N\) be the set of partial matchings of
\(\{0,1,\ldots,N\}\), and write \(V(\mathfrak m)\) for the matched vertices.
Define
\begin{align*}
 h_{N,j}&=H_{N+j}-H_j+H_{N-j},\\
 u_{N,j}(x)&=x-h_{N,j},\\
 B_N&=\frac{\prod_{k=1}^Nk!}{\prod_{m=N}^{2N}m!}.
\end{align*}

\begin{theorem}[Reciprocal-Gamma matching formula]
\label{thm:reciprocal-gamma-matching}
For every \(N\geq1\),
\begin{equation}
 \boxed{
 K_N(x)=(-1)^{N+1+M_N}B_N Z_N(x),}
 \label{eq:KN-matching-prefactor}
\end{equation}
where
\begin{equation}
 Z_N(x)=
 \sum_{\mathfrak m\in\mathfrak M_N}
 \left(\prod_{\{j,k\}\in\mathfrak m}\frac1{(j-k)^2}\right)
 \left(\prod_{j\notin V(\mathfrak m)}u_{N,j}(x)\right).
 \label{eq:ZN-matching}
\end{equation}
In particular \(\deg K_N=N+1\), with leading coefficient
\((-1)^{N+1+M_N}B_N\).
\end{theorem}

\begin{proof}
Introduce one auxiliary variable per column and set
\[
 F_N(\boldsymbol\alpha)=
 \det\left(\frac1{\Gamma(i+j+1+\alpha_j)}\right)_{0\leq i,j\leq N}.
\]
A reciprocal-Gamma Vandermonde evaluation gives
\begin{equation}
 F_N(\boldsymbol\alpha)=(-1)^{M_N}
 \frac{\prod_{0\leq j<k\leq N}
 (k-j+\alpha_k-\alpha_j)}
 {\prod_{j=0}^N\Gamma(N+j+1+\alpha_j)}.
 \label{eq:reciprocal-gamma-Vandermonde}
\end{equation}
Since
\[
 -\left.(\partial_{\alpha_j}+x-\gamma)
 \frac1{\Gamma(i+j+1+\alpha_j)}\right|_{\alpha_j=0}
 =\frac{H_{i+j}-x}{(i+j)!},
\]
we apply one such operator in every column.  At
\(\boldsymbol\alpha=0\), one has
\(F_N(\boldsymbol0)=(-1)^{M_N}B_N\neq0\), so a local logarithm can be chosen.
Equivalently, put
\[
 \widetilde F_N(\boldsymbol\alpha)
 =e^{(x-\gamma)\sum_j\alpha_j}F_N(\boldsymbol\alpha).
\]
Then
\[
 \left.\prod_{j=0}^N(\partial_{\alpha_j}+x-\gamma)F_N
 \right|_{\boldsymbol0}
 =
 \left.\partial_{\alpha_0}\cdots\partial_{\alpha_N}\widetilde F_N
 \right|_{\boldsymbol0}.
\]
The distinct-variable logarithmic cumulants of \(\widetilde F_N\) are
\begin{align*}
 \kappa_{\{j\}}
 &=\left.\partial_{\alpha_j}\log F_N\right|_{\boldsymbol0}+x-\gamma
   =u_{N,j}(x),\\
 \kappa_{\{j,k\}}
 &=\left.\partial_{\alpha_j}\partial_{\alpha_k}\log F_N
   \right|_{\boldsymbol0}=(j-k)^{-2},\\
 \kappa_S
 &=\left.\prod_{j\in S}\partial_{\alpha_j}\log F_N
   \right|_{\boldsymbol0}=0\qquad(|S|\geq3).
\end{align*}
The multivariate moment--cumulant formula therefore sums precisely over set
partitions into singletons and pairs, which are the partial matchings in
\eqref{eq:ZN-matching}.  Evaluation of \(F_N(\boldsymbol0)\), together with
the \(N+1\) minus signs, gives
\eqref{eq:KN-matching-prefactor}.
\end{proof}

The matching sum has two exact cross-field interpretations.  First, it is the
monomer--dimer partition function of the complete graph on
\(\{0,\ldots,N\}\), with monomer activities \(u_{N,j}(x)\) and positive dimer
activities \((j-k)^{-2}\).  Second, let \(A_N\) be a random real
skew-symmetric matrix whose upper-triangular entries are independent and
centered, with
\[
 \mathbb E(A_N)_{jk}^2=(j-k)^{-2}\qquad(j<k),
\]
and let
\(\mathsf H_N=\operatorname{diag}(h_{N,0},\ldots,h_{N,N})\).  Expansion by
permutations gives
\begin{equation}
 \boxed{Z_N(x)=\mathbb E\det(xI-\mathsf H_N+A_N).}
 \label{eq:skew-RMT-matching}
\end{equation}
Indeed, every cycle of length at least three contains independent centered
entries to the first power and has zero expectation, while the transpositions
give the dimer factors.  Formula \eqref{eq:skew-RMT-matching} is an expected
characteristic-polynomial representation, not a random-matrix limit theorem.

\begin{theorem}[Heilmann--Lieb complex-zero strip]
\label{thm:KN-zero-strip}
Every complex zero \(\zeta\) of \(K_N\) satisfies
\begin{equation}
 \boxed{H_{2N}-H_N\leq\Re\zeta\leq2H_N.}
 \label{eq:KN-zero-strip}
\end{equation}
\end{theorem}

\begin{proof}
A direct difference calculation shows that \(j\mapsto h_{N,j}\) is strictly
decreasing, with
\[
 h_{N,0}=2H_N,
 \qquad
 h_{N,N}=H_{2N}-H_N.
\]
If \(\Re x>2H_N\), every monomer activity
\(u_{N,j}(x)=x-h_{N,j}\) lies in the open right half-plane.  If
\(\Re x<H_{2N}-H_N\), they all lie in the open left half-plane.  The
Heilmann--Lieb left--right half-plane theorem for a monomer--dimer partition
function with nonnegative dimer activities~\cite[Theorem~4.6 and
Lemma~4.7]{HeilmannLieb1972} says that \(Z_N(x)\neq0\) in either case; see also
the precise multivariate restatement in
\cite[Theorem~10.1]{ChoeOxleySokalWagner2004}.  Equation
\eqref{eq:KN-matching-prefactor} proves the claim.
\end{proof}

\begin{remark}
Theorem~\ref{thm:KN-zero-strip} does not assert real-rootedness.  In fact
\[
 K_1(x)=-\frac{2x^2-5x+4}{4}
\]
already has nonreal roots.  Heilmann--Lieb supplies a zero-free complement of
a vertical strip, not a real-rooted matching polynomial.
\end{remark}

\begin{proposition}[Cleared determinants below the strip]
\label{prop:KN-divergence}
Fix a real \(x<\log2\).  If \(q_N\) is any positive integer that clears all
coefficient denominators of \(K_N\), then, for all sufficiently large \(N\),
\begin{equation}
 q_N|K_N(x)|
 \geq\prod_{j=0}^N(h_{N,j}-x)
 \longrightarrow+\infty.
 \label{eq:KN-cleared-divergence}
\end{equation}
\end{proposition}

\begin{proof}
For large \(N\), \(x<h_{N,N}=H_{2N}-H_N\).  Hence all singleton factors
\(u_{N,j}(x)\) are negative.  A matching with \(r\) pairs leaves
\(N+1-2r\) singletons, so every term in \eqref{eq:ZN-matching} has the same
sign \((-1)^{N+1}\).  There is no cancellation, and the empty matching gives
\[
 |Z_N(x)|\geq\prod_{j=0}^N(h_{N,j}-x).
\]
The leading coefficient of \(K_N\) is \(\pm B_N\).  Since \(q_NK_N\) has
integer coefficients, \(q_NB_N\) is a positive integer and is at least one.
This proves the inequality.  Finally, for \(j\leq N-1\),
\[
 h_{N,j}\geq h_{N,N-1}
 =H_{2N-1}-H_{N-1}+1\longrightarrow1+\log2,
\]
whereas \(h_{N,N}-x\to\log2-x>0\).  At least \(N\) factors are therefore
eventually bounded below by a fixed number greater than one, proving
divergence.
\end{proof}

In particular, \eqref{eq:KN-cleared-divergence} applies at \(x=\gamma\): an
elementary estimate gives \(\gamma<\log2\).  For completeness, using
\[
 \gamma=\sum_{k\geq1}\left(\frac1k-\log\left(1+\frac1k\right)\right),
 \qquad
 u-\log(1+u)<\frac{u^2}{2},
\]
one obtains \(\gamma<11/8-\log2<\log2\); the last inequality follows from
\(\log2>2(1/3+1/81)=56/81>11/16\).  Hence
\[
 q_N|K_N(\gamma)|\longrightarrow+\infty.
\]
This is an obstruction for this cleared-determinant normalization, not an
irrationality or transcendence statement about \(\gamma\).

\begin{remark}[Priority and scope]
The Geronimus reduction and the Laguerre function of the second kind are
classical.  So are the continued fraction, the Uvarov modification, and the
Painlev\'e V embedding, or they are direct specializations of cited results.
The zero strip is a direct application of the classical Heilmann--Lieb theorem
to the exact matching formula.  These statements give rigorous structural
closures.  They also rule out the specified determinant normalizations as
sources of small linear forms.
They do not solve a named open problem in orthogonal polynomials, moment
theory, continued fractions, isomonodromy, or random matrices.  The exact
priority of the differentiated reciprocal-Gamma matching formula itself
requires a specialist literature search before any novelty claim.
\end{remark}

% Suggested bibliography entries/keys for the main manuscript:
% DeanoHuertasRoman2018:
% A. Deano, E. J. Huertas, P. Roman, Asymptotics of orthogonal polynomials
% generated by a Geronimus perturbation of the Laguerre measure,
% J. Math. Anal. Appl. 433 (2016), 732--746; arXiv:1504.05976.
%
% MuLyu2025:
% X. Mu, S. Lyu, Hankel determinants for a deformed Laguerre weight with
% multiple variables and generalized Painleve V equation,
% arXiv:2410.22780; J. Math. Phys. 66 (2025), 073506.
%
% JimboMiwa1981:
% M. Jimbo, T. Miwa, Monodromy preserving deformation of linear ordinary
% differential equations with rational coefficients. II,
% Physica D 2 (1981), 407--448.
%
% HeilmannLieb1972:
% O. J. Heilmann, E. H. Lieb, Theory of monomer-dimer systems,
% Comm. Math. Phys. 25 (1972), 190--232.
%
% ChoeOxleySokalWagner2004:
% Y.-B. Choe, J. G. Oxley, A. D. Sokal, D. G. Wagner,
% Homogeneous multivariate polynomials with the half-plane property,
% Adv. in Appl. Math. 32 (2004), 88--187.
%
% Lagarias2013:
% J. C. Lagarias, Euler's constant: Euler's work and modern developments,
% Bull. Amer. Math. Soc. 50 (2013), 527--628; arXiv:1303.1856.
%
% EMHS2023:
% A.-M. Ernvall-Hytonen, T. Matala-aho, L. Seppala,
% Euler's factorial series, Hardy integral, and continued fractions,
% J. Number Theory 244 (2023), 224--250; arXiv:2111.13649.
%
% DLMF:
% NIST Digital Library of Mathematical Functions, Chapters 5, 13, and 18.


\begin{proposition}[Exact adelic denominator wall for the one-pole lift]
\label{prop:borel-laguerre-adelic-wall}
Let $M\geq N+1$, put $m=M-N-1$, and let
\[
 K_{N,M}(z)=\frac{z^N}{N!}\,{}_1F_1(M;N+1;-z),
 \qquad
 \mathcal H[K_{N,M}](y)=\frac{y^N}{(1+y)^M}.
\]
Then
\begin{equation}\label{eq:borel-laguerre-clearing-expansion}
 e^zK_{N,M}(z)
 =\sum_{j=0}^{m}(-1)^j\binom mj
   \frac{z^{N+j}}{(N+j)!}.
\end{equation}
The least positive common denominator of this polynomial is exactly
$(M-1)!$.  Equivalently, at every prime $p$ its local denominator
exponent is exactly
\begin{equation}\label{eq:borel-laguerre-local-wall}
 \max_{0\leq j\leq m}
 \max\!\left\{0,-v_p\!\left(
 \frac{\binom mj}{(N+j)!}\right)\right\}
 =v_p((M-1)!).
\end{equation}
Consequently
\begin{equation}\label{eq:borel-laguerre-primitive-polynomial}
 \widetilde K_{N,M}(z):=(M-1)!e^zK_{N,M}(z)
 \in\mathbb Z[z]
\end{equation}
is primitive.  Its coefficient height satisfies
\begin{equation}\label{eq:borel-laguerre-height-bounds}
 \frac{(M-1)!}{N!}
 \leq H(\widetilde K_{N,M})
 \leq 2^m\frac{(M-1)!}{N!}.
\end{equation}
If $M=\alpha N$ with fixed $\alpha>1$ along an integral subsequence, then
\begin{align}
 \log H(\widetilde K_{N,M})
 &=(\alpha-1)N\log N+O_\alpha(N),
 \label{eq:borel-laguerre-height-asymptotic}\\
 \log\!\left((M-1)!
 \max_{y\geq0}\frac{y^N}{(1+y)^M}\right)
 &=\alpha N\log N+O_\alpha(N).
 \label{eq:borel-laguerre-normalized-bump}
\end{align}
In fact the restriction to a fixed ratio is unnecessary.  For fixed
$N$, the quantity
\[
 \mathfrak F_{N,M}:=(M-1)!
 \max_{y\geq0}\frac{y^N}{(1+y)^M}
\]
is strictly increasing for integers $M\geq N+1$.  Hence
\begin{equation}\label{eq:borel-laguerre-uniform-wall}
 \mathfrak F_{N,M}
 \geq \mathfrak F_{N,N+1}
 =N!\frac{N^N}{(N+1)^{N+1}}
 =\exp\bigl(N\log N-N+O(\log N)\bigr).
\end{equation}
Thus the archimedean exponential decay of the one-pole bump is defeated
by an exact all-prime factorial denominator for every possible choice of
the pole order $M=M(N)\geq N+1$.  This closes this scalar rational-pole
family, but it does not cover coupled determinants or nonseparable
Borel--Rodrigues kernels.
\end{proposition}

\begin{proof}
Kummer's transformation gives
\[
 K_{N,M}(z)=\frac{z^Ne^{-z}}{N!}
 {}_1F_1(-m;N+1;z).
\]
Expanding the terminating hypergeometric series yields
\eqref{eq:borel-laguerre-clearing-expansion}.  Every denominator
$(N+j)!$ divides $(M-1)!$, while the top coefficient, corresponding to
$j=m$, is $(-1)^m/(M-1)!$.  Hence the global least common denominator is
exactly $(M-1)!$ and its $p$-part is exactly
$p^{v_p((M-1)!)}$, proving \eqref{eq:borel-laguerre-local-wall}.
After clearing, the leading coefficient is $(-1)^m$, so the integer
polynomial is primitive.

The coefficient with $j=0$ is $(M-1)!/N!$, which proves the lower height
bound.  For every $j$,
\[
 \binom mj\frac{(M-1)!}{(N+j)!}
 \leq\binom mj\frac{(M-1)!}{N!}
 \leq2^m\frac{(M-1)!}{N!},
\]
giving the upper bound.  Stirling's formula proves
\eqref{eq:borel-laguerre-height-asymptotic}.  Finally,
\[
 \max_{y\geq0}\frac{y^N}{(1+y)^M}
 =\frac{N^N(M-N)^{M-N}}{M^M};
\]
its logarithm is $O_\alpha(N)$ when $M=\alpha N$, whereas
$\log((M-1)!)=\alpha N\log N+O_\alpha(N)$.  This proves
\eqref{eq:borel-laguerre-normalized-bump}.

For the uniform assertion, write $m=M-N$.  Direct division gives
\[
 \frac{\mathfrak F_{N,M+1}}{\mathfrak F_{N,M}}
 =M\frac{(m+1)^{m+1}}{m^m}
   \frac{M^M}{(M+1)^{M+1}}.
\]
Its logarithm is
\[
 \log M-\bigl(\phi(M)-\phi(m)\bigr),
 \qquad
 \phi(x)=(x+1)\log(x+1)-x\log x.
\]
Since
\[
 \phi(M)-\phi(m)
 =\int_m^M\log(1+x^{-1})\,dx
 <\int_m^M\frac{dx}{x}
 =\log(M/m)\leq\log M,
\]
the ratio is strictly greater than one.  Its minimum is therefore at
$M=N+1$, and Stirling's formula gives
\eqref{eq:borel-laguerre-uniform-wall}.
\end{proof}


\section{The explicit mixed family}\label{sec:mixed}

Let
\[
 \Eul(w)=\sum_{n\ge0}(-1)^nn!w^n,
 \qquad
 \Phi_\beta(z)=e^z\bigl(\Ein(z)-\beta\bigr)-\Eul(1/z),
 \quad \beta\in\Qb.
\]
The direction-zero Borel sum of \(\Eul(1/z)\) on \(|\arg z|<\pi\) is
\[
 \mathcal S_0\Eul(1/z)=ze^zE_1(z).
\]
Using \(\Ein(z)=\gamma+\log z+E_1(z)\) gives
\begin{equation}\label{eq:mixed-sum}
 \mathcal S_0\Phi_\beta(z)
 =e^z\bigl[\gamma-\beta+\log z+(1-z)E_1(z)\bigr],
\qquad
 \mathcal S_0\Phi_\beta(1)=e(\gamma-\beta).
\end{equation}

Let \(\mathcal M\) denote the mixed space consisting of formal objects
\(A(z)+a(1/z)\), where \(A\) is an \(E\)-function, \(a\) is arithmetic
Gevrey of order one, and \(a(0)=0\).  The harmless normalization \(a(0)=0\)
makes the splitting unique up to an algebraic constant.

\begin{proposition}[A mixed divisibility obstruction]\label{prop:nonmixed}
For every \(\beta\in\Qb\),
\[
 \Phi_\beta\notin(z-1)\mathcal M.
\]
\end{proposition}

\begin{proof}
Suppose
\(\Phi_\beta=(z-1)(A(z)+a(1/z))\) with \(A(z)+a(1/z)\in\mathcal M\).
Rearranging gives
\[
 e^z(\Ein(z)-\beta)-(z-1)A(z)
 =\Eul(1/z)+(z-1)a(1/z).
\]
The left side has only nonnegative powers of \(z\).  Because \(a(0)=0\),
the right side has only nonpositive powers.  Their common value must therefore
be a constant \(c\in\Qb\).  Evaluation at \(z=1\) gives
\[
 e(\Ein(1)-\beta)=c,
\]
contradicting the algebraic independence of \(e\) and \(\Ein(1)\) supplied
by \cref{thm:master}.
\end{proof}

There is an independent coefficient-level form of the obstruction.  In the
canonical split expansion of the formal quotient, its inverse tail is
\(w\Eul(w)/(w-1)\).  Theorem~\ref{thm:denominator}, with \(\xi=1\), shows
that even after allowing an arbitrary algebraic overlap constant \(c\), the
series \((\Eul(w)-c)/(w-1)\) has superexponential normalized denominators.

Consider the restricted zero-removal assertion
\begin{equation}\label{eq:zr-beta}
 \mathsf{ZR}(\beta):\qquad
 \mathcal S_0\Phi_\beta(1)=0
 \ \Longrightarrow\ 
 \Phi_\beta\in(z-1)\mathcal M.
\end{equation}

\begin{theorem}[Exact logical status of the proposed bridge]\label{thm:bridge}
For a fixed \(\beta\in\Qb\), \(\mathsf{ZR}(\beta)\) holds if and only if
\(\gamma\ne\beta\).  Consequently
\[
 \bigl[\mathsf{ZR}(\beta)\text{ for every }\beta\in\Qb\bigr]
 \quad\Longleftrightarrow\quad
 \gamma\text{ is transcendental}.
\]
\end{theorem}

\begin{proof}
By \eqref{eq:mixed-sum}, the antecedent of \(\mathsf{ZR}(\beta)\) holds
exactly when \(\gamma=\beta\).  By \cref{prop:nonmixed}, its consequent is
always false.  Thus the implication holds exactly when its antecedent is
false.  Taking the conjunction over algebraic \(\beta\) gives the result.
\end{proof}

This theorem explains why the apparently modest one-family lemma is not a
shortcut: it is the original transcendence problem in vacuous-implication
form.  There is also an analytic reflection of the obstruction.  Under the
hypothetical equality \(\beta=\gamma\), the summed quotient is locally
holomorphic at \(1\) and equals
\begin{equation}\label{eq:q-monodromy}
 q(z)=e^z\left(\frac{\log z}{z-1}-E_1(z)\right).
\end{equation}
A positive circuit about \(0\) sends \(\log z\) to \(\log z+2\pi i\) and
\(E_1(z)\) to \(E_1(z)-2\pi i\).  Therefore
\[
 \Delta q(z)=\frac{2\pi i\,ze^z}{z-1},
\]
which has a genuine pole at \(z=1\).  In the Fourier--Laplace/\(E\)-operator
formulation of mixed arithmetic-Gevrey functions
\cite{Andre2000,FR2024ED}, nonzero finite points are ordinary; every
continued solution germ must remain holomorphic there.  Thus
the local monodromy records the same obstruction as the denominator theorem.

\section{A compact Euler integral and its unavoidable logarithmic
contamination}
\label{sec:compact-integral-contamination}

The arithmetic framework of this section---Beukers--Sondow double
integrals and Nesterenko-type rational functions producing linear forms in
\(1,\gamma\), and logarithms---goes back at least to Sondow
\cite{Sondow2003Criteria,Sondow2009Hypergeometric}.  We make no priority
claim for that framework.  We instead isolate the exact asymmetric
contiguous family needed here, correct its asymptotic and height
normalization, and prove that its logarithmic contamination is strictly
positive for every index.

\subsection{The floor-weighted rational transform}

For a rational function \(g\) which is integrable against the indicated
weight, define
\begin{equation}\label{eq:floor-transform}
 \mathcal T(g)=\int_0^\infty(\lfloor u\rfloor+1)g(u)\,du.
\end{equation}
Write \(H_m^{(r)}=\sum_{j=1}^m j^{-r}\), with \(H_0^{(r)}=0\), and
\(H_m=H_m^{(1)}\).

\begin{proposition}[Complete pole-order transducer]
\label{prop:complete-pole-transducer}
Suppose that \(g\in\Qbar(u)\) is \(O(u^{-3})\) at infinity, all its poles
belong to \(\{-1,-2,\ldots\}\), and
\[
 g(u)=\sum_{p\ge1}\sum_{r=1}^{m_p}
       \frac{c_{p,r}}{(u+p)^r}
\]
is its finite partial-fraction expansion.  Then
\begin{align}
 \mathcal T(g)
 ={}&\sum_p c_{p,1}
 \left(\log((p-1)!)+p+1-\frac12\log(2\pi)\right)
 \notag\\
 &+\sum_p c_{p,2}\left(\gamma-1-H_{p-1}\right)
 \label{eq:complete-pole-transducer}\\
 &+\sum_p\sum_{r=3}^{m_p}
 \frac{c_{p,r}}{r-1}
 \left(\zeta(r-1)-H_{p-1}^{(r-1)}\right).
 \notag
\end{align}
In particular, simple poles carry factorial logarithms, double poles carry
\(\gamma\), and poles of order \(r\ge3\) carry \(\zeta(r-1)\).
\end{proposition}

\begin{proof}
For a simple pole, summation over the intervals \([j,j+1]\) gives
\[
 \sum_{j=0}^J(j+1)
 \log\frac{j+p+1}{j+p}
 =(J+1)\log(J+p+1)-\log\Gamma(J+p+1)+\log\Gamma(p).
\]
Stirling's formula isolates the finite part
\(\log\Gamma(p)+p+1-\frac12\log(2\pi)\).
For a double pole the corresponding finite part is
\[
 \frac1p-1-\psi(p+1)=\gamma-1-H_{p-1}.
\]
For \(r\ge3\), absolute convergence and summation by parts give
\begin{align*}
 &\sum_{j\ge0}(j+1)\int_j^{j+1}\frac{du}{(u+p)^r}\\
 &\hspace{2cm}=\frac1{r-1}\sum_{j\ge0}(j+p)^{1-r}
 =\frac{\zeta(r-1,p)}{r-1}.
\end{align*}
Since \(p\) is a positive integer,
\(\zeta(r-1,p)=\zeta(r-1)-H_{p-1}^{(r-1)}\).

The hypothesis \(g(u)=O(u^{-3})\) gives
\[
 \sum_pc_{p,1}=0,
 \qquad
 \sum_pc_{p,2}=\sum_ppc_{p,1}.
\]
These identities cancel the linear and logarithmic divergences of the
simple- and double-pole partial sums, so the sum of their finite parts is
exactly \eqref{eq:complete-pole-transducer}.
\end{proof}

\begin{remark}[Necessary correction of scope]
A product of Beta factors can have poles of multiplicity greater than two.
In that situation a partial-fraction expansion containing only simple and
double poles is incomplete; the last line of
\eqref{eq:complete-pole-transducer} is indispensable.  Restricting to
products with pole multiplicity at most two recovers the
\((1,\gamma,\log)\)-valued special case.
\end{remark}

There is also a one-variable endpoint obstruction which explains why the
compact construction below is genuinely two-dimensional.

\begin{lemma}[Endpoint removal loses \(\gamma\)]
\label{lem:one-variable-endpoint-no-go}
Let \(P(x)=\sum_{k=0}^d a_kx^k\in\Qbar[x]\) satisfy \(P(1)=0\).  Then
\begin{equation}\label{eq:endpoint-no-go-exact}
 \int_0^1P(x)\left(\frac1{1-x}+\frac1{\log x}\right)dx
 =\int_0^1\frac{P(x)}{1-x}\,dx
  +\sum_{k=0}^da_k\log(k+1).
\end{equation}
The first term on the right is algebraic; in particular no \(\gamma\)-term
survives.
\end{lemma}

\begin{proof}
The first integrand is a polynomial because \(P(1)=0\).  The same condition
allows one to write \(P(x)=\sum_ka_k(x^k-1)\), and Frullani's identity
\[
 \int_0^1\frac{x^k-1}{\log x}\,dx=\log(k+1)
\]
proves \eqref{eq:endpoint-no-go-exact}.
\end{proof}

\subsection{The asymmetric compact family}

For \(n\ge1\), set
\begin{equation}\label{eq:asymmetric-Jn}
 J_n=\int_0^1\!\int_0^1
 \frac{[x(1-x)y(1-y)]^n(1-x)}
 {(1-xy)(-\log(xy))}\,dx\,dy.
\end{equation}
Using
\(1/(-\log(xy))=\int_0^\infty(xy)^t\,dt\) and
\((1-xy)^{-1}=\sum_{j\ge0}(xy)^j\), Tonelli's theorem gives
\begin{equation}\label{eq:Jn-floor-beta}
 J_n=\mathcal T(g_n),\qquad
 g_n(u)=B(u+n+1,n+2)B(u+n+1,n+1).
\end{equation}
Put \(L=n+1\), \(M=2n+2\), and \(C=n!(n+1)!\).  Then
\begin{equation}\label{eq:gn-rational-form}
 g_n(u)=\frac{C}{(u+M)\prod_{p=L}^{M-1}(u+p)^2}.
\end{equation}

\begin{theorem}[Exact coefficients and strict logarithmic contamination]
\label{thm:Jn-exact-strict-contamination}
The partial-fraction expansion of \(g_n\) is
\[
 g_n(u)=\sum_{p=L}^{M}\frac{\alpha_p}{u+p}
       +\sum_{p=L}^{M-1}\frac{\beta_p}{(u+p)^2},
\]
where, for \(0\le k\le n\),
\begin{align}
 \beta_{L+k}
 &=\binom nk\binom{n+1}{k},
 \label{eq:Jn-beta}\\
 \alpha_{L+k}
 &=\beta_{L+k}\left(
 2H_k-2H_{n-k}-\frac1{n+1-k}\right),
 \label{eq:Jn-alpha}\\
 \alpha_M&=\frac1{n+1}.
 \label{eq:Jn-alpha-last}
\end{align}
Consequently
\begin{equation}\label{eq:Jn-exact-linear-form}
 J_n=\binom{2n+1}{n}\gamma+r_n+\mathcal L_n,
\end{equation}
where
\begin{align}
 r_n
 &=\sum_{p=L}^{M}\alpha_p(p+1)
   +\sum_{p=L}^{M-1}\beta_p
      \left(\frac1p-1-H_p\right)\in\Q,
 \label{eq:Jn-rational-part}\\
 \mathcal L_n
 &=\sum_{p=L}^{M}\alpha_p\log((p-1)!)
   =\sum_{j=L}^{M-1}T_j\log j,
 \qquad
 T_j=\sum_{p>j}\alpha_p.
 \label{eq:Jn-log-part}
\end{align}
Every proper tail is strictly positive:
\begin{equation}\label{eq:Jn-positive-tails}
 T_j>0\qquad(L\le j\le M-1).
\end{equation}
In particular
\begin{equation}\label{eq:Jn-log-positive}
 \boxed{\mathcal L_n>0\quad\text{for every }n\ge1.}
\end{equation}
Thus this entire compact family is never logarithm-free.
\end{theorem}

\begin{proof}
At the double pole \(p=L+k\), evaluation of
\((u+p)^2g_n(u)\) at \(u=-p\) gives
\[
 \beta_{L+k}
 =\frac{n!(n+1)!}
 {(n+1-k)(k!)^2((n-k)!)^2}
 =\binom nk\binom{n+1}{k}.
\]
Logarithmic differentiation of the same expression gives
\[
 \frac{\alpha_{L+k}}{\beta_{L+k}}
 =-\frac1{n+1-k}
  -2\sum_{\substack{0\le j\le n\\j\ne k}}\frac1{j-k}
 =2H_k-2H_{n-k}-\frac1{n+1-k},
\]
which proves \eqref{eq:Jn-alpha}.  The simple pole at \(M\) gives
\eqref{eq:Jn-alpha-last}.  Since \(g_n(u)=O(u^{-2n-3})\), comparison at
infinity gives
\begin{equation}\label{eq:Jn-alpha-sum-zero}
 \sum_{p=L}^{M}\alpha_p=0.
\end{equation}

Proposition~\ref{prop:complete-pole-transducer} now yields
\eqref{eq:Jn-exact-linear-form}--\eqref{eq:Jn-log-part}; the
\(\log(2\pi)\)-term disappears by
\eqref{eq:Jn-alpha-sum-zero}.  Vandermonde's identity gives the
\(\gamma\)-coefficient:
\[
 \sum_{k=0}^n\beta_{L+k}
 =\sum_{k=0}^n\binom nk\binom{n+1}{k}
 =\binom{2n+1}{n}.
\]

It remains to prove the strict obstruction.  Write
\[
 b_k=2H_k-2H_{n-k}-\frac1{n+1-k}.
\]
Then
\[
 b_{k+1}-b_k
 =\frac2{k+1}+\frac1{n-k}+\frac1{n+1-k}>0,
\]
while \(b_0<0\) and \(b_n=2H_n-1>0\).  Since every
\(\beta_{L+k}\) is positive, the sequence
\(\alpha_L,\ldots,\alpha_{M-1}\) changes sign at most once, from negative
to positive, and \(\alpha_M>0\).  Together with the zero total sum
\eqref{eq:Jn-alpha-sum-zero}, this implies that every proper tail \(T_j\)
is positive.  Formula \eqref{eq:Jn-log-part} then gives
\eqref{eq:Jn-log-positive}.
\end{proof}

\subsection{Corrected asymptotic and integer height}

\begin{theorem}[Corrected asymptotic and integerization]
\label{thm:Jn-asymptotic-integerization}
As \(n\to\infty\),
\begin{equation}\label{eq:Jn-correct-asymptotic}
 \boxed{\displaystyle
 J_n\sim\frac{\pi}{12\log2}\,\frac{16^{-n}}n.}
\end{equation}
Let \(d_n=\operatorname{lcm}(1,\ldots,2n+2)\).  There are integers
\(A_n,B_n,C_{n,q}\), with \(q\) running over the primes at most \(2n+1\),
such that
\begin{equation}\label{eq:Jn-integer-linear-form}
 d_nJ_n=A_n+B_n\gamma+
 \sum_{\substack{q\le2n+1\\q\ \mathrm{prime}}}C_{n,q}\log q,
 \qquad
 B_n=d_n\binom{2n+1}{n}.
\end{equation}
For every prime \(q\in[n+1,2n+1]\),
\begin{equation}\label{eq:Jn-prime-coefficient-positive}
 C_{n,q}=d_nT_q>0.
\end{equation}
Thus the number of nonzero logarithmic coordinates tends to infinity and is
at least
\(\pi_{\mathrm{pr}}(2n+1)-\pi_{\mathrm{pr}}(n)\sim n/\log n\), where
\(\pi_{\mathrm{pr}}\) is the prime-counting function.

If
\[
 \mathcal H_n=\max\left(
 1,|A_n|,|B_n|,
 \max_{q\le2n+1}|C_{n,q}|
 \right),
\]
then
\begin{align}
 \log\mathcal H_n&=(2+\log4)n+o(n),
 \label{eq:Jn-integer-height}\\
 -\log|d_nJ_n|&=(\log16-2)n+o(n).
 \label{eq:Jn-integer-smallness}
\end{align}
Hence the exponent of this explicit integer form is
\begin{equation}\label{eq:Jn-correct-tau}
 \frac{-\log|d_nJ_n|}{\log\mathcal H_n}
 \longrightarrow
 \frac{\log16-2}{\log4+2}
 =0.2281516725\ldots.
\end{equation}
The exponent is positive, but the ambient dimension grows; this is not a
fixed-dimensional Nesterenko criterion and proves neither irrationality nor
transcendence of \(\gamma\).
\end{theorem}

\begin{proof}
The function
\[
 \phi(x,y)=\log(x(1-x)y(1-y))
\]
has its unique maximum \(-\log16\) at \((1/2,1/2)\), with Hessian
\(-8I_2\).  The remaining factor in \eqref{eq:asymmetric-Jn} has value
\[
 \frac{1-1/2}{(1-1/4)(-\log(1/4))}=\frac1{3\log2}
\]
there.  The two-dimensional Laplace method therefore gives
\[
 J_n\sim16^{-n}\frac{2\pi/n}{8}\frac1{3\log2},
\]
which is \eqref{eq:Jn-correct-asymptotic}.

Equations \eqref{eq:Jn-beta}--\eqref{eq:Jn-alpha-last} show that the
denominators of all \(\alpha_p\) divide
\(\operatorname{lcm}(1,\ldots,n+1)\), while the denominators in
\eqref{eq:Jn-rational-part} divide \(d_n\).  Expanding every
\(\log((p-1)!)\) into prime logarithms proves
\eqref{eq:Jn-integer-linear-form}.  If \(q\in[n+1,2n+1]\) is prime, then
\(2q>2n+1\), so
\[
 v_q((p-1)!)=\mathbf1_{p>q}\qquad(L\le p\le M).
\]
Equation \eqref{eq:Jn-prime-coefficient-positive} follows from
\eqref{eq:Jn-positive-tails}; the prime-number theorem gives the count.

Finally
\[
 \log d_n=2n+o(n),\qquad
 \log\binom{2n+1}{n}=n\log4+o(n).
\]
The coefficient formulas show that every coefficient in
\eqref{eq:Jn-integer-linear-form} is at most
\(d_n4^n\exp(o(n))\), while \(B_n\) has exactly that exponential scale.
This proves \eqref{eq:Jn-integer-height}.  Combining
\eqref{eq:Jn-correct-asymptotic} with \(\log d_n=2n+o(n)\) proves
\eqref{eq:Jn-integer-smallness}, and division gives
\eqref{eq:Jn-correct-tau}.
\end{proof}

\begin{remark}[Comparison with the reciprocal height barrier]
The reciprocal truncations studied earlier saturate their full defining
height so completely that primitive integerization destroys the small
value.  The compact forms \eqref{eq:Jn-integer-linear-form} behave
differently: after integerization they remain exponentially small.  Their
obstruction is instead the rigorously growing set of positive logarithmic
coordinates in \eqref{eq:Jn-prime-coefficient-positive}.  These are two
distinct boundary mechanisms and should not be conflated.
\end{remark}


\section{A recent logarithmic consequence}\label{sec:log}

The preceding results concern algebraic relations.  A recent theorem of
Fischler and Rivoal on logarithms of \(E\)-values \cite{FRlog} gives a complementary
one-number statement.

\begin{proposition}\label{prop:log-ein}
The number
\[
 \log\Ein(1)=\log\left(\gamma+\frac{\delta}{e}\right)
\]
is transcendental.  It is also not an ultra-Liouville number.
\end{proposition}

\begin{proof}
The specialization of Fischler--Rivoal's logarithm theorem \cite{FRlog} to \(\Ein\) leaves
only the exceptional alternatives \(\Ein(1)=1\) and \(\Ein(1)=e^{-1}\).
They are excluded by the elementary integral bounds
\[
 e^{-1}
 <\int_0^1\frac{1-e^{-t}}{t}\,dt
 <1.
\]
Indeed \((1-e^{-t})/t\) is the average of \(e^{-u}\) on \([0,t]\), so it is
strictly above \(e^{-1}\) on a set of positive measure and strictly below
\(1\).  Their quantitative theorem supplies the non-ultra-Liouville
conclusion.
\end{proof}

This proposition is an application of the cited general theorem, not a new
proof of an older open problem.  It does not imply the transcendence of either
summand in \(\Ein(1)=\gamma+\delta/e\) separately.

\section{Joukowski descent and the arithmetic of reciprocal levels}
\label{sec:joukowski-descent}

The reciprocal equation admits a useful quotient which separates its
complex-analytic geometry from its arithmetic level.  Put
\[
 J(z)=z+z^{-1}\qquad(z\in\C^\times).
\]
The involution $z\mapsto z^{-1}$ is the deck involution of $J$.  In this
section we descend the reciprocal $\Ein$-sum through $J$, determine the
critical points of the descended entire function, and then apply
Theorem~\ref{thm:master} to its values at algebraic points.

\subsection{The entire quotient and its critical points}

Let $C_0(w)=2$, $C_1(w)=w$, and
\[
 C_{n+1}(w)=wC_n(w)-C_{n-1}(w)\qquad(n\geq1).
\]
Thus $C_n(w)=2T_n(w/2)$ and
\begin{equation}\label{eq:joukowski-chebyshev}
 C_n(J(z))=z^n+z^{-n}.
\end{equation}

\begin{proposition}[Entire Joukowski descent]\label{prop:entire-joukowski-descent}
There is a unique entire function $H$ such that
\begin{equation}\label{eq:H-descent}
 H(J(z))=\Ein(z)+\Ein(z^{-1})\qquad(z\in\C^\times).
\end{equation}
It is given by the locally uniformly convergent series
\begin{equation}\label{eq:H-chebyshev-series}
 H(w)=\sum_{n\geq1}\frac{(-1)^{n-1}}{n\,n!}C_n(w).
\end{equation}
Moreover
\begin{align}
 H'(w)
 &=2e^{-w/2}
   \frac{\sinh\!\left(\frac12\sqrt{w^2-4}\right)}{\sqrt{w^2-4}}
   \label{eq:Hprime-sinh}\\
 &=e^{-w/2}\prod_{n\geq1}
   \left(1+\frac{w^2-4}{4\pi^2n^2}\right).
   \label{eq:Hprime-product}
\end{align}
The right side of \eqref{eq:Hprime-sinh} is understood by its removable
continuation at $w=\pm2$; in particular $H'(2)=e^{-1}$ and
$H'(-2)=e$.
\end{proposition}

\begin{proof}
If $K\subset\C$ is compact, there is an $R_K>1$ such that both roots of
$X^2-wX+1$ have modulus at most $R_K$ for every $w\in K$.  Hence
\eqref{eq:joukowski-chebyshev} gives
$|C_n(w)|\leq2R_K^n$ on $K$.  The series
\eqref{eq:H-chebyshev-series} therefore converges locally uniformly and
defines an entire function.  Substitution of $w=J(z)$ and the power series
of $\Ein$ proves \eqref{eq:H-descent}.  Uniqueness follows because $J$ is
surjective: the equation $z^2-wz+1=0$ has a nonzero root for every
$w\in\C$.

For $z\ne\pm1$, differentiation of \eqref{eq:H-descent} gives
\[
 H'(J(z))
 =\frac{e^{-1/z}-e^{-z}}{z-z^{-1}}.
\]
Write $w=z+z^{-1}$ and $v=z-z^{-1}$, so $v^2=w^2-4$.  Then
\[
 e^{-1/z}-e^{-z}
 =e^{-w/2}\bigl(e^{v/2}-e^{-v/2}\bigr),
\]
which proves \eqref{eq:Hprime-sinh} away from $w=\pm2$.  The quotient
$2\sinh(v/2)/v$ is an entire even function of $v$, so the apparent
singularities are removable.  Finally Euler's product
\[
 \frac{\sinh u}{u}=\prod_{n\geq1}
   \left(1+\frac{u^2}{\pi^2n^2}\right)
\]
with $u=\frac12\sqrt{w^2-4}$ gives
\eqref{eq:Hprime-product}.
\end{proof}

\begin{corollary}[Critical lattice and the real diffeomorphism]
\label{cor:H-critical-lattice}
The critical points of $H$ are exactly the simple points
\begin{equation}\label{eq:H-critical-points}
 \xi_n^\pm=\pm2i\sqrt{\pi^2n^2-1}\qquad(n\geq1).
\end{equation}
In particular every critical point is transcendental.  On the real axis
$H'(x)>0$, and $H$ maps $\mathbb R$ increasingly and bijectively onto
$\mathbb R$.
\end{corollary}

\begin{proof}
The product \eqref{eq:Hprime-product} gives
\eqref{eq:H-critical-points}; each zero is simple because it occurs in
exactly one factor and is nonzero.  If one of these points were algebraic,
then $\pi^2=((\xi_n^\pm)^2/(-4)+1)/n^2$ would be algebraic, a contradiction.
For real $x$, every factor in \eqref{eq:Hprime-product} is positive, since
\[
 1+\frac{x^2-4}{4\pi^2n^2}
 \geq1-\frac1{\pi^2n^2}>0.
\]
Thus $H'(x)>0$.

For $w\to+\infty$, choose the root $z>1$ of $J(z)=w$.  Then
$z=w+O(w^{-1})$ and
\[
 H(w)=\Ein(z)+\Ein(z^{-1})
     =\gamma+\log w+O(w^{-1})\longrightarrow+\infty.
\]
For $w\to-\infty$, write the root of modulus greater than one as $z=-a$,
$a\to+\infty$.  Since
$\Ein(-a)=\gamma+\log a-\Ei(a)$ and
$\Ei(a)\sim e^a/a$, one has $H(w)\to-\infty$.  Strict monotonicity now
proves the last assertion.
\end{proof}

\begin{proposition}[A rational differential equation]
\label{prop:H-ode}
The function $H$ satisfies
\begin{equation}\label{eq:H-third-order-ode}
 w(w^2-4)H'''
 +(w^3+2w^2-4w+4)H''
 +(w^2-w+2)H'=0.
\end{equation}
Consequently every additive translate $H-2c$ satisfies the same equation.
\end{proposition}

\begin{proof}
Put $P(w)=e^{w/2}H'(w)$.  By \eqref{eq:Hprime-sinh},
\[
 P(w)=2\frac{\sinh(\frac12\sqrt{w^2-4})}{\sqrt{w^2-4}}.
\]
Direct differentiation gives
\[
 w(w^2-4)P''+2(w^2+2)P'-\frac{w^3}{4}P=0.
\]
Substitution of $P=e^{w/2}H'$ and collection of terms gives
\eqref{eq:H-third-order-ode}.  The equation contains no undifferentiated
$H$, so it is invariant under additive translation.
\end{proof}

\subsection{Algebraic values after descent}

For $w\in\Qbar$, let
\[
 \Lambda(w)=\{\lambda\in\Qbar^\times:
                    \lambda+\lambda^{-1}=w\}
\]
as a set.  At $w=2$ or $w=-2$ this set has one element, but the identity
\eqref{eq:H-descent} counts that element twice.

\begin{theorem}[Multipoint arithmetic descent]
\label{thm:H-multipoint-descent}
Let $w_1,\ldots,w_s\in\Qbar$ be distinct, let
$\Lambda=\bigcup_{j=1}^s\Lambda(w_j)$, and put
\[
 d=\dim_{\Q}\Span_{\Q}\Lambda.
\]
Then
\begin{equation}\label{eq:H-multipoint-trdeg}
 \operatorname{trdeg}_{\Qbar}
 \Qbar\bigl(\{e^\lambda\}_{\lambda\in\Lambda},
             H(w_1),\ldots,H(w_s)\bigr)=d+s.
\end{equation}
In particular $H(w_1),\ldots,H(w_s)$ are algebraically independent over
the field generated by the exponential coordinates.  Hence $H(w)$ is
transcendental for every algebraic $w$.
\end{theorem}

\begin{proof}
By Theorem~\ref{thm:master}, the coordinates
$Y_\lambda=\Ein(\lambda)$, $\lambda\in\Lambda$, are algebraically
independent over
$E=\Qbar(\{e^\lambda\}_{\lambda\in\Lambda})$, and
$\operatorname{trdeg}_{\Qbar}E=d$.  Each $H(w_j)$ is a nonzero linear form
in the $Y_\lambda$ supported on $\Lambda(w_j)$, with coefficient $2$ at a
double root.  The supports for distinct $j$ are disjoint.  These $s$ linear
forms can therefore be extended to a linear coordinate system in the
$Y_\lambda$, and are algebraically independent over $E$.  This proves
\eqref{eq:H-multipoint-trdeg}.  The last assertion is the case $s=1$.
\end{proof}

\begin{corollary}[Algebraic levels have transcendental preimages]
\label{cor:H-algebraic-levels}
For every $c\in\Qbar$, every solution of
\begin{equation}\label{eq:H-algebraic-level}
 H(w)=2c
\end{equation}
is transcendental.  Equivalently, every solution of
\begin{equation}\label{eq:reciprocal-algebraic-level}
 \Ein(z)+\Ein(z^{-1})=2c
\end{equation}
is transcendental.
\end{corollary}

\begin{proof}
An algebraic solution of \eqref{eq:H-algebraic-level} would contradict the
last assertion of Theorem~\ref{thm:H-multipoint-descent}.  If $z$ in
\eqref{eq:reciprocal-algebraic-level} were algebraic, then
$w=J(z)$ would be algebraic and would solve \eqref{eq:H-algebraic-level}.
\end{proof}

\begin{corollary}[The unique real orbit of a real level]
\label{cor:H-real-levels}
For every $c\in\mathbb R$ there is a unique $w_c\in\mathbb R$ with
$H(w_c)=2c$.  It satisfies
\[
 \begin{array}{c|c|c}
 \text{condition on }c&\text{location of }w_c&
 \text{preimage under }J\\ \hline
 \Ein(-1)<c<\Ein(1)&-2<w_c<2&
 \text{one conjugate pair on }|z|=1,\\
 c>\Ein(1)&w_c>2&\text{two positive reciprocal points},\\
 c<\Ein(-1)&w_c<-2&\text{two negative reciprocal points}.
 \end{array}
\]
At $c=\Ein(1)$ or $c=\Ein(-1)$, respectively, $z=1$ or $z=-1$ is a
double zero of \eqref{eq:reciprocal-algebraic-level}.  If $c$ is algebraic,
the points in the table are all transcendental.

More generally, for arbitrary real $c$, every nonreal algebraic solution
of $H(w)=2c$ is impossible.  Thus at most the unique real preimage $w_c$
can be algebraic.
\end{corollary}

\begin{proof}
Existence and uniqueness follow from Corollary~\ref{cor:H-critical-lattice}.
The thresholds follow from
$H(2)=2\Ein(1)$ and $H(-2)=2\Ein(-1)$, and the three descriptions are the
elementary inverse images of the intervals $(-2,2)$, $(2,\infty)$ and
$(-\infty,-2)$ under $J$.  Since
\[
 J(z)-2=\frac{(z-1)^2}{z},\qquad
 J(z)+2=\frac{(z+1)^2}{z},
\]
and $H'(\pm2)\ne0$, the endpoint zeros have order two.  The assertion for
algebraic $c$ is Corollary~\ref{cor:H-algebraic-levels}.

Finally, if a nonreal algebraic $w$ solved $H(w)=2c$ with $c$ real, then
$\bar w$ would be a distinct algebraic point with
$H(\bar w)=\overline{H(w)}=2c$.  This contradicts the algebraic
independence of $H(w)$ and $H(\bar w)$ in
Theorem~\ref{thm:H-multipoint-descent}.
\end{proof}

\subsection{The Euler level and witness amplification}

Put
\[
 G_\gamma(w)=H(w)-2\gamma.
\]

\begin{theorem}[Euler level in the quotient]
\label{thm:H-euler-level}
The function $G_\gamma$ has exactly one real zero $w_*$, and
\begin{equation}\label{eq:w-star-location}
 -2<w_*<2,\qquad w_*=2\cos\theta_*,
\end{equation}
A high-precision Newton iteration gives the non-rigorous decimal values
\[
 w_*=1.0073046182687490652\ldots,
 \qquad \theta_*=1.0429750684400255695\ldots.
\]
The two preimages of $w_*$ are the unit-circle zeros
$\alpha_*=e^{i\theta_*}$ and $\overline{\alpha_*}$ of the reciprocal
$\Ein$ equation.  Every algebraic zero of $G_\gamma$ must equal $w_*$;
equivalently every algebraic zero of the reciprocal equation belongs to
$\{\alpha_*,\overline{\alpha_*}\}$.

If $w_*$, equivalently $\alpha_*$, is algebraic, then $\gamma$ and
$\delta$ are algebraically independent.  More precisely, put
$a=\alpha_*$, $b=a^{-1}$,
\[
 D_*=\Ein(a)-\Ein(b),\qquad
 d_*=\dim_{\Q}\Span_{\Q}\{1,a,b\}.
\]
Then
\begin{equation}\label{eq:witness-amplification}
 \operatorname{trdeg}_{\Qbar}
 \Qbar(e,e^a,e^b,\gamma,\delta,D_*)=d_*+3,
\end{equation}
and $\gamma,\delta,D_*$ are algebraically independent over
$\Qbar(e,e^a,e^b)$.
\end{theorem}

\begin{proof}
Since
\[
 H(-2)=2\gamma-2\Ei(1)<2\gamma
 <2\gamma+2E_1(1)=H(2),
\]
Corollary~\ref{cor:H-critical-lattice} gives the unique zero $w_*$ and
\eqref{eq:w-star-location}.  If an algebraic zero $w$ were nonreal, the
last assertion of Corollary~\ref{cor:H-real-levels} would exclude it; if it
were real, uniqueness would force $w=w_*$.  Passing between $w$ and the
roots of $z^2-wz+1$ proves the equivalent statement in the $z$-plane.

Now assume that $w_*$ is algebraic.  Then $a$ and $b$ are distinct
algebraic numbers.  Apply Theorem~\ref{thm:master} to
$\Lambda=\{1,a,b\}$.  Over
$E=\Qbar(e,e^a,e^b)$ the three numbers
\[
 Y_1=\Ein(1),\qquad Y_a=\Ein(a),\qquad Y_b=\Ein(b)
\]
are algebraically independent.  The reciprocal equation and
$\delta=eE_1(1)$ give
\[
 \gamma=\frac{Y_a+Y_b}{2},\qquad
 \frac{\delta}{e}=Y_1-\frac{Y_a+Y_b}{2},\qquad
 D_*=Y_a-Y_b.
\]
This is an invertible linear change of the three $Y$-coordinates; replacing
$\delta/e$ by $\delta$ is invertible over $E$.  This proves
\eqref{eq:witness-amplification}.
\end{proof}

\begin{corollary}[Conditional completion of the six-constant field]
\label{cor:witness-six-constant-field}
If $\alpha_*$ is algebraic, then
\[
 \operatorname{trdeg}_{\Qbar}
 \Qbar\bigl(\gamma,\delta,E_1(1),\Ei(1),\Ci(1),\Chi(1)\bigr)=5.
\]
The complete ideal of algebraic relations among these six displayed
coordinates is generated by
\begin{equation}\label{eq:six-constant-only-relation}
 \Ei(1)-E_1(1)-2\Chi(1)=0.
\end{equation}
\end{corollary}

\begin{proof}
Use the notation
\[
 A=\Ein(1),\qquad B=\Ein(-1),\qquad
 C=\frac{\Ein(i)+\Ein(-i)}2.
\]
The point $\alpha_*$ is distinct from $\pm1,\pm i$.  Indeed
$w_*\in(-2,2)$ excludes $\pm1$, while
\[
 H(0)=2\gamma-2\Ci(1)<2\gamma
\]
excludes $\pm i$; here
$\Ci(1)=\gamma+\int_0^1(\cos t-1)\,dt/t>\gamma-1/4>0$.
Theorem~\ref{thm:master}, applied to
$\{1,-1,i,-i,\alpha_*,\alpha_*^{-1}\}$, now shows that
\[
 \gamma,\quad e,\quad A,\quad B,\quad C
\]
are algebraically independent: they are one exponential coordinate and
four independent linear forms with disjoint supports in the relevant
$\Ein$-coordinates.  The six-constant field equals
$\Qbar(\gamma,e,A,B,C)$, so its transcendence degree is five.

Identity \eqref{eq:six-constant-only-relation} follows from
$E_1(1)=A-\gamma$, $\Ei(1)=\gamma-B$, and
$\Chi(1)=\gamma-(A+B)/2$.  The kernel of the evaluation map from a
polynomial ring in the six displayed coordinates is a height-one prime.
It contains the irreducible linear polynomial in
\eqref{eq:six-constant-only-relation}; hence it is precisely the principal
ideal generated by that polynomial.
\end{proof}

\subsection{Every level has asymptotic transcendental fibres}

The following strengthens the existence part of Picard's theorem for this
specific function: there is no exceptional finite level, and an explicit
sequence of preimages is available at every level.

\begin{theorem}[Asymptotic level fibres]
\label{thm:H-level-asymptotics}
Fix $c\in\C$.  For every sufficiently large integer $k$ there is a simple
solution $w_{k,c}$ of $H(w)=2c$ such that
\begin{equation}\label{eq:H-level-zero-asymptotic}
 w_{k,c}
 =-\log(2\pi k)-\log\log(2\pi k)
  +i\left(2\pi k+\frac\pi2\right)
  +O_c\!\left(\frac1{\log k}\right).
\end{equation}
Consequently every finite value of $H$ is assumed infinitely often.  If
$c$ is algebraic, every member of these fibres is transcendental.  If
$c\in\mathbb R$, the conjugates $\overline{w_{k,c}}$ give a second family.
\end{theorem}

\begin{proof}
It is convenient first to work in the $z$-plane.  Put
\[
 \mathcal F_c(z)=\Ein(z)+\Ein(z^{-1})-2c
                =H(J(z))-2c.
\]
On $\C\setminus(-\infty,0]$, with principal logarithm, set
\[
 L_c(z)=\Log z+\gamma-2c.
\]
The connection formula for $\Ein$ gives
\begin{equation}\label{eq:Fc-connection}
 \mathcal F_c(z)=L_c(z)+E_1(z)+\Ein(z^{-1}).
\end{equation}
Uniformly in a fixed closed sector about the positive imaginary axis,
the standard expansion of $E_1$ \cite{DLMF} and the Taylor expansion of
$\Ein$ give
\begin{equation}\label{eq:Fc-asymptotic-input}
 E_1(z)=\frac{e^{-z}}z\bigl(1+O(z^{-1})\bigr),
 \qquad
 \Ein(z^{-1})=z^{-1}+O(z^{-2}).
\end{equation}

Let
\[
 R_k=2\pi k,\qquad
 a_k=\log R_k+\log\log R_k,\qquad
 \zeta_k=-a_k+i\left(R_k+\frac\pi2\right).
\]
On $D_k=\{z:|z-\zeta_k|\leq1\}$, for all sufficiently large $k$, the
principal branches of $\Log z$ and $\Log L_c(z)$ are analytic.  Uniformly
on $D_k$,
\begin{align}
 \Log z&=\log R_k+\frac{\pi i}{2}
          +O(a_k/R_k),\label{eq:log-z-on-Dk}\\
 \Log L_c(z)&=\log\log R_k+O_c((\log R_k)^{-1}).
 \label{eq:log-L-on-Dk}
\end{align}
Define
\[
 \Phi_{k,c}(z)=(2k+1)\pi i-\Log z-\Log L_c(z).
\]
Equations \eqref{eq:log-z-on-Dk}--\eqref{eq:log-L-on-Dk} give
\[
 \Phi_{k,c}(z)=\zeta_k+O_c((\log R_k)^{-1}),
 \qquad
 \Phi_{k,c}'(z)=-\frac1z-\frac1{zL_c(z)}=O(R_k^{-1}).
\]
Thus $\Phi_{k,c}$ maps $D_k$ into itself and is a contraction for large
$k$.  Its fixed point $u_{k,c}$ satisfies
\begin{equation}\label{eq:model-root-asymptotic}
 u_{k,c}=\zeta_k+O_c((\log R_k)^{-1})
\end{equation}
and
\[
 u_{k,c}e^{u_{k,c}}L_c(u_{k,c})=-1.
\]
Hence $u_{k,c}$ is a zero of the model function
\[
 h_c(z)=1+ze^zL_c(z).
\]

Write $q_c(z)=z+\Log z+\Log L_c(z)$, so
$h_c=1+e^{q_c}$ on $D_k$.  At $u_{k,c}$,
\begin{equation}\label{eq:model-derivative}
 q_c'(u_{k,c})
 =1+\frac1{u_{k,c}}+\frac1{u_{k,c}L_c(u_{k,c})}
 =1+O(R_k^{-1}).
\end{equation}
In particular the model zero is simple.

Now multiply the exact function by its nonvanishing normalizing factor:
\[
 A_c(z)=ze^z\mathcal F_c(z).
\]
From \eqref{eq:Fc-connection}--\eqref{eq:Fc-asymptotic-input},
\begin{equation}\label{eq:exact-model-error}
 A_c(z)-h_c(z)
 =ze^z\left(E_1(z)-\frac{e^{-z}}z\right)
  +ze^z\Ein(z^{-1}).
\end{equation}
The first term is $O(R_k^{-1})$ near $u_{k,c}$.  The model equation gives
\[
 |e^{u_{k,c}}|
 =\frac1{|u_{k,c}|\,|L_c(u_{k,c})|}
 =O_c((R_k\log R_k)^{-1}),
\]
so the second term in \eqref{eq:exact-model-error} is
$O_c((R_k\log R_k)^{-1})$.  Consequently there is a constant $C_c$ such
that
\begin{equation}\label{eq:exact-model-uniform-bound}
 |A_c(z)-h_c(z)|\leq \frac{C_c}{R_k}
\end{equation}
uniformly on $|z-u_{k,c}|\leq1$ for all sufficiently large $k$.

Fix $M>2C_c$.  Equation \eqref{eq:model-derivative} and Taylor's theorem
give, on $|z-u_{k,c}|=M/R_k$,
\[
 |h_c(z)|\geq\frac{M}{2R_k}
\]
for large $k$.  A second application of the same Taylor estimate, now
comparing $h_c$ with
$-q_c'(u_{k,c})(z-u_{k,c})$, shows that $h_c$ has exactly one zero, counted
with multiplicity, in this circle.  By
\eqref{eq:exact-model-uniform-bound} and Rouch\'e's theorem, $A_c$ also has
exactly one zero there.  Since $ze^z$ does not vanish, this is a simple zero
$z_{k,c}$ of $\mathcal F_c$, and
\[
 z_{k,c}=u_{k,c}+O_c(R_k^{-1}).
\]

Finally put $w_{k,c}=J(z_{k,c})$.  Since
$z_{k,c}^{-1}=O(R_k^{-1})$ and $J'(z_{k,c})=1+O(R_k^{-2})$, this is a
simple zero of $H(w)-2c$ and has the same asymptotic
\eqref{eq:H-level-zero-asymptotic}.  The arithmetic assertion follows from
Corollary~\ref{cor:H-algebraic-levels}, and conjugation gives the lower
family for real $c$.
\end{proof}

\begin{corollary}[Order and a zero-count lower bound]
\label{cor:H-order-zero-count}
For every $c\in\C$, the entire function $H-2c$ has order one and
exponential type one.  If $n_c(r)$ counts its zeros in $|w|\leq r$, with
multiplicity, then
\[
 \frac{r}{\pi}-O_c(\log r)\leq n_c(r)=O_c(r).
\]
\end{corollary}

\begin{proof}
Formula \eqref{eq:Hprime-sinh} gives exponential type at most one, and the
negative-real asymptotic used in Corollary~\ref{cor:H-critical-lattice}
gives type at least one.  Jensen's formula then gives $n_c(r)=O_c(r)$.
For real $c$, Theorem~\ref{thm:H-level-asymptotics} and conjugation give the
lower bound directly.  For general $c$, apply the same contraction argument
with negative indices as well as positive ones; equivalently repeat the
proof in a fixed sector about the negative imaginary axis.  The two
families have modulus $2\pi k+O_c(\log k)$, which gives the stated bound.
\end{proof}

\subsection{The derivative-level de Branges structure}

Define
\begin{equation}\label{eq:K-rotated-derivative}
 K(t)=e^{it/2}H'(it).
\end{equation}

\begin{proposition}[Rotated derivative]
\label{prop:H-rotated-LP}
One has
\begin{align}
 K(t)
 &=\frac{2\sin\!\left(\frac12\sqrt{t^2+4}\right)}{\sqrt{t^2+4}}
 \label{eq:K-sine}\\
 &=\sin(1)\prod_{n\geq1}
   \left(1-\frac{t^2}{4(\pi^2n^2-1)}\right).
 \label{eq:K-real-product}
\end{align}
Thus $K$ belongs to the Laguerre--P\'olya class and has only the simple real
zeros
\[
 \pm2\sqrt{\pi^2n^2-1}\qquad(n\geq1).
\]
Furthermore
\begin{equation}\label{eq:HB-function}
 E(t)=K(t)+iK'(t)
\end{equation}
is a Hermite--Biehler function and therefore defines a de Branges space.
\end{proposition}

\begin{proof}
Equations \eqref{eq:K-sine} and \eqref{eq:K-real-product} follow from
\eqref{eq:Hprime-sinh} and Euler's product for $\sin u/u$; the constant is
$\prod_{n\geq1}(1-1/(\pi^2n^2))=\sin(1)$.  The partial products in
\eqref{eq:K-real-product} are real-rooted polynomials and converge locally
uniformly, which is exactly the required Laguerre--P\'olya approximation.

For $\Im t>0$, the logarithmic derivative of the product satisfies
\[
 \Im\frac{K'(t)}{K(t)}
 =\sum_{n\geq1}\Im\left(
   \frac1{t-a_n}+\frac1{t+a_n}\right)<0,
 \qquad a_n=2\sqrt{\pi^2n^2-1}.
\]
The paired series converges locally uniformly off the real zero set.  With
$E^\#(t)=\overline{E(\bar t)}=K(t)-iK'(t)$, it follows that
\[
 |E(t)|^2-|E^\#(t)|^2
 =-4|K(t)|^2\Im\frac{K'(t)}{K(t)}>0
 \qquad(\Im t>0).
\]
This is the Hermite--Biehler inequality.
\end{proof}

\begin{remark}[Scope of the real-zero analogy]
The terminology above is classical in the sense of P\'olya--Schur and de
Branges.\footnote{G.~P\'olya and J.~Schur, \emph{Über zwei Arten von
Faktorenfolgen in der Theorie der algebraischen Gleichungen}, J. Reine
Angew. Math. 144 (1914), 89--113; L.~de Branges, \emph{Hilbert Spaces of
Entire Functions}, Prentice--Hall, 1968.}
It applies one derivative below the reciprocal-level problem.  Every
translate $H-2c$ has the same derivative, hence the same $K$, the same
Hermite--Biehler function \eqref{eq:HB-function}, and the same canonical de
Branges space.  This construction therefore forgets the level $c$ and in
particular cannot distinguish $c=\gamma$ from another real level.

Conversely, Theorem~\ref{thm:H-level-asymptotics} gives infinitely many
nonreal zeros of $H-2c$.  For real $c$, the function $H-2c$ is therefore not
in the Laguerre--P\'olya class and cannot be the real part of a
Hermite--Biehler function, whose real and imaginary parts have only real
zeros.  Likewise, although the symmetry $z\leftrightarrow z^{-1}$ resembles
a self-inversive or Lee--Yang symmetry, the off-circle fibres show that no
full circle theorem holds here.  The statement that algebraic zeros at the
Euler level must lie on the unit circle is an arithmetic restriction, not a
Lee--Yang theorem.
\end{remark}

\begin{remark}[Novelty and logical status]
The analytic descent and product formulas above are elementary consequences
of the Joukowski symmetry and Euler's sine product.  The arithmetic
statements are corollaries of Theorem~\ref{thm:master}; their novelty is the
packaging, not a new transcendence mechanism.  The conditional hypothesis
that $\alpha_*$ is algebraic is not known, so
Theorem~\ref{thm:H-euler-level} does not prove the transcendence of
$\gamma$ or $\delta$.  No claim of priority should be made for the
Joukowski packaging before a specialist literature search.
\end{remark}


\providecommand{\Eexp}{\mathcal E_{\mathrm{exp}}}

\section{Spectral stability, arithmetic Clark data, and a second
recovery of \texorpdfstring{$\gamma$}{gamma}}
\label{sec:ssc-spectral-stability-clark}

The Joukowski quotient has a positive Laplace representation which places
its zero divisor strictly in the left half-plane.  After rotation and a
centering exponential, this stability becomes a Hermite--Biehler
inequality.  The resulting meromorphic inner function has two complementary
features: its values at algebraic real points are algebraically independent,
whereas their large-point phase recovers Euler's constant.  We develop these
statements in a form which keeps the unconditional arithmetic assertions
separate from the still open arithmetic of \(\gamma\).

\subsection{A positive Bessel--Laplace representation}

\begin{proposition}[Bessel--Laplace representation]
\label{prop:ssc-bessel-laplace}
For every \(w\in\C\),
\begin{equation}\label{eq:ssc-bessel-laplace}
 H'(w)=\int_0^1 e^{-sw}
 J_0\!\left(2\sqrt{s(1-s)}\right)\,ds.
\end{equation}
The kernel in \eqref{eq:ssc-bessel-laplace} is strictly positive on
\([0,1]\).
\end{proposition}

\begin{proof}
Set \(t=2s-1\).  The classical Sonine integral, obtained for example by
expanding the exponential and the Bessel function and integrating
termwise, is
\[
 \int_{-1}^{1}e^{at}J_0\!\left(b\sqrt{1-t^2}\right)\,dt
 =2\frac{\sinh\sqrt{a^2-b^2}}{\sqrt{a^2-b^2}},
\]
with removable continuation when \(a^2=b^2\).  Taking
\(a=-w/2\) and \(b=1\) gives
\[
 \int_0^1 e^{-sw}J_0\!\left(2\sqrt{s(1-s)}\right)\,ds
 =2e^{-w/2}
   \frac{\sinh(\frac12\sqrt{w^2-4})}{\sqrt{w^2-4}},
\]
which is \eqref{eq:Hprime-sinh}.  Finally
\(0\leq2\sqrt{s(1-s)}\leq1\), and the alternating series for \(J_0\)
shows that \(J_0(x)>0\) for \(0\leq x\leq1\).
\end{proof}

\begin{theorem}[Rightmost point of every real fibre]
\label{thm:ssc-rightmost-fibre}
Let \(a\in\mathbb R\).  If
\[
 H(w)=H(a),\qquad w\ne a,
\]
then
\begin{equation}\label{eq:ssc-rightmost-fibre}
 \Re w<a.
\end{equation}
Thus \(a\) is the unique rightmost point of its full complex fibre.
\end{theorem}

\begin{proof}
Integrating \eqref{eq:ssc-bessel-laplace} from \(a\) to \(w\) gives
\begin{equation}\label{eq:ssc-fibre-integral}
 H(w)-H(a)=\int_0^1 e^{-as}
 \frac{1-e^{-s(w-a)}}{s}
 J_0\!\left(2\sqrt{s(1-s)}\right)\,ds,
\end{equation}
where the integrand is continued at \(s=0\).  Put \(d=w-a\).  If
\(\Re d\geq0\), then
\[
 \Re\bigl(1-e^{-sd}\bigr)
 =1-e^{-s\Re d}\cos(s\Im d)
 \geq1-e^{-s\Re d}\geq0.
\]
All remaining factors in \eqref{eq:ssc-fibre-integral} are positive.
If the integral vanishes, its real part therefore vanishes identically.
This first forces \(\Re d=0\), and then
\(1-\cos(s\Im d)=0\) for every \(s\in[0,1]\), which forces
\(\Im d=0\).  Hence \(d=0\), contrary to the hypothesis.  Thus
\(\Re d<0\).
\end{proof}

\begin{corollary}[Hurwitz stability of the level-zero divisors]
\label{cor:ssc-hurwitz-stability}
Let \(\omega_0\) denote the unique real zero of \(H\).  Then
\(\omega_0<0\), and every other zero \(\rho\) of \(H\) satisfies
\[
 \Re\rho<\omega_0<0.
\]
All zeros of \(H\) are simple and transcendental.  Moreover every zero
\(z\) of
\[
 \Psi(z)=H(z+z^{-1})
\]
is simple, transcendental, and lies in the open left half-plane.
\end{corollary}

\begin{proof}
The real diffeomorphism in Corollary~\ref{cor:H-critical-lattice}, together
with \(H(0)=2\Cin(1)>0\), gives the unique real zero \(\omega_0<0\).
Theorem~\ref{thm:ssc-rightmost-fibre} gives the asserted strict inequality
for every other zero.  The critical points of \(H\) are the purely
imaginary points \eqref{eq:H-critical-points}, so none is a zero; the real
zero is simple because \(H'(\omega_0)>0\).  Transcendence follows from
Corollary~\ref{cor:H-algebraic-levels} at the algebraic level zero.

If \(\Psi(z)=0\), put \(\rho=J(z)\).  Since
\[
 \Re J(z)=\Re z\,(1+|z|^{-2}),
\]
the inequality \(\Re\rho<0\) implies \(\Re z<0\).  The points
\(z=\pm1\), at which \(J'\) vanishes, are not zeros because
\(H(2)>0\) and \(H(-2)<0\).  Hence simplicity follows from the chain
rule.  Finally an algebraic zero \(z\) would make \(J(z)\) an algebraic
zero of \(H\), which has just been excluded.
\end{proof}

\subsection{The Euler point as an explicit
\texorpdfstring{$E$--$D$}{E--D} collision}

Recall that \(w_*\in(-2,2)\) is the unique real solution of
\(H(w)=2\gamma\), and let \(\alpha_*\) and \(\alpha_*^{-1}\) be its two
preimages under \(J\).  Define
\begin{align}
 \mathcal G(u)
 &:=eH'(2+u)
   =e^{-u/2}\sum_{m\geq0}
     \frac{\bigl(u(u+4)/4\bigr)^m}{(2m+1)!},
     \label{eq:ssc-G-definition}\\
 \mathcal K(u)
 &:=-\frac12\int_0^u\mathcal G(t)\,dt.
     \label{eq:ssc-K-definition}
\end{align}
Both functions are entire, $D$-finite, and have rational Taylor
coefficients.  Direct integration gives the useful identity
\begin{equation}\label{eq:ssc-K-H-identity}
 \mathcal K(u)=\frac e2\bigl(H(2)-H(2+u)\bigr).
\end{equation}

There is a tempting but false shortcut here: rational Taylor coefficients
and holonomicity do not make \(\mathcal K\) a Siegel \(E\)-function.

\begin{proposition}[The auxiliary primitive is not an \(E\)-function]
\label{prop:ssc-K-not-E}
Neither \(\mathcal G\) nor \(\mathcal K\) is a Siegel \(E\)-function.
Nevertheless, if \(u\in\Qbar\), then \(\mathcal K(u)\) belongs to the ring
\(\mathbf E\) of finite \(E\)-values.
\end{proposition}

\begin{proof}
Writing \(y=\mathcal G\), the differential equation
\eqref{eq:H-third-order-ode} shifted by \(w=2+u\) gives
\begin{equation}\label{eq:ssc-G-minimal-operator}
 u(u+2)(u+4)y''
 +(u^3+8u^2+16u+12)y'
 +(u^2+3u+4)y=0.
\end{equation}
The function \(y=eH'(2+u)\) has infinitely many zeros, by
\eqref{eq:H-critical-points}.  It therefore cannot satisfy a first-order
homogeneous equation over \(\Qbar(u)\): a nonzero entire solution of such an
equation can have zeros only among the finitely many poles of its rational
logarithmic derivative.  Thus the order-two operator in
\eqref{eq:ssc-G-minimal-operator} is minimal.  At \(u=-4\) its indicial
polynomial is
\[
 8r(r-1)+12r=4r(2r+1),
\]
so its local exponents are \(0\) and \(-1/2\).  Hence \(-4\) is a
non-apparent nonzero finite singularity.  The finite-singularity theorem for
minimal annihilators of \(E\)-functions says that every nonzero finite
singularity is apparent \cite{Andre2000}.  Consequently
\(\mathcal G\) is not an \(E\)-function.  Since \(E\)-functions are closed
under differentiation and \(\mathcal K'=-\mathcal G/2\), neither is
\(\mathcal K\).

For the last assertion, choose algebraic \(a\ne0\) satisfying
\(a+a^{-1}=2+u\).  Equations \eqref{eq:H-descent} and
\eqref{eq:ssc-K-H-identity} give
\begin{equation}\label{eq:ssc-K-finite-E-values}
 \mathcal K(u)=\frac e2
 \bigl(2\Ein(1)-\Ein(a)-\Ein(a^{-1})\bigr).
\end{equation}
The right side belongs to \(\mathbf E\), since \(e\) and the three displayed
finite algebraic-point values are \(E\)-values and \(\mathbf E\) is a ring.
\end{proof}

\begin{proposition}[Equivalent algebraic witnesses]
\label{prop:ssc-euler-equivalences}
The following assertions are equivalent:
\begin{enumerate}[label=\textup{(\roman*)}]
 \item \(w_*\in\Qbar\);
 \item \(\alpha_*\in\Qbar\);
 \item there exists
       \(a\in\Qbar^\times\setminus(-\infty,0]\) such that, on the
       principal branch,
       \[
        E_1(a)+E_1(a^{-1})=0;
       \]
 \item there exists \(u\in\Qbar\cap(-4,0)\) such that
       \[
        \mathcal K(u)=\delta.
       \]
\end{enumerate}
If these conditions hold, then \(\gamma\) and \(\delta\) are
algebraically independent.
\end{proposition}

\begin{proof}
The equivalence of (i) and (ii) follows from
\(\alpha_*^2-w_*\alpha_*+1=0\) and
\(w_*=\alpha_*+\alpha_*^{-1}\).  For a point in the domain specified in
(iii), the inverse-compatible principal branches satisfy
\[
 E_1(z)+E_1(z^{-1})=H(J(z))-2\gamma
\]
and this identity shows that (ii) implies (iii).  Conversely an algebraic
solution in (iii)
gives an algebraic zero \(J(a)\) of \(H-2\gamma\).  By
Theorem~\ref{thm:H-euler-level}, every algebraic zero of this Euler-level
translate must equal \(w_*\), proving (i); the same theorem then forces
\(a\in\{\alpha_*,\alpha_*^{-1}\}\).

Put \(u_*=w_*-2\in(-4,0)\).  Since
\[
 H(2)=2\Ein(1)=2\gamma+\frac{2\delta}{e},
\]
identity \eqref{eq:ssc-K-H-identity} gives
\(\mathcal K(u_*)=\delta\), so (i) implies (iv).  Conversely, if
\(u\in(-4,0)\) satisfies \(\mathcal K(u)=\delta\), the same identity
gives \(H(2+u)=2\gamma\).  Uniqueness on the real axis forces
\(2+u=w_*\), proving (i).  The final assertion is the witness
amplification statement in Theorem~\ref{thm:H-euler-level}.
\end{proof}

\begin{corollary}[Exact \(E\)--\(D\) consequence and the Euler dichotomy]
\label{cor:ssc-ED-dichotomy}
Let \(\mathbf E\) and \(\mathbf D\) denote the rings of \(E\)-values and
arithmetic-Gevrey values of order one, respectively, in the sense of
Fischler--Rivoal \cite{FR2024ED}.  If \(w_*\) is algebraic, then
\[
 \delta\in(\mathbf E\cap\mathbf D)\setminus\Qbar.
\]
Consequently the conjectural equality
\(\mathbf E\cap\mathbf D=\Qbar\) implies that \(w_*\) is
transcendental.  Unconditionally, exactly one of the following alternatives
holds:
\begin{enumerate}[label=\textup{(\alph*)}]
 \item every zero of \(H-2\gamma\), and every zero of the reciprocal
       Euler equation, is transcendental;
 \item the Euler level has the unique algebraic reciprocal pair
       \(\{\alpha_*,\alpha_*^{-1}\}\), and \(\gamma,\delta\) are
       algebraically independent.
\end{enumerate}
\end{corollary}

\begin{proof}
Under the algebraicity hypothesis,
Proposition~\ref{prop:ssc-euler-equivalences} writes \(\delta\) as
\(\mathcal K(u_*)\) at an algebraic point.  Formula
\eqref{eq:ssc-K-finite-E-values}, rather than the false assertion that
\(\mathcal K\) itself is an \(E\)-function, gives
\(\delta\in\mathbf E\).  The inclusion
\(\delta\in\mathbf D\) is the standard Euler-series realization of the
Gompertz constant, and the algebraic independence conclusion in
Proposition~\ref{prop:ssc-euler-equivalences} shows that \(\delta\) is
transcendental.  This proves the first assertion.

Theorem~\ref{thm:H-euler-level} says that \(w_*\) is the only possible
algebraic zero at the Euler level.  If it is transcendental, an algebraic
reciprocal zero is impossible because its Joukowski image would be
algebraic.  If it is algebraic, its two Joukowski preimages are algebraic,
they are the only algebraic reciprocal zeros, and witness amplification
gives the stated algebraic independence.
\end{proof}

\subsection{A level-zero de Branges--Clark model}

For an entire function \(E\), write
\(E^\#(z)=\overline{E(\bar z)}\).  Define
\begin{equation}\label{eq:ssc-E0-Theta}
 E_0(z)=e^{-iz/2}H(-iz),\qquad
 \Theta(z)=\frac{E_0^\#(z)}{E_0(z)}
          =e^{iz}\frac{H(iz)}{H(-iz)}.
\end{equation}

\begin{theorem}[Hermite--Biehler and Clark structure]
\label{thm:ssc-HB-Clark}
The function \(E_0\) is a Cartwright-class Hermite--Biehler function of
exponential type \(1/2\).  Consequently \(\Theta\) is a meromorphic inner
function in the upper half-plane and has zero mean type.  The zeros of
\(\Theta\) are the points \(-i\rho\), where \(\rho\) ranges over the zeros
of \(H\).
Thus \(E_0\) defines a de Branges space, and for each
\(\alpha\in\mathbb T\) the real solutions of
\begin{equation}\label{eq:ssc-Clark-level}
 \Theta(x)=\alpha
\end{equation}
form the corresponding discrete Clark, equivalently self-adjoint, spectrum.
\end{theorem}

\begin{proof}
Corollary~\ref{cor:ssc-hurwitz-stability} places every zero \(\rho\) of
\(H\) in the open left half-plane.  Hence the zeros \(i\rho\) of \(E_0\)
lie in the lower half-plane, and \(E_0\) has no real zeros.  The connection
formula and the standard sectorial expansion of \(E_1\) give, as
\(x\to+\infty\),
\begin{equation}\label{eq:ssc-H-imaginary-asymptotic}
 H(ix)=\gamma+\log x+\frac{\pi i}{2}+O(x^{-1}),
 \qquad H(-ix)=\overline{H(ix)}.
\end{equation}
The two-exponential formula \eqref{eq:Hprime-sinh}, followed by one
integration, gives the indicator
\[
 h_{E_0}(\vartheta)=\frac12|\sin\vartheta|.
\]
In particular \(E_0\) has exponential type \(1/2\).  Equation
\eqref{eq:ssc-H-imaginary-asymptotic} shows that
\(|E_0(x)|=O(\log(2+|x|))\) on the real axis, so \(E_0\) belongs to the
Cartwright class.

It follows that \(\Theta=E_0^\#/E_0\) is analytic and of bounded type in
the upper half-plane, with unimodular boundary values.  For
\(y\to+\infty\), the real-axis asymptotics obtained from the Joukowski
representation are
\[
 H(y)=\gamma+\log y+O(y^{-1}),\qquad
 H(-y)=-\frac{e^y}{y}\bigl(1+O(y^{-1})\bigr).
\]
Therefore
\begin{equation}\label{eq:ssc-Theta-mean-type}
 \Theta(iy)=e^{-y}\frac{H(-y)}{H(y)}
 =-\frac{1+o(1)}{y\log y}.
\end{equation}
Thus \(\Theta\) has zero mean type.  Bounded-type factorization, together
with its unimodular boundary values, now shows that \(\Theta\) is inner:
meromorphic continuation across the real axis removes the singular-inner
measure, while zero mean type removes the exponential factor in the
canonical factorization \cite{deBranges1968,Clark1972}.  It is nonconstant,
so the Hermite--Biehler inequality is strict.
Meromorphic continuation follows from \eqref{eq:ssc-E0-Theta}.  The
description of its zeros is immediate from \(H(iz)=0\), and the final
assertion is the standard meromorphic-inner Clark/de Branges
correspondence \cite{Clark1972,deBranges1968}.
\end{proof}

\begin{theorem}[Arithmetic independence of the Clark data]
\label{thm:ssc-arithmetic-Clark}
Let \(x_1,\ldots,x_m\in\Qbar\cap\mathbb R^\times\), and suppose that
\(x_j^2\ne x_k^2\) for \(j\ne k\).  Then the \(m\) numbers
\[
 R(x_j):=\frac{H(ix_j)}{H(-ix_j)}
\]
are algebraically independent over \(\Eexp\).  Separately, the \(m\)
numbers
\[
 \Theta(x_j)=e^{ix_j}R(x_j)
\]
are algebraically independent over \(\Eexp\).  In particular each of
\[
 \{R(n):n\geq1\},\qquad \{\Theta(n):n\geq1\}
\]
is a countably algebraically independent family.

If \(\alpha\in\Qbar\cap\mathbb T\), every real solution of
\(\Theta(x)=\alpha\) is transcendental, apart from the unavoidable
solution \(x=0\) when \(\alpha=1\).
\end{theorem}

\begin{proof}
The \(2m\) algebraic points \(ix_j,-ix_j\) are distinct.  Their Joukowski
fibres are pairwise disjoint.  Theorem~\ref{thm:pure-exp-base}, followed by
the same disjoint-support linear projection used in
Theorem~\ref{thm:H-multipoint-descent}, shows that
\[
 H(ix_1),H(-ix_1),\ldots,H(ix_m),H(-ix_m)
\]
are algebraically independent over \(\Eexp\).  Ratios made from disjoint
pairs of independent indeterminates remain algebraically independent:
indeed, a polynomial relation among the ratios becomes, after clearing
the denominator monomial, a nonzero polynomial relation among the original
indeterminates.  Multiplication by the nonzero coefficients
\(e^{ix_j}\in\Eexp\) preserves independence.

If a nonzero algebraic real \(x\) satisfied \(\Theta(x)=\alpha\), then
\[
 e^{ix}H(ix)-\alpha H(-ix)=0
\]
would be a nontrivial linear relation over \(\Eexp\) between two
independent values.  This is impossible.  At \(x=0\), one has
\(\Theta(0)=1\), which explains the unique exception.
\end{proof}

\begin{remark}[Why signs must not be duplicated]
The restriction \(x_j^2\ne x_k^2\) is necessary: on the real axis
\[
 \Theta(-x)=\Theta(x)^{-1},\qquad R(-x)=R(x)^{-1}.
\]
Thus a family containing both \(x\) and \(-x\) has an immediate algebraic
relation.
\end{remark}

\subsection{Recovery of \texorpdfstring{$\gamma$}{gamma} from arithmetic
scattering data}

For \(n\geq1\), put
\begin{equation}\label{eq:ssc-an-Rn-Tn}
 a_n=\frac{n+\sqrt{n^2+4}}2,\qquad
 R_n=R(n)=\frac{H(in)}{H(-in)}=e^{-in}\Theta(n),\qquad
 T_n=\frac{\pi i}{2}\frac{R_n+1}{R_n-1}.
\end{equation}
Theorem~\ref{thm:ssc-arithmetic-Clark} implies \(R_n\ne1\), so \(T_n\)
is well defined.  Since \(|R_n|=1\), it is real.

\begin{theorem}[Scattering recovery and its first correction]
\label{thm:ssc-gamma-scattering}
One has
\begin{equation}\label{eq:ssc-gamma-raw-recovery}
 \gamma=\lim_{n\to\infty}\bigl(T_n-\log n\bigr),
 \qquad
 T_n-\log n
 =\gamma+O\!\left(\frac{\log n}{n}\right).
\end{equation}
More precisely, the completely explicit corrected estimator
\begin{equation}\label{eq:ssc-gamma-corrected-estimator}
 \widehat\gamma_n=
 \left(1-\frac{2(1+\cos a_n)}{\pi a_n}\right)T_n
 +\frac{\sin a_n}{a_n}-\log a_n
\end{equation}
satisfies
\begin{equation}\label{eq:ssc-gamma-corrected-rate}
 \widehat\gamma_n
 =\gamma+O\!\left(\frac{\log n}{n^2}\right).
\end{equation}
Thus \(\gamma\) is recovered from ratios of values generated by
\(\Ein\) at moving algebraic points, even though every finite family of
those ratios is algebraically independent over \(\Eexp\).
\end{theorem}

\begin{proof}
The identity \(a_n-a_n^{-1}=n\) gives \(J(ia_n)=in\), hence
\begin{equation}\label{eq:ssc-Hin-Ein}
 H(in)=\Ein(ia_n)+\Ein(-i/a_n).
\end{equation}
Put \(L_n=\gamma+\log a_n\) and \(b=\pi/2\).  On the positive imaginary
axis, the connection formula and the standard asymptotic expansion
\[
 E_1(ia)=\frac{e^{-ia}}{ia}
 \left(1-\frac1{ia}+O(a^{-2})\right)
\]
combine with
\(\Ein(-i/a)=-i/a+1/(4a^2)+O(a^{-3})\).  Separating real and imaginary
parts in \eqref{eq:ssc-Hin-Ein} gives
\begin{align}
 \Re H(in)
 &=L_n-\frac{\sin a_n}{a_n}
   +\frac{\cos a_n+1/4}{a_n^2}+O(a_n^{-3}),
   \label{eq:ssc-Hin-real-expansion}\\
 \Im H(in)
 &=b-\frac{1+\cos a_n}{a_n}
   -\frac{\sin a_n}{a_n^2}+O(a_n^{-3}).
   \label{eq:ssc-Hin-imag-expansion}
\end{align}
Because \(H(-in)=\overline{H(in)}\), direct algebra gives
\begin{equation}\label{eq:ssc-Tn-real-ratio}
 T_n=b\frac{\Re H(in)}{\Im H(in)}.
\end{equation}
Equations \eqref{eq:ssc-Hin-real-expansion}--\eqref{eq:ssc-Tn-real-ratio}
first yield
\(T_n=L_n+O(L_n/a_n)\).  Since
\(a_n=n+O(n^{-1})\), this proves
\eqref{eq:ssc-gamma-raw-recovery}.

For the sharper statement, multiply \eqref{eq:ssc-Tn-real-ratio} by
\[
 1-\frac{1+\cos a_n}{b a_n}
\]
and add \(\sin(a_n)/a_n\).  The terms of order \(a_n^{-1}\) in both the
numerator and denominator cancel, leaving
\[
 \left(1-\frac{1+\cos a_n}{b a_n}\right)T_n
 +\frac{\sin a_n}{a_n}
 =L_n+O\!\left(\frac{L_n}{a_n^2}\right).
\]
Since \(b=\pi/2\), subtraction of \(\log a_n\) gives
\eqref{eq:ssc-gamma-corrected-estimator} and
\eqref{eq:ssc-gamma-corrected-rate}.
\end{proof}

\subsection{An exact inverse-resonance trace}

\begin{theorem}[Inverse-resonance trace]
\label{thm:ssc-inverse-resonance}
Let \(Z(H)\) denote the multiset of all zeros of \(H\).  Then
\begin{equation}\label{eq:ssc-inverse-resonance}
 \boxed{
 \sum_{\rho\in Z(H)}\frac{-2\Re\rho}{|\rho|^2}
 =1+\frac{\sin1}{\Cin(1)}.}
\end{equation}
The series converges absolutely, every summand is positive, and its value
is transcendental over \(\Eexp\), hence also over
\(\Qbar(e,e^i)\).
\end{theorem}

\begin{proof}
By Theorem~\ref{thm:ssc-HB-Clark}, \(\Theta\) is a meromorphic inner
function of zero mean type.  Meromorphic continuation across the real axis
excludes a finite singular measure in its inner factorization, and zero
mean type excludes an exponential inner factor.  Thus \(\Theta\) is a
pure Blaschke product.  Its zeros are
\[
 \zeta_\rho=-i\rho\in\mathbb C_+
 \qquad(\rho\in Z(H)).
\]
Logarithmic differentiation of the Blaschke product at zero gives
\begin{equation}\label{eq:ssc-Blaschke-log-derivative}
 \frac1i\frac{\Theta'(0)}{\Theta(0)}
 =\sum_{\rho\in Z(H)}
   \frac{2\Im\zeta_\rho}{|\zeta_\rho|^2}
 =\sum_{\rho\in Z(H)}
   \frac{-2\Re\rho}{|\rho|^2}.
\end{equation}
The Blaschke condition gives absolute convergence; equivalently, all terms
are positive and there are no zeros near the origin.

On the other hand, \(\Theta(0)=1\),
\(H(0)=2\Cin(1)\), and \(H'(0)=\sin1\).  Differentiating
\eqref{eq:ssc-E0-Theta} directly yields
\begin{equation}\label{eq:ssc-origin-phase-derivative}
 \frac1i\frac{\Theta'(0)}{\Theta(0)}
 =1+2\frac{H'(0)}{H(0)}
 =1+\frac{\sin1}{\Cin(1)}.
\end{equation}
Together with \eqref{eq:ssc-Blaschke-log-derivative}, this proves the
trace identity.

Finally \(\sin1\in\Eexp^\times\), whereas
\[
 \Cin(1)=\frac{\Ein(i)+\Ein(-i)}2
\]
is transcendental over \(\Eexp\) by
Theorem~\ref{thm:pure-exp-base}.  Hence the right side of
\eqref{eq:ssc-inverse-resonance} is transcendental over \(\Eexp\).
\end{proof}

\begin{remark}[Logical scope]
The stability and Clark results above are unconditional consequences of
the Joukowski descent and the multipoint \(\Ein\) theorem.  The recovery
formula \eqref{eq:ssc-gamma-raw-recovery} is an asymptotic encoding of
\(\gamma\), not an irrationality criterion: passing from algebraic
independence at each finite set of moving points to an arithmetic statement
about the limit requires a genuinely uniform measure.  Likewise
Proposition~\ref{prop:ssc-euler-equivalences} isolates an exact algebraic
witness but does not establish that the witness exists.
\end{remark}


\section{Reciprocal traces, phase recovery, and height saturation}
\label{sec:reciprocal-trace-height}

This section separates two reciprocal functions which have different
arithmetic meanings.  Put
\begin{equation}\label{eq:Psi-and-Fgamma}
 \Psi(z)=\Ein(z)+\Ein(z^{-1})=H(J(z)),
 \qquad
 \mathcal F_\gamma(z)=E_1(z)+E_1(z^{-1})=\Psi(z)-2\gamma,
\end{equation}
where branches are chosen coherently so that
$\Log(z^{-1})=-\Log z$.  The first function has rational Laurent
coefficients and is genuinely independent of $\gamma$.  The second is the
Euler-level translate of the first; the opposite monodromies of its two
$E_1$ terms cancel, so the displayed sum denotes its single-valued analytic
continuation on $\C^\times$.

This distinction is also essential for effectivity.  The entire function
$\Ein$ is a Siegel $E$-function, whereas $E_1$ has a logarithmic
singularity at the origin and is not an $E$-function.  In particular,
an algebraic approximation to a zero of $\mathcal F_\gamma$ does not put all
terms of
\[
 \Log z+E_1(z)+\Ein(z^{-1})
\]
inside the effective theory of $E$-values.  Indeed the connection formula
gives the identity
\begin{equation}\label{eq:tautological-recovery}
 \Log z+E_1(z)+\Ein(z^{-1})
 =\Psi(z)-\gamma=\gamma+\mathcal F_\gamma(z).
\end{equation}
Thus evaluation at a zero of $\mathcal F_\gamma$ recovers $\gamma$, but
only tautologically.  The nontrivial recovery below uses instead the zero
divisor and the asymptotic phase of reciprocal functions.

\subsection{An exact Hadamard trace at the Euler level}

Let $\{u_k\}$ be the zeros of $H(u)-2\gamma$, repeated according to
multiplicity.  The index set includes the unique real zero from
Theorem~\ref{thm:H-euler-level} and both members of every nonreal conjugate
pair.

\begin{theorem}[Euler-level Hadamard trace]
\label{thm:euler-hadamard-trace}
The series $\sum_k |u_k|^{-2}$ converges and
\begin{equation}\label{eq:euler-hadamard-trace}
 \boxed{
 \sum_k\frac{1}{u_k^2}
 =\frac{\sin^2 1-\sin 1\,\Ci(1)}{4\Ci(1)^2}.}
\end{equation}
Consequently, if $T$ denotes the left side, then
\begin{equation}\label{eq:trace-inversion}
 \Ci(1)=
 \frac{\sin 1\bigl(\sqrt{1+16T}-1\bigr)}{8T},
 \qquad
 \gamma=\Cin(1)+
 \frac{\sin 1\bigl(\sqrt{1+16T}-1\bigr)}{8T}.
\end{equation}
\end{theorem}

\begin{proof}
Set $f(u)=H(u)-2\gamma$.  By
Corollary~\ref{cor:H-order-zero-count}, $f$ is entire of order one.  Hence
its zero divisor has exponent of convergence at most one and admits a
genus-one Hadamard factorization.  Since
\[
 f(0)=2\Cin(1)-2\gamma=-2\Ci(1)\ne0,
\]
logarithmic differentiation twice at the origin gives
\begin{equation}\label{eq:hadamard-log-derivative}
 -\left(\frac{f'}f\right)'(0)=\sum_k\frac1{u_k^2}.
\end{equation}
The product formula for $H'$ in
Proposition~\ref{prop:entire-joukowski-descent} gives the exact origin jet
\[
 H'(0)=\sin 1,
 \qquad
 H''(0)=-\frac12\sin 1.
\]
Therefore
\[
 -\left(\frac{f'}f\right)'(0)
 =\frac{f'(0)^2-f''(0)f(0)}{f(0)^2}
 =\frac{\sin^2 1-\sin 1\,\Ci(1)}{4\Ci(1)^2},
\]
which proves \eqref{eq:euler-hadamard-trace}.  The standard bounds
$0<\Ci(1)<\sin 1$ (which also follow directly from the corresponding
alternating series and elementary bounds for $\gamma$) show that the trace
$T$ is positive.  Solving
$4T X^2+\sin(1)X-\sin^2(1)=0$ for the positive root
$X=\Ci(1)$ gives \eqref{eq:trace-inversion}.
\end{proof}

\begin{remark}[What the trace does not prove]
Formula \eqref{eq:trace-inversion} is an exact spectral reparametrization,
not a transcendence proof: the divisor on the left is itself the divisor of
the $\gamma$-dependent level $H-2\gamma$.  No arithmetic assertion in this
paper relies on a numerical zero count.  In particular, an exact floor
formula suggested by finite computations is not used; the rigorous bounds
remain those of Corollary~\ref{cor:H-order-zero-count}.
\end{remark}

\subsection{A nontrivial phase recovery from the
\texorpdfstring{$\gamma$}{gamma}-free zero set}

The situation is different for zeros of $\Psi$, because the defining
Laurent series then contains no $\gamma$.  Let $z_{k,0}$ be the upper
half-plane zero with $|z_{k,0}|>1$ constructed in the proof of
Theorem~\ref{thm:H-level-asymptotics} at level $c=0$, and write
$\rho_k=z_{k,0}$.  Put
\[
 t_k=2\pi k,
 \qquad q_k=\log t_k.
\]

\begin{theorem}[Second-order zero phase]
\label{thm:gamma-free-zero-phase}
As $k\to\infty$,
\begin{align}
 \rho_k={}&-q_k-\frac12\log\left((q_k+\gamma)^2+\frac{\pi^2}{4}\right)
 \label{eq:rho-second-order}\\
 &\quad+i\left[t_k+\arctan\left(\frac{2(q_k+\gamma)}{\pi}\right)\right]
 +O\left(\frac{q_k}{t_k}\right).\nonumber
\end{align}
In particular the zero set of the $\gamma$-free function $\Psi$ determines
$\gamma$ through
\begin{equation}\label{eq:gamma-phase-recovery}
 \boxed{
 \gamma=\lim_{k\to\infty}
 \left\{
  \frac\pi2\tan\bigl(\Im\rho_k-2\pi k\bigr)-\log(2\pi k)
 \right\}.}
\end{equation}
More precisely, the expression in braces is
\begin{equation}\label{eq:gamma-phase-rate}
 \gamma+O\!\left(\frac{(\log k)^3}{k}\right).
\end{equation}
Every $\rho_k$ is transcendental.
\end{theorem}

\begin{proof}
At level $c=0$, the fixed-point equation in the proof of
Theorem~\ref{thm:H-level-asymptotics} is
\begin{equation}\label{eq:refined-fixed-point}
 z=(2k+1)\pi i-\Log z-\Log(\Log z+\gamma).
\end{equation}
The exact zero $\rho_k$ differs from its model fixed point by $O(t_k^{-1})$.
Uniformly in the unit disc used there,
\[
 \Log z=q_k+\frac{\pi i}{2}+O(q_k/t_k),
\]
and hence
\begin{align*}
 \Log(\Log z+\gamma)
  ={}&\frac12\log\left((q_k+\gamma)^2+\frac{\pi^2}{4}\right)\\
   &+i\arctan\left(\frac{\pi}{2(q_k+\gamma)}\right)
     +O(q_k/t_k).
\end{align*}
Substitution in \eqref{eq:refined-fixed-point}, followed by
\[
 \frac\pi2-\arctan\left(\frac{\pi}{2x}\right)
 =\arctan\left(\frac{2x}{\pi}\right)
 \qquad(x>0),
\]
proves \eqref{eq:rho-second-order}.

Write
$\Im\rho_k-t_k=\arctan(2(q_k+\gamma)/\pi)+\varepsilon_k$, where
$\varepsilon_k=O(q_k/t_k)$.  Since $q_k/t_k=o(q_k^{-1})$, the whole
mean-value segment remains a distance asymptotic to $q_k^{-1}$ from the
nearest pole of $\tan$.  Consequently $\sec^2=O(q_k^2)$ uniformly there,
\[
 \frac\pi2\tan(\Im\rho_k-t_k)
 =q_k+\gamma+O(q_k^3/t_k).
\]
The error tends to zero and proves both
\eqref{eq:gamma-phase-recovery} and \eqref{eq:gamma-phase-rate}.
Finally $\Psi(\rho_k)=0$, so $\rho_k$ lies over the algebraic level
$c=0$; its transcendence follows from
Corollary~\ref{cor:H-algebraic-levels}.
\end{proof}

Theorem~\ref{thm:gamma-free-zero-phase} is analytically nonvacuous: unlike
\eqref{eq:tautological-recovery}, the object whose zeros are used has
rational Laurent coefficients.  It still gives no irrationality result,
because no effective arithmetic control of these transcendental zeros is
established here.

\subsection{Laurent truncations saturate the height budget}

Define
\begin{equation}\label{eq:Psi-truncation}
 \Psi_N(z)=\sum_{m=1}^N
 \frac{(-1)^{m+1}}{m\,m!}(z^m+z^{-m}),
 \qquad
 L_N=\operatorname{lcm}(1,\ldots,N),
\end{equation}
and
\begin{equation}\label{eq:QN-definition}
 D_N^*=\operatorname{lcm}_{1\le m\le N}(m\,m!),
 \qquad
 Q_N(z)=D_N^*z^N\Psi_N(z).
\end{equation}

\begin{lemma}[The denominator scale is intrinsic]
\label{lem:truncation-denominator-scale}
One has
\[
 Q_N\in\Z[z],\qquad Q_N\ \text{primitive},\qquad \deg Q_N=2N,
 \qquad H(Q_N)=D_N^*,
\]
and therefore
\begin{equation}\label{eq:QN-height}
 \log H(Q_N)=N\log N+o(N).
\end{equation}
Moreover,
\begin{equation}\label{eq:minimal-denominator-sandwich}
 N!\prod_{p\le N}p\ \mid\ D_N^*\ \mid\ N!L_N,
\end{equation}
so $\log D_N^*=N\log N+o(N)$ as well.
\end{lemma}

\begin{proof}
The coefficients in \eqref{eq:QN-definition} are
$D_N^*/(m\,m!)$, hence are integral and have maximum absolute value
$D_N^*$, attained at $m=1$.  Their gcd is one: otherwise division by a
common prime would give a smaller common denominator than $D_N^*$.
This proves the first assertions.

For the right divisibility in \eqref{eq:minimal-denominator-sandwich}, use
$m!\mid N!$ and $m\mid L_N$.  For the left divisibility, fix a prime
$p\le N$ and take $m=p\lfloor N/p\rfloor$.  There is no multiple of $p$
strictly between $m$ and $N$, so
$v_p(m!)=v_p(N!)$ and
$v_p(m\,m!)\ge v_p(N!)+1$.  Taking the least common multiple over $m$
adds at least one copy of every prime $p\le N$ to $N!$.  The prime number
theorem and Stirling's formula give \eqref{eq:QN-height}.
\end{proof}

\begin{theorem}[Factorial approximation and defining-height saturation]
\label{thm:truncation-height-saturation}
Let $\rho$ be a fixed simple zero of $\Psi$ with $|\rho|>1$.  For all
sufficiently large $N$, $\Psi_N$ has a zero $\alpha_N$ converging to
$\rho$, and
\begin{equation}\label{eq:truncation-root-asymptotic}
 \boxed{
 \alpha_N-\rho=
 \frac{(-1)^{N+2}\rho^{N+1}}
 {\Psi'(\rho)(N+1)(N+1)!}
 \left(1+O_\rho(N^{-1})\right).}
\end{equation}
Consequently
\begin{equation}\label{eq:defining-height-ratio}
 \frac{-\log|\alpha_N-\rho|}{\log H(Q_N)}\longrightarrow1.
\end{equation}
Here $Q_N$ is the primitive integer polynomial obtained by minimal
denominator clearing; no claim about the minimal-polynomial height of
$\alpha_N$ is made.
\end{theorem}

\begin{proof}
On a fixed neighbourhood of $\rho$, the tail
$R_N=\Psi-\Psi_N$ satisfies
\begin{equation}\label{eq:Psi-tail-leading}
 R_N(z)=
 \frac{(-1)^{N+2}z^{N+1}}{(N+1)(N+1)!}
 \left(1+O_\rho(N^{-1})\right).
\end{equation}
Indeed the ratio of consecutive exterior terms is $O_\rho(N^{-1})$, while
the terms in $z^{-m}$ are exponentially smaller.  Rouch\'e's theorem on a
small circle about the simple zero $\rho$ gives a unique nearby zero
$\alpha_N$ of $\Psi_N$.  The standard perturbation formula yields
\[
 \alpha_N-\rho
 =\frac{R_N(\rho)}{\Psi'(\rho)}
  \left(1+O_\rho(N^{-1})\right),
\]
which is \eqref{eq:truncation-root-asymptotic}.  Stirling's formula and
Lemma~\ref{lem:truncation-denominator-scale} give
\[
 -\log|\alpha_N-\rho|=N\log N+O_\rho(N),
 \qquad
 \log H(Q_N)=N\log N+o(N),
\]
and hence \eqref{eq:defining-height-ratio}.
\end{proof}

\begin{corollary}[Full primitive integerization destroys the small value]
\label{cor:integerized-truncation-growth}
Under the hypotheses of
Theorem~\ref{thm:truncation-height-saturation},
\begin{equation}\label{eq:cleared-value-asymptotic}
 \boxed{
 Q_N(\rho)=
 (-1)^{N+3}\frac{D_N^*}{(N+1)(N+1)!}\rho^{2N+1}
 \left(1+O_\rho(N^{-1})\right).}
\end{equation}
In particular
\begin{equation}\label{eq:cleared-value-growth}
 \log|Q_N(\rho)|
 =N\bigl(1+2\log|\rho|\bigr)+o(N)>0.
\end{equation}
Also
\begin{equation}\label{eq:moving-point-value-scale}
 |\Psi(\alpha_N)|=H(Q_N)^{-1+o(1)}.
\end{equation}
\end{corollary}

\begin{proof}
Since $\Psi(\rho)=0$, one has
$Q_N(\rho)=-D_N^*\rho^N R_N(\rho)$.  Substitution of
\eqref{eq:Psi-tail-leading} gives \eqref{eq:cleared-value-asymptotic}.
Lemma~\ref{lem:truncation-denominator-scale} and Stirling's formula then
give \eqref{eq:cleared-value-growth}.  Finally
$\Psi_N(\alpha_N)=0$, so the same tail estimate at $\alpha_N$ and
\eqref{eq:QN-height} give \eqref{eq:moving-point-value-scale}.
\end{proof}

The last corollary closes the naive use of this full primitive clearing
polynomial.  Taken alone, it does not exclude systematic factorization or a
minimal-height compression for the selected root; the dyadic theorem in
Section~\ref{sec:dyadic-capacity} closes precisely that escape on
\(N=2^r\).  The factorial analytic
error is real, but after denominator clearing the full primitive polynomial
value grows rather than tends to zero.  It also explains
why the pointwise measure of Adamczewski--Faverjon cannot simply be compared
with $H(Q_N)$: in their theorem the height belongs to the relation
polynomial in fixed $E$-values, whereas here the algebraic evaluation point
and its degree move with $N$.

\subsection{The exterior-rank barrier}

Let $C_m$ be the monic Chebyshev polynomials from
\eqref{eq:joukowski-chebyshev} and set
\begin{equation}\label{eq:HN-truncation}
 H_N(w)=\sum_{m=1}^N\frac{(-1)^{m+1}}{m\,m!}C_m(w).
\end{equation}

\begin{theorem}[The level first appears in the top norm]
\label{thm:top-norm-barrier}
Fix a level $\ell\in\C$ and let
$w_{1,N},\ldots,w_{N,N}$ be the roots, with multiplicity, of
$H_N(w)-\ell$.  Then the elementary symmetric functions
\[
 e_1(w_{1,N},\ldots,w_{N,N}),\ldots,
 e_{N-1}(w_{1,N},\ldots,w_{N,N})
\]
are independent of $\ell$, while
\begin{equation}\label{eq:top-norm-exact}
 \boxed{
 \prod_{j=1}^N w_{j,N}
 =N\,N!\bigl(\ell-H_N(0)\bigr).}
\end{equation}
Thus the first elementary symmetric invariant which sees the level is the full
rank-$N$ norm, and it carries the factorial factor $N!$.
\end{theorem}

\begin{proof}
Since $C_N$ is monic, the leading coefficient of $H_N-\ell$ is
$(-1)^{N+1}/(N\,N!)$.  The level $\ell$ changes only the constant
coefficient.  Vieta's formulas therefore make $e_1,\ldots,e_{N-1}$
level-independent, and give
\[
 e_N=(-1)^N
 \frac{H_N(0)-\ell}{(-1)^{N+1}/(N\,N!)}
 =N\,N!\bigl(\ell-H_N(0)\bigr).
\]
\end{proof}

At the Euler level $\ell=2\gamma$, since
$H_N(0)\to H(0)=2\Cin(1)$, one has
\[
 \prod_{j=1}^Nw_{j,N}=2\Ci(1)\,N\,N!\,(1+o(1)),
 \qquad
 \log\left|\prod_{j=1}^Nw_{j,N}\right|
 =N\log N-N+O(\log N).
\]
Thus, by Newton's identities, the power sums of order below $N$ and all
proper exterior traces of the finite zero set forget $\gamma$; the first
elementary symmetric invariant which
remembers it pays the full determinant cost.  A possible route not excluded
by the theorem
would have to remove the universal $N\log N$ contribution canonically, for
example through a zeta-regularized determinant or an arithmetic spectral
projector.  No such arithmetic regularization is constructed here.


\section{Dyadic irreducibility and exact capacity limitations}
\label{sec:dyadic-capacity}

\subsection{Dyadic irreducibility and minimal Weil-height saturation}
\label{subsec:dyadic-irreducibility}

The defining-height result above leaves open, for a general truncation
index, the possibility that the root selected near a zero of \(\Psi\)
belongs to a factor of much smaller degree or height.  On the dyadic
subsequence this possibility can be excluded completely.  Recall
\(H_N\), \(D_N^*\), and \(Q_N\) from
\eqref{eq:HN-truncation} and \eqref{eq:QN-definition}, and put
\begin{equation}\label{eq:PN-AN-dyadic}
 P_N(w)=D_N^*H_N(w),
 \qquad
 A_N=\frac{D_N^*}{N\,N!}.
\end{equation}
Thus, with \(b_{m,N}=D_N^*/(m\,m!)\),
\begin{align}
 P_N(w)
 &=\sum_{m=1}^N(-1)^{m+1}b_{m,N}C_m(w),
 \label{eq:PN-chebyshev}\\
 Q_N(z)
 &=z^NP_N(z+z^{-1})
   =\sum_{m=1}^N(-1)^{m+1}b_{m,N}
       \bigl(z^{N+m}+z^{N-m}\bigr).
 \label{eq:QN-dyadic-explicit}
\end{align}

\begin{theorem}[Dyadic Eisenstein and reciprocal irreducibility]
\label{thm:dyadic-reciprocal-irreducibility}
Let \(N=2^r\geq2\).  Then \(A_N\) is an odd integer and
\begin{equation}\label{eq:PN-mod-four}
 \boxed{P_N(w)\equiv-A_N(w^N+2)\pmod 4.}
\end{equation}
Consequently \(P_N\) is a primitive \(2\)-Eisenstein polynomial and is
irreducible over \(\Q\) of degree \(N\).  Moreover, \(Q_N\) is primitive
and irreducible over \(\Q\) of degree \(2N\).
\end{theorem}

\begin{proof}
Legendre's formula gives
\[
 v_2(N!)=N-s_2(N)=N-1,
\]
where \(s_2(N)\) is the sum of the binary digits of \(N\).  Hence
\begin{equation}\label{eq:dyadic-top-valuation}
 v_2(N\,N!)=N+r-1.
\end{equation}
If \(1\leq m<N\), then \(v_2(m)\leq r-1\) and
\(v_2(m!)\leq m-1\), so
\[
 v_2(m\,m!)\leq m+r-2\leq N+r-3.
\]
It follows that the maximum defining \(v_2(D_N^*)\) is attained at
\(m=N\), that \(A_N\) is odd, and that
\begin{equation}\label{eq:lower-dyadic-coefficients-four}
 4\mid b_{m,N}\qquad(1\leq m<N).
\end{equation}

The identity \(C_{2m}=C_m^2-2\), beginning with
\(C_2(w)=w^2-2\), gives by induction
\begin{equation}\label{eq:dyadic-chebyshev-mod-four}
 C_{2^r}(w)\equiv w^{2^r}+2\pmod4.
\end{equation}
Since \(N\) is even, the coefficient of \(C_N\) in
\eqref{eq:PN-chebyshev} is \(-A_N\).  Equations
\eqref{eq:lower-dyadic-coefficients-four} and
\eqref{eq:dyadic-chebyshev-mod-four} prove
\eqref{eq:PN-mod-four}.  Thus the leading coefficient of \(P_N\) is odd,
all its remaining coefficients are even, and its constant coefficient is
congruent to \(-2A_N\) modulo \(4\).  The polynomial \(Q_N\) is primitive
by Lemma~\ref{lem:truncation-denominator-scale}; hence \(P_N\) is
primitive as well, because any common divisor of the coefficients of
\(P_N\) would divide every coefficient of \(z^NP_N(z+z^{-1})=Q_N\).
Eisenstein's criterion now proves the irreducibility of \(P_N\).

It remains to check that the reciprocal substitution does not split the
polynomial.  Since \(C_m(2)=2\) and \(C_m(-2)=2(-1)^m\),
\begin{align}
 P_N(2)
 &=2D_N^*\sum_{m=1}^N\frac{(-1)^{m+1}}{m\,m!}>0,
 \label{eq:PN-at-two-positive}\\
 P_N(-2)
 &=-2D_N^*\sum_{m=1}^N\frac{1}{m\,m!}<0.
 \label{eq:PN-at-minus-two-negative}
\end{align}
The first sign follows by pairing consecutive terms of the decreasing
alternating sum, since \(N\) is even.

Let \(\theta\) be a root of \(P_N\), let
\(K=\Q(\theta)\), and write \(a_N=\operatorname{lc}(P_N)=-A_N\).
Because \(P_N\) is irreducible of degree \(N\),
\begin{equation}\label{eq:dyadic-discriminant-norm}
 \operatorname N_{K/\Q}(\theta^2-4)
 =\frac{P_N(2)P_N(-2)}{a_N^2}<0.
\end{equation}
A square in \(K\) has norm equal to a square in \(\Q\), so
\(\theta^2-4\) is not a square in \(K\).  Therefore
\(X^2-\theta X+1\) is irreducible over \(K\).  If \(\alpha\) is one of
its roots, then \(\theta=\alpha+\alpha^{-1}\) and consequently
\[
 [\Q(\alpha):\Q]
 =2[K:\Q]=2N.
\]
The number \(\alpha\) is a zero of \(Q_N\), whose degree is \(2N\).
Thus \(Q_N\) is irreducible; its primitivity was already established in
Lemma~\ref{lem:truncation-denominator-scale}.
\end{proof}

\begin{theorem}[Dyadic minimal-height saturation]
\label{thm:dyadic-minimal-height-saturation}
Let \(\rho\) be a fixed simple zero of \(\Psi\) with \(|\rho|>1\), and,
for all sufficiently large \(N\), let \(\alpha_N\) be the zero of
\(\Psi_N\) converging to \(\rho\) in
Theorem~\ref{thm:truncation-height-saturation}.  Along
\(N=2^r\to\infty\),
\begin{equation}\label{eq:dyadic-degree-height}
 [\Q(\alpha_N):\Q]=2N,
 \qquad
 h(\alpha_N)=\frac12\log N+O(1).
\end{equation}
More precisely,
\begin{align}
 \frac{-\log|\alpha_N-\rho|}
 {[\Q(\alpha_N):\Q]h(\alpha_N)}&\longrightarrow1,
 \label{eq:dyadic-root-height-ratio}\\
 \frac{-\log|\Psi(\alpha_N)|}
 {[\Q(\alpha_N):\Q]h(\alpha_N)}&\longrightarrow1.
 \label{eq:dyadic-value-height-ratio}
\end{align}
Here \(h\) is the absolute logarithmic Weil height.  In particular,
\(Q_N\), up to sign, is the primitive integer minimal polynomial of
\(\alpha_N\); the full clearing polynomial is not concealing a factor of
smaller degree or height on the dyadic subsequence.
\end{theorem}

\begin{proof}
By Theorem~\ref{thm:dyadic-reciprocal-irreducibility}, \(Q_N\) is the
primitive irreducible integer polynomial of \(\alpha_N\), up to sign.
The coefficient--Mahler inequalities for a polynomial of degree \(2N\),
together with Lemma~\ref{lem:truncation-denominator-scale}, give
\begin{equation}\label{eq:dyadic-mahler-bounds}
 4^{-N}D_N^*
 \leq M(Q_N)
 \leq \lVert Q_N\rVert_2
 \leq \sqrt{2N+1}\,D_N^*.
\end{equation}
Indeed, the lower bound follows from
\(H(Q_N)\leq2^{2N}M(Q_N)\), while the upper bound is Landau's inequality
and \(H(Q_N)=D_N^*\).  Since
\(\log D_N^*=N\log N+o(N)\), it follows that
\begin{equation}\label{eq:dyadic-mahler-asymptotic}
 \log M(Q_N)=N\log N+O(N).
\end{equation}
The degree assertion in \eqref{eq:dyadic-degree-height} follows from
Theorem~\ref{thm:dyadic-reciprocal-irreducibility}, and the standard height
identity gives
\[
 2N\,h(\alpha_N)=\log M(Q_N).
\]
This proves the height estimate in \eqref{eq:dyadic-degree-height}.

The root asymptotic \eqref{eq:truncation-root-asymptotic} gives
\begin{equation}\label{eq:dyadic-root-log-scale}
 -\log|\alpha_N-\rho|=N\log N+O_\rho(N).
\end{equation}
Moreover \(\Psi_N(\alpha_N)=0\), so the uniform tail estimate used in
the proof of Theorem~\ref{thm:truncation-height-saturation} yields
\[
 \Psi(\alpha_N)
 =\frac{(-1)^{N+2}\alpha_N^{N+1}}
 {(N+1)(N+1)!}
 \left(1+O_\rho(N^{-1})\right).
\]
Since \(\alpha_N\to\rho\),
\begin{equation}\label{eq:dyadic-value-log-scale}
 -\log|\Psi(\alpha_N)|=N\log N+O_\rho(N).
\end{equation}
Dividing \eqref{eq:dyadic-root-log-scale} and
\eqref{eq:dyadic-value-log-scale} by
\(2N h(\alpha_N)=N\log N+O(N)\) proves
\eqref{eq:dyadic-root-height-ratio} and
\eqref{eq:dyadic-value-height-ratio}.
\end{proof}

\begin{corollary}[The precise moving-point threshold]
\label{cor:dyadic-moving-point-threshold}
For every real \(c<1\), infinitely many algebraic numbers in the dyadic
family above satisfy
\begin{equation}\label{eq:dyadic-lower-bound-failure}
 \log|\Psi(\alpha_N)|
 <-c\,[\Q(\alpha_N):\Q]h(\alpha_N).
\end{equation}
Equivalently, no estimate of the form
\[
 -\log|\Psi(\alpha)|
 \leq(c+o(1))[\Q(\alpha):\Q]h(\alpha),
 \qquad c<1,
\]
can hold uniformly on any moving family containing these dyadic points.
For every \(c>1\), the reverse strict inequality to
\eqref{eq:dyadic-lower-bound-failure} holds eventually along this family,
so this sequence alone cannot contradict a bound with coefficient greater
than one.
\end{corollary}

\begin{proof}
This is exactly the limit in
\eqref{eq:dyadic-value-height-ratio}.
\end{proof}

\begin{remark}[Scope of the obstruction]
The corollary identifies the normalized threshold for the term
\([\Q(\alpha):\Q]h(\alpha)\) alone.  It does not rule out estimates with
additional terms of comparable size, such as
\([\Q(\alpha):\Q]\log[\Q(\alpha):\Q]\), nor does it exclude nonlinear
auxiliary constructions or non-dyadic phenomena.  What it closes is the
specific escape through a smaller irreducible factor or smaller minimal
Weil height for the selected truncation root.
\end{remark}


\subsection{Two exact limitations of capacity normalization}
\label{subsec:capacity-pullback-limitations}

The critical constants appearing in the reciprocal and Hankel constructions
are stable under two natural operations.  The statements below are deliberately
limited to polynomial pullbacks of the standard adelic set and to changes of
basis of the full Hankel determinant.  They do not constitute a no-go theorem
for weighted, place-dependent, or nonlinear auxiliary constructions.  This is
the boundary at which the adelic-capacity mechanism used in, for example,
Dimitrov's Schinzel--Zassenhaus argument ceases to improve the present
normalization \cite{Dimitrov2019}; we use the standard local capacity
transformation law in the form reviewed by Bost and Chambert-Loir
\cite{BostChambertLoir2009}.

We normalize the absolute values of \(\Q\) by
\(|p|_p=p^{-1}\) and the usual absolute value at infinity.  At the
archimedean place put
\[
 \mathcal E_\infty=[-2,2],
\]
and, at every finite place, let
\[
 \mathcal E_p=
 \bigl\{z\in\mathbb A_{\C_p}^{1,\mathrm{an}}:|z|_p\leq1\bigr\}
\]
be the full closed Berkovich unit disk.  Capacities below are logarithmic
capacities relative to infinity.

\begin{proposition}[Polynomial pullbacks preserve critical adelic capacity]
\label{prop:adelic-polynomial-capacity}
One has
\[
 \capc_\infty(\mathcal E_\infty)=1,
 \qquad
 \capc_p(\mathcal E_p)=1
 \quad\text{for every prime }p.
\]
Consequently the global capacity of the adelic set
\(\mathcal E=(\mathcal E_v)_{v\leq\infty}\) is
\[
 \prod_{v\leq\infty}\capc_v(\mathcal E_v)=1.
\]
More generally, let
\[
 R(X)=a_dX^d+\cdots+a_0\in\Q[X],
 \qquad d\geq1,\quad a_d\ne0.
\]
Then, at every place \(v\),
\begin{equation}\label{eq:local-capacity-pullback}
 \capc_v\bigl(R^{-1}(\mathcal E_v)\bigr)
 =
 \left(
   \frac{\capc_v(\mathcal E_v)}{|a_d|_v}
 \right)^{1/d},
\end{equation}
and hence
\begin{equation}\label{eq:global-capacity-pullback}
 \boxed{
 \prod_{v\leq\infty}
 \capc_v\bigl(R^{-1}(\mathcal E_v)\bigr)=1.}
\end{equation}
Thus rational affine changes, rational polynomial pullbacks, and in
particular rational Chebyshev pullbacks cannot turn this specific critical
adelic capacity into a strict inequality below one.
\end{proposition}

\begin{proof}
The classical identity
\(\capc([-2,2])=1\) gives the archimedean assertion, and the closed
Berkovich disk of radius one has nonarchimedean capacity one.  The local
polynomial transformation law for logarithmic capacity gives
\eqref{eq:local-capacity-pullback}.  Multiplying it over all places yields
\[
 \prod_v\capc_v\bigl(R^{-1}(\mathcal E_v)\bigr)
 =
 \left(
   \frac{\prod_v\capc_v(\mathcal E_v)}
        {\prod_v|a_d|_v}
 \right)^{1/d}.
\]
The first product in the numerator is one, while the product formula for
\(a_d\in\Q^\times\) gives \(\prod_v|a_d|_v=1\).  This proves
\eqref{eq:global-capacity-pullback}.
\end{proof}

\begin{remark}[Why the finite-place set is not \(\mathbb Z_p\)]
The full Berkovich unit disk in Proposition~\ref{prop:adelic-polynomial-capacity}
must not be replaced silently by the classical compact subset
\(\mathbb Z_p\subset\Q_p\).  With the standard logarithmic kernel,
\begin{equation}\label{eq:capacity-Zp}
 \capc_p(\mathbb Z_p)=p^{-1/(p-1)},
\end{equation}
not one.  Indeed, if \(\mu_p\) is normalized Haar measure on
\(\mathbb Z_p\), then
\[
 \int_{\mathbb Z_p^2}v_p(x-y)\,d\mu_p(x)d\mu_p(y)
 =\sum_{k\geq1}\Pr\bigl(v_p(x-y)\geq k\bigr)
 =\sum_{k\geq1}p^{-k}=\frac1{p-1}.
\]
The corresponding equilibrium energy is
\(\log(p)/(p-1)\), which gives \eqref{eq:capacity-Zp}.
In particular, the literal family with finite components \(\mathbb Z_p\)
does not have global capacity one.
\end{remark}

We next make the critical cancellation in
Theorem~\ref{thm:hankel} integral.  For \(n\geq1\), put
\begin{equation}\label{eq:hankel-integral-clearing}
 \mathcal B_n
 =n!\prod_{j=1}^n
   \binom{2j}{j}\binom{2j-1}{j},
 \qquad
 \ell_n=\operatorname{lcm}(1,\ldots,n),
 \qquad
 \mathcal I_n=\ell_n\mathcal B_n.
\end{equation}

\begin{proposition}[Integral Hankel saturation and full-determinant rigidity]
\label{prop:hankel-integral-saturation}
For every \(n\geq1\) and every \(x\in\C\),
\begin{equation}\label{eq:hankel-cleared-linear-form}
 \mathcal I_nD_n(x)
 =(-1)^n\bigl(2\ell_nH_n-\ell_nx\bigr).
\end{equation}
In particular, the two coefficients on the right side are integers.  One
has
\begin{equation}\label{eq:hankel-clearing-asymptotics}
 \mathcal B_n^{1/n^2}\longrightarrow4,
 \qquad
 \mathcal I_n^{1/n^2}\longrightarrow4.
\end{equation}
For every fixed \(x\in\C\),
\begin{equation}\label{eq:hankel-critical-product}
 |D_n(x)|^{1/n^2}\longrightarrow\frac14,
 \qquad
 \boxed{
 |\mathcal I_nD_n(x)|^{1/n^2}\longrightarrow1.}
\end{equation}

Let
\[
 M_n(x)=(H_{i+j}-x)_{0\leq i,j\leq n}.
\]
If \(U_n,V_n\in\operatorname{GL}_{n+1}(\Z)\), then
\begin{equation}\label{eq:unimodular-hankel-invariance}
 \det\bigl(U_nM_n(x)V_n\bigr)=\pm D_n(x).
\end{equation}
Hence unimodular row and column lattice reduction cannot alter the
capacity constant of the full determinant.
\end{proposition}

\begin{proof}
Theorem~\ref{thm:hankel} gives
\[
 D_n(x)=(-1)^n\frac{2H_n-x}{\mathcal B_n}.
\]
Multiplication by \(\mathcal I_n=\ell_n\mathcal B_n\) proves
\eqref{eq:hankel-cleared-linear-form}; the integrality follows from
\(\ell_nH_n\in\Z\).  The Stirling estimate already used in the proof of
Theorem~\ref{thm:hankel} gives
\[
 \log\mathcal B_n=(\log4)n^2+O(n\log n).
\]
The elementary Chebyshev bound \(\log\ell_n=O(n)\) therefore proves
\eqref{eq:hankel-clearing-asymptotics}.  Since
\(H_n=\log n+O(1)\), for fixed \(x\) one has
\[
 \log|2H_n-x|=O(\log\log n)
\]
for all sufficiently large \(n\).  Formula
\eqref{eq:hankel-critical-product} now follows either directly or together
with Theorem~\ref{thm:hankel}.  Finally,
\[
 \det(U_nM_n(x)V_n)=\det(U_n)\det(V_n)D_n(x),
\]
and both integer determinants are \(\pm1\), proving
\eqref{eq:unimodular-hankel-invariance}.
\end{proof}

\begin{remark}[Exact scope of the limitation]
Proposition~\ref{prop:hankel-integral-saturation} concerns the covolume, or
equivalently the full determinant, of this particular Hankel lattice.
It does not exclude a short individual row or linear form produced by
lattice reduction, a useful proper minor, a change of the underlying
lattice, an \(x\)-dependent auxiliary construction, or a nonlinear
combination of the moments.  Likewise,
Proposition~\ref{prop:adelic-polynomial-capacity} does not exclude weighted
capacity, place-dependent changes of the local compact sets, proper local
subsets, rational maps with poles, or nonpolynomial pullbacks.  The proved
obstruction is precisely that rational polynomial pullback and unimodular
full-determinant normalization preserve the existing critical equality;
no broader capacity or lattice no-go is asserted.
\end{remark}


% Standalone insertion fragment: exact operator-theoretic consequences and scope.

\section{The operator carrier: normal form, determinants, and scope}

We first isolate the universal operator-theoretic content of the construction.
Throughout this section, inner products are linear in the second variable.

\begin{proposition}[Rank-one similarity collapse]\label{prop:rank-one-collapse}
Let $E$ be a complex Hilbert space, let $0\ne\psi,h\in E$, and put
\[
 C=|\psi\rangle\langle h|,
 \qquad \alpha=\langle h,\psi\rangle\ne0.
\]
Let $S$ be the unilateral shift on $\ell^2(\mathbb N_0)$ and define
\[
 \mathsf A=S^*\otimes C,
 \qquad
 \mathsf B=S\otimes C.
\]
There is an $R\in\operatorname{GL}(E)$ and an orthogonal rank-one
projection $P_1$ such that
\[
 RCR^{-1}=\alpha P_1.
\]
Consequently the pair $(\mathsf A,\mathsf B)$ is simultaneously similar to
\[
 ((\alpha S^*)\oplus0,(\alpha S)\oplus0).
\]
Moreover, if $P_0=|e_0\rangle\langle e_0|$, then
\[
 [\mathsf A,\mathsf B]=P_0\otimes C^2,
 \qquad
 \operatorname{Tr}[\mathsf A,\mathsf B]=\alpha^2,
\]
and hence
\begin{equation}\label{eq:universal-hhp}
 \det_F\!\left(
 e^{\mathsf A}e^{\mathsf B}e^{-\mathsf A}e^{-\mathsf B}
 \right)=e^{\alpha^2}.
\end{equation}
\end{proposition}

\begin{proof}
The identity $C^2=\alpha C$ shows that $P=C/\alpha$ is a bounded
rank-one idempotent.  The decomposition
$E=\operatorname{ran}P\dotplus\ker P$ is topological, so a bounded
invertible change of coordinates sends $P$ to an orthogonal rank-one
projection.  Tensoring this change of coordinates with the identity on
$\ell^2$ gives the asserted simultaneous normal form.  Finally,
$S^*S=I$, $SS^*=I-P_0$, and
$\operatorname{tr}(C^2)=\alpha^2$.  The last assertion is the
Helton--Howe--Pincus formula; see, for example, \cite{Migler2017}.
\end{proof}

For the Euler specialization, let
\[
 H=L^2(\mathbb R_+,e^{-x}\,dx),\qquad Xf(x)=xf(x),\qquad A=X+I,
\]
take $E=H\oplus\mathbb C$, and put
\[
 \psi=(-A^{-1}1,1),\qquad h_\gamma=(A\log X\,1,0).
\]
Both vectors belong to $E$, and direct integration gives
\[
 \langle h_\gamma,\psi\rangle
 =-\int_0^\infty e^{-x}\log x\,dx=\gamma.
\]
Thus $C=|\psi\rangle\langle h_\gamma|$ has $\alpha=\gamma$, and
\eqref{eq:universal-hhp} becomes $e^{\gamma^2}$.  The proposition also
makes the limitation precise: the simultaneous similarity class sees only
the scalar $\gamma$ and is blind to the zero-mode norm $2-\delta$ and to
the remaining Hilbert-complex geometry.

There is, however, a natural relative Fredholm determinant on the
$\delta$-side.  With the preceding $H$ and $A$, put $u=A^{-1/2}1$ and let
\[
 Q:\mathcal D(A^{1/2})\oplus\mathbb C\longrightarrow H,
 \qquad Q(f,c)=A^{1/2}f+cu.
\]

\begin{proposition}[Zero-mode perturbation determinant]
The operator $Q$ is a closed surjective Fredholm operator of index one,
\[
 QQ^*=A+|u\rangle\langle u|,
 \qquad
 \ker Q=\mathbb C\psi,
 \qquad
 \psi=(-A^{-1}1,1).
\]
For $z\in\mathbb C\setminus[1,\infty)$ define
\[
 \Delta_Q(z)=
 \det_F\!\left((QQ^*-z)(A-z)^{-1}\right).
\]
Here the product denotes its bounded extension from the natural common
domain.
Then
\begin{align}
 \Delta_Q(z)
 &=1+\langle u,(A-z)^{-1}u\rangle \notag\\
 &=1+\int_0^\infty
       \frac{e^{-x}}{(1+x)(1+x-z)}\,dx.\label{eq:relative-determinant}
\end{align}
Writing $\delta=eE_1(1)$ and taking the principal branch of $E_1$, this is
\[
 \Delta_Q(z)=1+
 \frac{e^{1-z}E_1(1-z)-\delta}{z}
 \qquad(z\ne0),
\]
with a removable value at zero given by
\begin{equation}\label{eq:delta-relative-determinant}
 \Delta_Q(0)=2-\delta=\|\psi\|^2.
\end{equation}
\end{proposition}

\begin{proof}
Surjectivity follows by taking $(f,c)=(A^{-1/2}g,0)$ for $g\in H$;
the kernel formula is immediate.  Direct computation gives
$QQ^*=A+|u\rangle\langle u|$.  The rank-one Fredholm determinant lemma
then yields \eqref{eq:relative-determinant}.  Finally,
\[
 \int_0^\infty\frac{e^{-x}}{(1+x)^2}\,dx=1-\delta,
\]
which proves \eqref{eq:delta-relative-determinant} and the index assertion.
\end{proof}

Let $\mathscr H=\ell^2(\mathbb N_0)\otimes E$, let
$\pi:\mathcal B(\mathscr H)\to
\mathcal B(\mathscr H)/\mathcal L^1(\mathscr H)$ be the quotient map, and
put $a=\pi(e^{\mathsf A})$, $b=\pi(e^{\mathsf B})$.  These elements commute.
If $\xi=\{a,b\}\in
K_2^{\mathrm{alg}}(\mathcal B(\mathscr H)/\mathcal L^1(\mathscr H))$, then
Brown's determinant invariant gives~\cite{Migler2017,Kaad2011}
\begin{equation}\label{eq:k2-boundary}
 \det\bigl(\partial_2\xi\bigr)
 =\det_F(e^{\mathsf A}e^{\mathsf B}
          e^{-\mathsf A}e^{-\mathsf B})
 =e^{\alpha^2}.
\end{equation}
In particular, when $e^{\alpha^2}\ne1$, both
$\partial_2\xi\in
K_1(\mathcal B(\mathscr H),\mathcal L^1(\mathscr H))$ and $\xi$ are
nontrivial.  Notice that \eqref{eq:k2-boundary} is an algebraic
$K_2$ statement detected on a relative $K_1$ boundary; it is not a
``relative $K_2$'' class.

\paragraph{Scope.}
The preceding statements do not produce an arithmetic relation between
$e^{\gamma^2}$ and $\delta$.  Indeed, replacing $h$ by any covector with
$\langle h,\psi\rangle=\alpha$ realizes the arbitrary scalar
$e^{\alpha^2}$, while Proposition~\ref{prop:rank-one-collapse} erases
$\delta$ completely.  Nor does \eqref{eq:k2-boundary} yield a new
topological $C^*$-algebra $K$-class: after passage to the Calkin algebra,
the two displayed invertibles are exponentials and hence have trivial
individual topological $K_1$ classes.

Under the Hardy-space identification, the active normal form is the
one-mode Toeplitz pair with symbols $e^{\alpha z^{-1}}$ and $e^{\alpha z}$.
For $\alpha\ne0$, the first symbol has an essential singularity at the
origin, so meromorphic tame-symbol formulas do not apply.  Finally, no
Wodzicki-residue interpretation follows from this model: no classical
pseudodifferential calculus or spectral triple has been specified, the
relevant commutator is finite rank and hence is annihilated by singular
traces, and a multiplicative residue determinant sends a multiplicative
commutator to $1$ whenever all factors lie in its multiplicative domain.
Thus \eqref{eq:universal-hhp} is an ordinary Fredholm
determinant anomaly, not a noncommutative-residue anomaly.


\section{Frobenius framing and a depth-one rapid-decay target}
\label{sec:frobenius-framing}

The additive normalization freedom visible in
\eqref{eq:tautological-recovery} and in the Kummer tower is often expressed
by the unipotent matrices
\[
 S_t=\begin{pmatrix}1&t\\0&1\end{pmatrix}.
\]
Inside an unframed Tannakian category, adding further tensors already fixed
by the motivic Galois group cannot remove this direction.  A Frobenius
operator with distinct graded eigenvalues does remove the corresponding
centralizer, but this fact has to be interpreted carefully: after forgetting
the filtration and comparison data, the same separation of slopes also makes
the local extension split.  We begin by recording both sides of this linear
algebra.

\begin{lemma}[Frobenius-framed rigidity]
\label{lem:frobenius-framed-rigidity}
Let $K$ be a field of characteristic zero and let $V$ be a two-dimensional
$K$-realization with a $K$-linearized Frobenius operator (for example, an
$f$th iterate of a semilinear Frobenius for which
$\sigma^f=\operatorname{id}_K$) which, in a basis compatible with its two
graded pieces, has matrix
\begin{equation}\label{eq:frobenius-matrix}
 \Phi_p=\begin{pmatrix}1&c_p\\0&q_p\end{pmatrix},
 \qquad q_p\ne1.
\end{equation}
Then the unipotent subgroup $U=\{S_t:t\in K\}$ has trivial intersection
with the centralizer of $\Phi_p$:
\begin{equation}\label{eq:frobenius-centralizer}
 U\cap Z_{\mathrm{GL}(V)}(\Phi_p)=\{1\}.
\end{equation}
In particular this applies when $q_p=p^f$ or $p^{-f}$, including $f=1$.
\end{lemma}

\begin{proof}
Direct multiplication gives
\[
 S_t\Phi_p-\Phi_pS_t
 =\begin{pmatrix}0&(q_p-1)t\\0&0\end{pmatrix}.
\]
Since $q_p\ne1$, commutation is equivalent to $t=0$.
\end{proof}

For a raw $\sigma$-semilinear Frobenius, the corresponding centralizer
condition is $q_pt=\sigma(t)$ rather than $(q_p-1)t=0$.  The associated
inhomogeneous equation has the opposite-looking consequence.

\begin{theorem}[Automatic splitting of a bare Frobenius extension]
\label{thm:semilinear-frobenius-splitting}
Let $K$ be a complete non-archimedean field, let
$\sigma\colon K\to K$ be an isometric automorphism, and let
$q\in K^\sigma$ satisfy $|q|\ne1$.  For every $c\in K$, consider the
$\sigma$-semilinear operator whose matrix is
\[
 A=\begin{pmatrix}1&c\\0&q\end{pmatrix}.
\]
There is a unique $t\in K$ such that
\begin{equation}\label{eq:semilinear-splitting-equation}
 qt-\sigma(t)=c.
\end{equation}
For this value of $t$ one has
\begin{equation}\label{eq:semilinear-diagonalization}
 S_t^{-1}A\,\sigma(S_t)
 =\begin{pmatrix}1&0\\0&q\end{pmatrix}.
\end{equation}
Thus the underlying unfiltered semilinear Frobenius extension is split, and
its upper-unipotent splitting is unique.
\end{theorem}

\begin{proof}
If $|q|<1$, the convergent series
\[
 t=-\sum_{n\ge0}q^n\sigma^{-n-1}(c)
\]
satisfies \eqref{eq:semilinear-splitting-equation}.  If $|q|>1$, the
corresponding solution is
\[
 t=\sum_{n\ge0}q^{-n-1}\sigma^n(c).
\]
Convergence follows because $\sigma$ is isometric.  If $t_1$ and $t_2$ are
two solutions, their difference $u$ satisfies $qu=\sigma(u)$; taking
absolute values gives $|q||u|=|u|$, and hence $u=0$.  Finally,
\[
 S_t^{-1}A\,\sigma(S_t)
 =\begin{pmatrix}1&c+\sigma(t)-qt\\0&q\end{pmatrix},
\]
which proves \eqref{eq:semilinear-diagonalization}.
\end{proof}

Theorem~\ref{thm:semilinear-frobenius-splitting} does not contradict
Lemma~\ref{lem:frobenius-framed-rigidity}.  The homogeneous centralizer
equation has only the zero solution, while the inhomogeneous splitting
equation has exactly one solution.  Consequently a marked Frobenius basis
may remove a gauge ambiguity, but the bare Frobenius isocrystal does not by
itself retain the global extension class.

\begin{proposition}[A Kummer local-to-global counterexample]
\label{prop:kummer-bare-frobenius-counterexample}
Let
\[
 M_2=[\mathbb Z\xrightarrow{1\mapsto2}\mathbb G_m]
\]
be the Kummer $1$-motive over $\mathbb Q$.  Its canonical sequence
\begin{equation}\label{eq:kummer-global-extension}
 0\longrightarrow[0\to\mathbb G_m]
 \longrightarrow M_2
 \longrightarrow[\mathbb Z\to0]
 \longrightarrow0
\end{equation}
is non-split in the isogeny category.  Nevertheless, $M_2$ has good
reduction at every prime $p\ne2$, and at each such prime its underlying
rational crystalline $F$-isocrystal is split.  Hence the implication
\[
 \left.
 \begin{array}{c}
 \text{the unfiltered rational $F$-isocrystal is split}\\
 \text{at every good prime}
 \end{array}
 \right\}
 \quad\Longrightarrow\quad
 \text{the global motive is split}
\]
is false.
\end{proposition}

\begin{proof}
A splitting of \eqref{eq:kummer-global-extension} up to isogeny would make
the point $2\in\mathbb G_m(\mathbb Q)$ torsion, equivalently would force
$2^m=1$ for some $m\ge1$.  Thus the global extension is non-split.  For
$p\ne2$, the point $2$ extends to a unit over $\mathbb Z_p$, so the motive
has good reduction.  Its reduction $\overline{2}\in\mathbb F_p^\times$ has
finite order.  The reduced Kummer $1$-motive therefore splits up to
isogeny, and applying the rational crystalline realization gives a split
sequence with graded slopes $0$ and $1$, up to the convention for signs.
This is the root-of-unity case of the Kummer calculation in
\cite[Example~6.3]{MohajerSebbar2025}; the surrounding slope formalism may be
compared with \cite{KedlayaIsocrystals}.
\end{proof}

The qualifier \emph{unfiltered} in the proposition is essential.  Neither
the $p$-adic Galois realization nor the filtered Frobenius module is asserted
to split.  The relative position of the slope splitting and the Hodge--de
Rham filtration can retain the Kummer extension; in the preceding example
its coordinate is represented, subject to conventions, by $\log_p2$.
Accordingly, a viable arithmetic enhancement of the Euler extension must
use filtered-Frobenius, $p$-adic integration, or syntomic comparison data,
not the bare $F$-isocrystal alone.

The rapid-decay pairing itself is well established: Bloch and Esnault
construct it for irregular connections on curves, while Hien develops the
surface and arbitrary-dimensional theories
\cite{BlochEsnault2004,Hien2007,Hien2009}.  The classical arithmetic
comparison is also instructive.  Huber and W\"ustholz determine the
algebraic linear relations among periods of classical $1$-motives
\cite{HuberWustholz}.  For $1$-motives over number fields with good
reduction, Mohajer and Sebbar construct canonical rational structures and a
$p$-adic analytic-subgroup framework at depths one and two
\cite{MohajerSebbar2025}.  On the irregular side, Snodgrass constructs
rational rapid-decay realizations and $E$-periods in the $E$-operator
setting under his stated hypotheses \cite{Snodgrass2026}.  None of these
results supplies the numerical-to-motivic implication needed below.

Let $M_\gamma$ denote the Euler--Mascheroni exponential motive in the
normalization of Fres\'an and Jossen.  They prove that it is a non-split
extension
\begin{equation}\label{eq:euler-motive-extension}
 0\longrightarrow\mathbb Q(0)
 \longrightarrow M_\gamma
 \longrightarrow\mathbb Q(-1)
 \longrightarrow0,
\end{equation}
that $\operatorname{Hom}(M_\gamma,\mathbb Q(0))=0$, and that its complex
extension coordinate is $\gamma$ \cite[Sec.~12.8]{FresanJossen}.

The relevant motivic separation from classical periods is already known;
it is not an additional mixed-extension conjecture.

\begin{theorem}[Fres\'an--Jossen classical separation]
\label{thm:fresan-jossen-classical-separation}
Let $N$ be any classical motive over $\mathbb Q$, put
$N^+=N\oplus M_\gamma$, and denote their motivic Galois groups by
$G_N$ and $G_{N^+}$.  Then the natural faithfully flat morphism
\[
 G_{N^+}\longrightarrow G_N
\]
has a subgroup isomorphic to $\mathbb G_a$ in its kernel.  In particular,
\begin{equation}\label{eq:fresan-jossen-dimension-jump}
 \dim G_{N^+}>\dim G_N.
\end{equation}
The statement applies, in particular, to every finite direct sum of Tate
and Kummer motives.
\end{theorem}

\begin{proof}
Fres\'an and Jossen show that the perverse realization of a classical
motive has trivial fundamental group, whereas adjoining $M_\gamma$
introduces the full unipotent group $\mathbb G_a$.  Their comparison diagram
embeds this copy of $\mathbb G_a$ into $G_{N^+}$ and maps it trivially to
$G_N$.  This is precisely the proof of
\cite[Cor.~12.8.12]{FresanJossen}; see also their
Proposition~12.8.7 and Corollary~12.8.8.
\end{proof}

Theorem~\ref{thm:fresan-jossen-classical-separation} settles the motivic
part of the Euler--Kummer separation problem, object by object for every
finite classical family.  What remains open is the comparison step that
turns a numerical relation among complex period coordinates into a motivic
relation.  The minimal form needed for $\gamma$ is the following degree-one
fullness statement.

\begin{conjecture}[Euler rapid-decay comparison fullness]
\label{conj:rapid-decay-analytic-subgroup}
Assume that $M_\gamma$ is equipped with its rational rapid-decay realization
and with coherent filtered-Frobenius or syntomic comparison data at the good
places.  If
\begin{equation}\label{eq:euler-splitting-relation}
 \lambda\gamma+\beta=0,
 \qquad \lambda,\beta\in\Qbar,\quad \lambda\ne0,
\end{equation}
then this numerical relation is induced by a morphism in the
exponential-motivic category; equivalently in the present rank-two setting,
it forces a splitting of \eqref{eq:euler-motive-extension}.
\end{conjecture}

This is deliberately a rank-two, degree-one specialization of the
exponential-period philosophy.  It is not a consequence of the existence
of rapid-decay homology, a rational structure, or bare Frobenius
realizations.  The filtered or syntomic data are possible tools for proving
the comparison assertion, not substitutes for it.

\begin{proposition}[Conditional transcendence]
\label{prop:rapid-decay-implies-gamma}
Conjecture~\ref{conj:rapid-decay-analytic-subgroup} implies that $\gamma$ is
transcendental.
\end{proposition}

\begin{proof}
If $\gamma$ were algebraic, then
\eqref{eq:euler-splitting-relation} would hold with $\lambda=1$ and
$\beta=-\gamma$.  The conjecture would force
\eqref{eq:euler-motive-extension} to split, contradicting the non-splitting
proved by Fres\'an and Jossen
\cite[Prop.~12.8.7]{FresanJossen}.
\end{proof}

\begin{conjecture}[Euler--Baker linear period theorem]
\label{conj:euler-baker}
Let $\alpha_1,\ldots,\alpha_r\in\Qbar^\times$ be multiplicatively
independent, and fix branches of their logarithms.  If
\begin{equation}\label{eq:euler-baker-relation}
 \lambda\gamma+\beta+\nu(2\pi i)
 +\sum_{j=1}^r\mu_j\Log\alpha_j=0,
 \qquad
 \lambda,\beta,\nu,\mu_j\in\Qbar,
\end{equation}
then
\[
 \lambda=\beta=\nu=\mu_1=\cdots=\mu_r=0.
\]
\end{conjecture}

For clarity, let $N_{\boldsymbol\alpha}$ be a finite classical motive
containing the relevant Tate and Kummer objects, so that its period matrix
contains $1$, $2\pi i$, and the chosen $\Log\alpha_j$.

\begin{proposition}[The exact relative comparison gap]
\label{prop:relative-fullness-implies-euler-baker}
Assume degree-one comparison fullness for
$N_{\boldsymbol\alpha}\oplus M_\gamma$: every $\Qbar$-linear relation among
the displayed comparison-matrix entries is induced by motivic
functoriality.  Then Conjecture~\ref{conj:euler-baker} holds for
$\alpha_1,\ldots,\alpha_r$.
\end{proposition}

\begin{proof}
Under comparison fullness, a relation
\eqref{eq:euler-baker-relation} is motivic.  The copy of $\mathbb G_a$ in
the kernel furnished by
Theorem~\ref{thm:fresan-jossen-classical-separation} fixes all period
coordinates coming from $N_{\boldsymbol\alpha}$ and translates the Euler
extension coordinate.  Invariance of the relation therefore forces
$\lambda=0$.  The remaining relation
\[
 \beta+\nu(2\pi i)+\sum_{j=1}^r\mu_j\Log\alpha_j=0
\]
is classical.  The Baker--Lindemann theorem \cite{Baker1975}, together with the assumed
multiplicative independence, gives
$\beta=\nu=\mu_1=\cdots=\mu_r=0$.
\end{proof}

Thus no additional theorem asserting motivic independence of the Euler
extension from the span of Kummer extensions is missing: that separation is
already contained in
Theorem~\ref{thm:fresan-jossen-classical-separation}.  The genuine missing
input for the Euler--Baker statement is the relative
\emph{numerical-to-motivic comparison theorem} used in
Proposition~\ref{prop:relative-fullness-implies-euler-baker}.

\begin{remark}[An adelic formulation of the missing estimate]
Let $K$ now be a number field.  The local Dwork scales
$p^{-1/(p-1)}$ should not be identified numerically across different
primes.  A viable slope strategy would first have to place coherent
filtered-Frobenius or syntomic norms on one rapid-decay extension and then
form an adelic degree
\[
 \widehat{\deg}(Ks)=-\sum_{w\in M_K}n_w\log\|s\|_w.
\]
It would additionally require a theorem showing that an algebraic complex
period relation produces a global section of impossible adelic degree
unless the extension splits.  The bare slope decomposition cannot provide
this: Proposition~\ref{prop:kummer-bare-frobenius-counterexample} shows
that it can split at every good place while the global motive remains
non-split.  Neither the required coherence of the filtered local norms nor
the global slope inequality is presently proved.
\end{remark}

The conclusion is diagnostic.  Frobenius centralizer rigidity is a correct
uniqueness statement, but the rank-two scalar semilinear model of
Theorem~\ref{thm:semilinear-frobenius-splitting} splits automatically when
\(|q|\ne1\).  Fres\'an and Jossen already prove that the
Euler $\mathbb G_a$-direction survives after adjoining every fixed
classical motive.  The remaining depth-one problem is therefore exactly a
comparison/fullness theorem transferring a complex rapid-decay relation to
the motivic category.  Its minimal rank-two case would prove the
transcendence of $\gamma$; its finite relative Kummer form would prove the
Euler--Baker linear period theorem.


\section{Corrections, closure ledger, and remaining problems}

The proof audit changes the status of several statements in the preceding
draft.  The multipoint theorem now includes the Laurent-denominator argument
needed to pass from a rational differential field to a finite Laurent sum.
The Pad\'e recurrence is proved for every degree from the Laguerre recurrence
and the Pad\'e Casoratian, rather than inferred from finite symbolic checks.
The Laguerre--Hankel identity is proved for every \(n\) by the
Geronimus--Markov formula.  The reciprocal-Gamma calculation uses only
square-free mixed derivatives in distinct variables; it makes no assertion
that all ordinary higher derivatives vanish.

Three former interpretations are explicitly withdrawn.  A formal shift of a
differential equation is not an arithmetic impossibility theorem unless it
preserves the Stieltjes boundary condition.  The reciprocal Euler equation
has infinitely many off-circle complex zeros, so its single unit-circle orbit
is not the whole zero set.  Finally, a membership equivalence between two
formal classes is not a reduction of the transcendence problem unless the
new membership question is independently decidable.

Two further effectivity corrections are recorded in the present version.
First, \(E_1\) is not a Siegel \(E\)-function: its logarithmic singularity
at the origin contains the same connection constant which one is trying to
control.  The exact recovery identity
\eqref{eq:tautological-recovery} is therefore not an effective
\(E\)-value representation of \(\gamma\).  Second, the height in the
Adamczewski--Faverjon measure is the height of a relation polynomial for
fixed functions at a fixed algebraic point, not the height of a moving
algebraic approximation point.  A former comparison of a putative
\(\kappa<1\) requirement with a Dirichlet exponent is consequently
withdrawn.  The correct moving-point difficulty is quantified instead by
Theorem~\ref{thm:truncation-height-saturation}.

Version 6.2 makes four additional cross-field corrections.  The classical
compact set \(\mathbb Z_p\) has capacity \(p^{-1/(p-1)}\), not one; the
capacity-one local object used here is the full closed Berkovich unit disk.
The auxiliary primitive \(\mathcal K\) used at the Euler point is entire and
holonomic but is not an \(E\)-function; its algebraic-point values nevertheless
belong to the ring of finite \(E\)-values by the explicit finite
\(\Ein\)-value formula \eqref{eq:ssc-K-finite-E-values}.
The Hermite--Biehler construction yields a de Branges space and discrete
Clark spectra, but no unproved regularity of a canonical-system Hamiltonian
is asserted.  Finally, in the proved rank-two scalar semilinear model,
\(|q|\ne1\) makes the underlying bare extension split locally.  Thus raw
good-prime splitting cannot
recover a global extension class.  Conversely, Fres\'an--Jossen already prove
the required motivic \(\mathbb G_a\)-separation from every fixed classical
motive; the remaining gap is numerical-to-motivic comparison fullness, not a
new mixed-\(\operatorname{Ext}^1\) independence theorem.

Version 6.3 records a further five-document audit.  The Bell--Gould
prime-power lift and its coefficientwise CRT synthesis survive as exact
divisibility theorems, but neither implies small primitive linear forms.
A proposed compactifier estimate is false already at its first index, and a
quoted finite denominator bound lacks the recurrence, primitive denominator,
and strict numerical inequality needed for certification; neither is used
here.  In a separate compact double-integral family, the formerly stated
\(16^{-n}\) asymptotic misses a factor \(n^{-1}\), and the partial-fraction
formula omits higher-order poles.  The corrected family and its unavoidable
positive logarithmic contamination are stated only in their proved range.
Crossing the surviving multipoint theorem with the classical
polyexponential differential chain yields the resonant all-order theorem,
the all-order Gamma-jet split, and the exact \(2N-1\) rank of each finite
Jacobi spectral tower.  These deductions have complete proofs here, but the
one-point nonresonant polyexponential phenomenon is classical and no
historical novelty is claimed for it.  A proposed multiplicative-dependence
criterion for general algebraic arguments is false and is not used.

\begingroup
\footnotesize
\renewcommand{\arraystretch}{1.12}
\begin{longtable}{>{\raggedright\arraybackslash}p{0.20\linewidth}
                >{\raggedright\arraybackslash}p{0.34\linewidth}
                >{\raggedright\arraybackslash}p{0.34\linewidth}}
\caption{Cross-field closure ledger.  Here \emph{closure} means that the
stated finite classification or obstruction is proved, not that a named
long-standing problem is solved.}
\label{tab:closure-ledger}\\
\toprule
Imported area & Exact closure in this paper & What remains open \\
\midrule
\endfirsthead
\toprule
Imported area & Exact closure in this paper & What remains open \\
\midrule
\endhead
\bottomrule
\endfoot
Differential algebra and toric geometry & Full finite relation ideal for
exponential and \(\Ein\) coordinates; all-order resonant one-level
polyexponential tower; pure-exponential base change; common-shift defect one
and the exact residual dichotomy & The functions and the one-point
nonresonant phenomenon are classical; exact priority of the multipoint
resonant toric formulation still requires specialist review. \\

Graph, matroid, and Grassmannian geometry & Cycle-space ideals; every
\(\Qbar\)-representable algebraic matroid; exact Pl\"ucker ideal for explicit
\(\Ein\)-matrices & These are arithmetic realizations of classical generic
constructions, not new classifications of matroids or Grassmannians. \\

Classical special functions and moment theory & Countable independence of
the integral and quadratic \(E_1\) families; all-algebraic-level four-value
exclusion; all-order positive Hausdorff systems with strict total positivity,
independent Hankel/Favard coordinates, and exact finite spectral-tower rank
\(2N-1\); all-order Gamma-jet/tail split & These constructions do not isolate
\(\gamma\), \(\delta\), individual zeta values, or individual tail
integrals; exact priority of the arithmetic realizations remains to be
audited. \\

\(p\)-adic and umbral congruences & Exact all-prime radius phase, denominator
law, \(p\)-curvature, Borel logarithm, prime-power Bell--Gould annihilators,
and arbitrary-modulus CRT synthesis & Individual irrationality of \(p\)-adic
Euler sums and \(\delta\), as well as analytic smallness of CRT
representatives, remain open. \\

Orthogonal polynomials and isomonodromy & All-degree Kummer determinant,
Riccati/\(P_V\) boundary, positive mixed determinants, finite recovery of
\(\Q(\gamma,\delta)\), Nevai limits, and arcsine spectral distribution & The
classical normalizations produce no small linear forms for \(\delta\). \\

Statistical mechanics & Monomer--dimer representation and a complex vertical
zero strip & The determinant is not real-rooted, and no Lee--Yang theorem is
obtained. \\

Entire-function and Clark geometry & Finite \(\Ein\)-witness rigidity;
Joukowski descent, rightmost-fibre stability, simple transcendental zero
sets, meromorphic inner quotient, algebraically independent scattering
samples, two \(\gamma\)-recovery limits, and exact \(S_2,S_3\) trace
inversion & The recoveries do not provide the uniform moving-point
arithmetic measure needed to decide \(\gamma\). \\

Reciprocal spectral arithmetic & Exact Euler-level Hadamard and
inverse-resonance traces, \(\gamma\)-free phase recovery, top-norm barrier, and
dyadic irreducibility with actual minimal-Weil-height saturation & No
arithmetic regularized determinant or uniform moving-point \(E\)-value
measure is available; non-dyadic and nonlinear auxiliaries remain open. \\

Adelic capacity and Hankel lattices & Rational polynomial pullback preserves
global capacity one; primitive Hankel clearing saturates at one; the scalar
one-pole Borel--Laguerre family has an exact factorial wall; the compact
Sondow variant has strictly positive logarithmic contamination & Weighted or
place-dependent capacity, coupled determinants, general rational maps with
poles, short vectors, and nonlinear auxiliaries are not excluded. \\

Operator theory & Exact Fredholm determinant and relative determinant, plus
similarity collapse & The carrier forgets the coupling needed to relate
\(e^{\gamma^2}\) to \(\delta\); no residue or topological \(K\)-theory
solution follows. \\

Frobenius and exponential motives & The rank-two scalar bare-Frobenius model
with \(|q|\ne1\) splits uniquely; local splitting does not imply global
splitting; the Euler
\(\mathbb G_a\)-direction is motivically separated from every fixed classical
motive & A filtered-Frobenius or syntomic numerical-to-motivic
comparison/fullness theorem remains open. \\
\end{longtable}
\endgroup

The remaining high-value targets are therefore sharply delimited:
\begin{enumerate}[leftmargin=2em]
\item determine whether either \(\gamma\) or \(\delta\) is irrational;
\item produce an actual relation or fixed-degree super-small sequence in a
detector such as \eqref{eq:Tm-delta-detector} or
\eqref{eq:fixed-degree-super-small};
\item prove a genuinely nonvacuous intersection theorem between finite-point
\(E\)-values and arithmetic-Gevrey values;
\item find a weighted, place-dependent, or nonlinear normalization outside
the rational polynomial pullbacks and full-determinant transformations proved
critical in Propositions~\ref{prop:adelic-polynomial-capacity} and
\ref{prop:hankel-integral-saturation};
\item decide the arithmetic nature of the unique real Euler-level point
\(w_*\) in Theorem~\ref{thm:H-euler-level}.  Its algebraicity would already
imply the algebraic independence of \(\gamma\) and \(\delta\).
\item construct an arithmetic zeta-regularized determinant or spectral
projector which removes the universal \(N\log N\) term in
Theorem~\ref{thm:top-norm-barrier} without removing the reciprocal phase;
\item prove a mixed \(E\)-value/Baker measure uniform in the algebraic sample
point for the Clark scattering sequence of
Theorem~\ref{thm:ssc-arithmetic-Clark};
\item prove the Euler rapid-decay splitting principle
\ref{conj:rapid-decay-analytic-subgroup}; Fres\'an--Jossen already provide the
required motivic separation from every fixed classical family, so the missing
input is precisely numerical-to-motivic comparison fullness.
\end{enumerate}

\section*{Acknowledgement of scope}

This document is a revised research draft, not a peer-reviewed publication.
The results labelled as applications or syntheses rely on the cited theorems,
and their exact priority should be checked by a specialist literature review
before submission.  The negative statements are intentionally local to the
specified formal classes and normalizations.

\begin{thebibliography}{99}
\small

\bibitem{AdamczewskiFaverjon2025}
B.~Adamczewski and C.~Faverjon,
\newblock Algebraic independence measures for values of $E$-functions and
$M$-functions,
\newblock arXiv:2502.09999 (2025).
\newblock \url{https://arxiv.org/abs/2502.09999}.

\bibitem{Andre1989}
Y.~Andr\'e,
\newblock \emph{$G$-Functions and Geometry},
\newblock Aspects of Mathematics E13, Vieweg, Braunschweig, 1989.

\bibitem{Andre2000}
Y.~Andr\'e,
\newblock S\'eries Gevrey de type arithm\'etique, I: th\'eor\`emes de puret\'e et de dualit\'e,
\newblock \emph{Ann. of Math.} \textbf{151} (2000), 705--740.
\newblock \url{https://arxiv.org/abs/math/0003238}.

\bibitem{Baker1975}
A.~Baker,
\newblock \emph{Transcendental Number Theory},
\newblock Cambridge University Press, Cambridge, 1975.

\bibitem{Beukers2006}
F.~Beukers,
\newblock A refined version of the Siegel--Shidlovskii theorem,
\newblock \emph{Ann. of Math.} \textbf{163} (2006), 369--379.
\newblock \url{https://doi.org/10.4007/annals.2006.163.369}.

\bibitem{BostChambertLoir2009}
J.-B.~Bost and A.~Chambert-Loir,
\newblock Analytic curves in algebraic varieties over number fields,
\newblock in \emph{Algebra, Arithmetic, and Geometry}, vol.~I,
Progress in Mathematics 269, Birkh\"auser, 2009, 69--124.
\newblock \url{https://arxiv.org/abs/math/0702593}.

\bibitem{BlochEsnault2004}
S.~Bloch and H.~Esnault,
\newblock Homology for irregular connections,
\newblock \emph{J. Th\'eor. Nombres Bordeaux} \textbf{16} (2004), 357--371.
\newblock \url{https://doi.org/10.5802/jtnb.450}.

\bibitem{BrunsConcaVarbaro2013}
W.~Bruns, A.~Conca and M.~Varbaro,
\newblock Relations between the minors of a generic matrix,
\newblock \emph{Adv. Math.} \textbf{244} (2013), 171--206.
\newblock \url{https://arxiv.org/abs/1111.7263}.

\bibitem{ChoeOxleySokalWagner2004}
Y.-B.~Choe, J.~G.~Oxley, A.~D.~Sokal and D.~G.~Wagner,
\newblock Homogeneous multivariate polynomials with the half-plane property,
\newblock \emph{Adv. in Appl. Math.} \textbf{32} (2004), 88--187.
\newblock \url{https://arxiv.org/abs/math/0202034}.

\bibitem{ChristolEtAl1980}
G.~Christol, T.~Kamae, M.~Mend\`es France and G.~Rauzy,
\newblock Suites alg\'ebriques, automates et substitutions,
\newblock \emph{Bull. Soc. Math. France} \textbf{108} (1980), 401--419.
\newblock \url{https://archive.numdam.org/articles/10.24033/bsmf.1926/}.

\bibitem{Clark1972}
D.~N.~Clark,
\newblock One dimensional perturbations of restricted shifts,
\newblock \emph{J. Anal. Math.} \textbf{25} (1972), 169--191.
\newblock \url{https://doi.org/10.1007/BF02790036}.

\bibitem{Cigler2017}
J.~Cigler,
\newblock Hankel determinants of harmonic numbers and related topics,
\newblock arXiv:1703.00763 (2017).
\newblock \url{https://arxiv.org/abs/1703.00763}.

\bibitem{DeanoHuertasRoman2018}
A.~Dea\~no, E.~J.~Huertas and P.~Rom\'an,
\newblock Asymptotics of orthogonal polynomials generated by a Geronimus
perturbation of the Laguerre measure,
\newblock arXiv:1504.05976, revised 2018.
\newblock \url{https://arxiv.org/abs/1504.05976}.

\bibitem{deBranges1968}
L.~de Branges,
\newblock \emph{Hilbert Spaces of Entire Functions},
\newblock Prentice--Hall, Englewood Cliffs, 1968.

\bibitem{Delaygue2025}
\'E.~Delaygue,
\newblock A Lindemann--Weierstrass theorem for $E$-functions,
\newblock arXiv:2210.12046v2 (2025).
\newblock \url{https://arxiv.org/abs/2210.12046}.

\bibitem{Dieudonne1957}
J.~Dieudonn\'e,
\newblock On the Artin--Hasse exponential series,
\newblock \emph{Proc. Amer. Math. Soc.} \textbf{8} (1957), 210--214.
\newblock \url{https://doi.org/10.1090/S0002-9939-1957-0087034-9}.

\bibitem{DLMF}
NIST Digital Library of Mathematical Functions,
\newblock Chapters 13 and 18.
\newblock \url{https://dlmf.nist.gov/}.

\bibitem{Dimitrov2019}
V.~Dimitrov,
\newblock A proof of the Schinzel--Zassenhaus conjecture on polynomials,
\newblock arXiv:1912.12545 (2019).
\newblock \url{https://arxiv.org/abs/1912.12545}.

\bibitem{Dwork1960}
B.~Dwork,
\newblock On the rationality of the zeta function of an algebraic variety,
\newblock \emph{Amer. J. Math.} \textbf{82} (1960), 631--648.
\newblock \url{https://doi.org/10.2307/2372974}.

\bibitem{EMHS2019}
A.-M.~Ernvall-Hyt\"onen, T.~Matala-aho and L.~Sepp\"al\"a,
\newblock Euler's divergent series in arithmetic progressions,
\newblock \emph{J. Integer Seq.} \textbf{22} (2019), Article 19.2.8.
\newblock \url{https://arxiv.org/abs/1809.03859}.

\bibitem{EMHS2023}
A.-M.~Ernvall-Hyt\"onen, T.~Matala-aho and L.~Sepp\"al\"a,
\newblock Euler's factorial series, Hardy integral, and continued fractions,
\newblock \emph{J. Number Theory} \textbf{244} (2023), 224--250.
\newblock \url{https://arxiv.org/abs/2111.13649}.

\bibitem{FR2024ED}
S.~Fischler and T.~Rivoal,
\newblock Relations between values of arithmetic Gevrey series, and
\newblock applications to values of the Gamma function,
\newblock \emph{J. Number Theory} \textbf{261} (2024), 36--54.
\newblock \url{https://arxiv.org/abs/2301.13518}.

\bibitem{FRlog}
S.~Fischler and T.~Rivoal,
\newblock Transcendence of values of logarithms of \(E\)-functions,
\newblock \emph{Math. Ann.} (2026), Article 12.
\newblock \url{https://doi.org/10.1007/s00208-026-03374-z}.

\bibitem{FresanJossen}
J.~Fres\'an and P.~Jossen,
\newblock \emph{Exponential Motives},
\newblock research monograph manuscript, Sec.~12.8.
\newblock \url{http://javier.fresan.perso.math.cnrs.fr/expmot.pdf}.

\bibitem{HeilmannLieb1972}
O.~J.~Heilmann and E.~H.~Lieb,
\newblock Theory of monomer--dimer systems,
\newblock \emph{Comm. Math. Phys.} \textbf{25} (1972), 190--232.
\newblock \url{https://doi.org/10.1007/BF01877590}.

\bibitem{Hien2007}
M.~Hien,
\newblock Periods for irregular singular connections on surfaces,
\newblock \emph{Math. Ann.} \textbf{337} (2007), 631--669.
\newblock \url{https://doi.org/10.1007/s00208-006-0050-6}.

\bibitem{Hien2009}
M.~Hien,
\newblock Periods for flat algebraic connections,
\newblock \emph{Invent. Math.} \textbf{178} (2009), 1--22.
\newblock \url{https://doi.org/10.1007/s00222-009-0185-7}.

\bibitem{JimboMiwa1981}
M.~Jimbo and T.~Miwa,
\newblock Monodromy preserving deformation of linear ordinary differential
equations with rational coefficients, II,
\newblock \emph{Physica D} \textbf{2} (1981), 407--448.
\newblock \url{https://doi.org/10.1016/0167-2789(81)90021-X}.

\bibitem{HuberWustholz}
A.~Huber and G.~W\"ustholz,
\newblock \emph{Transcendence and Linear Relations of 1-Periods},
\newblock Cambridge Tracts in Mathematics 227, Cambridge University Press,
2022.
\newblock \url{https://arxiv.org/abs/1805.10104}.

\bibitem{Katz1970}
N.~M.~Katz,
\newblock Nilpotent connections and the monodromy theorem,
\newblock \emph{Publ. Math. IH\'ES} \textbf{39} (1970), 175--232.
\newblock \url{https://doi.org/10.1007/BF02684688}.

\bibitem{Kaad2011}
J.~Kaad,
\newblock Comparison of secondary invariants of algebraic (K)-theory,
\newblock \emph{J. K-Theory} \textbf{8} (2011), 217--269.

\bibitem{KedlayaIsocrystals}
K.~S.~Kedlaya,
\newblock Notes on isocrystals,
\newblock \emph{J. Number Theory} \textbf{237} (2022), 353--394.
\newblock \url{https://arxiv.org/abs/1606.01321}.


\bibitem{Kolchin}
E.~R.~Kolchin,
\newblock \emph{Differential Algebra and Algebraic Groups},
\newblock Academic Press, New York, 1973.

\bibitem{Lagarias2013}
J.~C.~Lagarias,
\newblock Euler's constant: Euler's work and modern developments,
\newblock \emph{Bull. Amer. Math. Soc.} \textbf{50} (2013), 527--628.

\bibitem{Mahler1968}
K.~Mahler,
\newblock Applications of a theorem by A.~B.~Shidlovski,
\newblock \emph{Proc. Roy. Soc. London Ser. A} \textbf{305} (1968), 149--173.
\newblock \url{https://carmamaths.org/resources/mahler/docs/169.pdf}.

\bibitem{MZ2018}
T.~Matala-aho and W.~Zudilin,
\newblock Euler's factorial series and global relations,
\newblock \emph{J. Number Theory} \textbf{186} (2018), 202--210.
\newblock \url{https://arxiv.org/abs/1703.02633}.

\bibitem{MohajerSebbar2025}
M.~Mohajer and A.~Sebbar,
\newblock \(p\)-adic period conjectures for \(1\)-motives: integration and
linear relations,
\newblock arXiv:2507.15020 (2025).
\newblock \url{https://arxiv.org/abs/2507.15020}.

\bibitem{Migler2017}
J.~Migler,
\newblock Functional calculus and joint torsion of pairs of almost commuting
operators,
\newblock \emph{J. Noncommut. Geom.} \textbf{11} (2017), 1377--1425.
\newblock \url{https://arxiv.org/abs/1409.6289}.

\bibitem{Moree2014}
P.~Moree,
\newblock Nicolaas Govert de Bruijn, the enchanter of friable integers,
\newblock arXiv:1212.1579.
\newblock \url{https://arxiv.org/abs/1212.1579}.

\bibitem{MuLyu2025}
X.~Mu and S.~Lyu,
\newblock Hankel determinants for a deformed Laguerre weight with multiple
variables and generalized Painlev\'e V equation,
\newblock \emph{J. Math. Phys.} \textbf{66} (2025), 073506.
\newblock \url{https://arxiv.org/abs/2410.22780}.

\bibitem{Powers2025}
M.~R.~Powers,
\newblock Transcendence results for \(\Gamma^{(n)}(1)\) and related sequences
of generalized constants,
\newblock preprint, arXiv:2511.01849v7 (2026).
\newblock \url{https://arxiv.org/abs/2511.01849}.

\bibitem{AminovArnaudo2026}
G.~Aminov and P.~Arnaudo,
\newblock Basics of multiple polyexponential integrals,
\newblock \emph{Lett. Math. Phys.} \textbf{116} (2026), Article 68.
\newblock \url{https://doi.org/10.1007/s11005-026-02090-8}.

\bibitem{Boyadzhiev2007}
K.~N.~Boyadzhiev,
\newblock Polyexponentials,
\newblock arXiv:0710.1332 (2007).
\newblock \url{https://arxiv.org/abs/0710.1332}.

\bibitem{Denisov2004}
S.~A.~Denisov,
\newblock On Rakhmanov's theorem for Jacobi matrices,
\newblock \emph{Proc. Amer. Math. Soc.} \textbf{132} (2004), 847--852.
\newblock \url{https://doi.org/10.1090/S0002-9939-03-07157-0}.

\bibitem{GouldQuaintance2007}
H.~W.~Gould and J.~Quaintance,
\newblock A linear binomial recurrence and the Bell numbers and polynomials,
\newblock \emph{Appl. Anal. Discrete Math.} \textbf{1} (2007), 371--385.
\newblock \url{https://doi.org/10.2298/AADM0702371G}.

\bibitem{Hardy1905PartII}
G.~H.~Hardy,
\newblock On the zeroes of certain classes of integral Taylor series,
Part II---on the integral function
\(\sum x^n/((n+a)^s n!)\) and other similar functions,
\newblock \emph{Proc. London Math. Soc.} (2) \textbf{2} (1905), 401--431.
\newblock \url{https://doi.org/10.1112/plms/s2-2.1.401}.

\bibitem{PrasadRajkumarReddy2019}
D.~Prasad, K.~Rajkumar and A.~S.~Reddy,
\newblock A survey on fixed divisors,
\newblock \emph{Confluentes Math.} \textbf{11} (2019), no.~1, 29--52.
\newblock \url{https://doi.org/10.5802/cml.54}.

\bibitem{Roman1984}
S.~Roman,
\newblock \emph{The Umbral Calculus},
\newblock Pure and Applied Mathematics 111, Academic Press, New York, 1984.

\bibitem{Shidlovskii1989}
A.~B.~Shidlovskii,
\newblock \emph{Transcendental Numbers},
\newblock De Gruyter Studies in Mathematics 12, Walter de Gruyter,
Berlin--New York, 1989.
\newblock \url{https://doi.org/10.1515/9783110889055}.

\bibitem{Sondow2003Criteria}
J.~Sondow,
\newblock Criteria for irrationality of Euler's constant,
\newblock \emph{Proc. Amer. Math. Soc.} \textbf{131} (2003), 3335--3344.
\newblock \url{https://doi.org/10.1090/S0002-9939-03-07081-3}.

\bibitem{Sondow2009Hypergeometric}
J.~Sondow and S.~Zlobin,
\newblock A hypergeometric approach, via linear forms involving logarithms,
to criteria for irrationality of Euler's constant,
\newblock \emph{Math. Slovaca} \textbf{59} (2009), 307--314.
\newblock \url{https://doi.org/10.2478/s12175-009-0127-2}.

\bibitem{SunZagier2011}
Z.-W.~Sun and D.~Zagier,
\newblock On a curious property of Bell numbers,
\newblock \emph{Bull. Aust. Math. Soc.} \textbf{84} (2011), 153--158.
\newblock \url{https://doi.org/10.1017/S0004972711002218}.

\bibitem{Touchard1933}
J.~Touchard,
\newblock Propri\'et\'es arithm\'etiques de certains nombres r\'ecurrents,
\newblock \emph{Ann. Soc. Sci. Bruxelles Ser. A} \textbf{53} (1933), 21--31.

\bibitem{Snodgrass2026}
B.~Snodgrass,
\newblock Periods of \(E\)-operators,
\newblock arXiv:2608.06005v2 (2026).
\newblock \url{https://arxiv.org/abs/2608.06005}.

\end{thebibliography}

\end{document}
